% Leçon 37 — Éléments socioculturels : les outils TICs de marques chinoises et camerounaises — Chinois Tle A — Cours signé OnBuch+
% Programme officiel MINESEC. Compiler avec : tectonic ZH37-socio-les-outils-tics-de-marques-chinoises-et-came.tex
\newcommand{\DOCMATIERE}{Chinois}
\newcommand{\DOCNIVEAU}{Tle A}
\newcommand{\DOCMODULE}{Module 5 : Mass médias et télécommunication}
\newcommand{\DOCLECON}{ZH37}
\newcommand{\DOCTITRE}{Leçon 37 — Éléments socioculturels : les outils TICs de marques chinoises et camerounaises}
% OnBuch+ — préambule « cours signés » ENRICHI (compilé avec tectonic / XeLaTeX)
% Système visuel « Pop » d'OnBuch+ : fond crème (jamais blanc), bordures encre
% épaisses, ombres « dures » décalées, titres Archivo Black en capitales.
% Version MATHÉMATIQUES : graphes (pgfplots/tikz), boîtes pédagogiques
% (définition, propriété, méthode, attention, le savais-tu, exercice résolu,
% exercice, TP, bilan), page de garde et signature OnBuch+ ancrée en bas.
%
% Macros attendues dans main.tex AVANT \input{preamble} :
%   \DOCMATIERE  \DOCNIVEAU  \DOCMODULE  \DOCLECON (numéro)  \DOCTITRE
\documentclass[11pt]{article}

\usepackage{fontspec}
\usepackage[french]{babel} % français = langue principale ; le chinois via \zh{..} / textebox (xeCJK)
\usepackage{xeCJK}
\usepackage{pifont} % \ding{51} \ding{55} utilisés par les modèles
\usepackage[a4paper,margin=2cm,top=3.4cm,bottom=2.8cm,headheight=1.4cm,footskip=1.3cm]{geometry}
\usepackage{amsmath,amssymb,amsfonts}
\usepackage{mathtools} % \xmapsto, \coloneqq... souvent utilisés par les modèles
\usepackage{cancel} % simplifications barrées (bilans de forces)
% \abs{z} (module, valeur absolue) et \norm{u} (norme), utilisés par les modèles
\providecommand{\abs}{}\let\abs\relax\DeclarePairedDelimiter{\abs}{\lvert}{\rvert}
\providecommand{\norm}{}\let\norm\relax\DeclarePairedDelimiter{\norm}{\lVert}{\rVert}
\usepackage[table]{xcolor} % [table] : fournit aussi \rowcolors (alternance de lignes)
\usepackage{pagecolor}
\usepackage{tikz}
\usetikzlibrary{arrows.meta,positioning,calc,shapes.geometric,shapes.misc,shapes.arrows,decorations.pathmorphing,decorations.pathreplacing,decorations.markings,patterns,shadows,mindmap,trees,fit,backgrounds,matrix,intersections,angles,quotes}
\usepackage{pgfplots}
\usepackage[version=4]{mhchem} % équations nucléaires \ce{^{235}_{92}U}
\usepackage[european,siunitx]{circuitikz} % schémas électriques
\pgfplotsset{compat=1.18}
\usepgfplotslibrary{fillbetween,groupplots} % aires entre courbes (calcul intégral)
\usepackage{enumitem}
\usepackage{fancyhdr}
% [most] charge toutes les bibliothèques tcolorbox (listings, minted, raster,
% documentation...) et gonfle inutilement la mémoire TeX sur un document long
% (~300 boîtes) : on ne charge que ce qui est réellement utilisé.
\usepackage[skins,breakable]{tcolorbox}
\usepackage{graphicx}
\usepackage{colortbl}
\usepackage{booktabs}
\usepackage{tabularx}
\usepackage{diagbox} % case d'en-tête diagonale des tableaux à double entrée
% Colonnes extensibles centrée / gauche / droite (C, L, R), souvent utilisées par les modèles
\newcolumntype{C}{>{\centering\arraybackslash}X}
\newcolumntype{L}{>{\raggedright\arraybackslash}X}
\newcolumntype{R}{>{\raggedleft\arraybackslash}X}
\usepackage{array}
\usepackage{multicol}
\usepackage{multirow}
\usepackage{siunitx}
% parse-numbers=false : accepte aussi \SI{2,5 \times 10^{-3}}{...}
\sisetup{locale=FR,output-decimal-marker={,},group-minimum-digits=4,parse-numbers=false}
\DeclareSIUnit\ohm{\ensuremath{\symup{\Omega}}} % U+2126 absent de la police
\DeclareSIUnit\division{div} % réglages d'oscilloscope (ms/div, V/div)
\DeclareSIUnit\tr{tr} % tours (vitesse de rotation, ex. \SI{1500}{\tr\per\minute})
\DeclareSIUnit\molar{M} % concentration molaire (ex. \SI{3}{\molar}, \SI{1}{\micro\molar})
\DeclareSIUnit\an{an} % année (le modèle confond parfois avec \year, un compteur TeX) — ex. \SI{5}{\kilo\meter\per\an}
\DeclareSIUnit\FCFA{FCFA} % franc CFA (ex. \SI{500}{\FCFA\per\kilo\gram})
\DeclareSIUnit\degreeBrix{\textdegree Brix} % teneur en sucre d'un sirop/jus (réfractométrie)
\DeclareSIUnit\month{mois} % le modèle utilise parfois \month, un compteur TeX, comme unité de durée
\usepackage{qrcode}
% \degree : commande fréquemment utilisée par les modèles pour un angle ou une
% température en dehors de \SI{}{} (non fournie par les paquets chargés).
\providecommand{\degree}{^{\circ}}
% \qed (fin de démonstration), utilisé par les modèles sans amsthm
\providecommand{\qed}{\hfill$\square$}
% \barre : double barre des valeurs interdites dans un tableau de variations (tkz-tab)
\providecommand{\barre}{\ensuremath{\|}}
% environnement proof (amsthm non chargé) utilisé par les modèles
\ifdefined\proof\else\newenvironment{proof}[1][Démonstration]{\par\noindent\textit{#1.}\ }{\qed\par}\fi
\usepackage{lastpage}
\usepackage{hyperref}
\hypersetup{hidelinks}

% ============================================================
% Palette « Pop » OnBuch+ — mêmes hex que la classe P (plus_ui.dart)
% ============================================================
\definecolor{poporange}{HTML}{F59321}
\definecolor{popdark}{HTML}{E07A0C}
\definecolor{popcream}{HTML}{FFF6E8}
\definecolor{popink}{HTML}{1C1714}
\definecolor{poppurple}{HTML}{7A5AE0}
\definecolor{popgreen}{HTML}{2E9E6A}
\definecolor{poppink}{HTML}{E0457B}
\definecolor{popblue}{HTML}{2F7DE1}
\definecolor{popgold}{HTML}{C98A12}
\definecolor{popmuted}{HTML}{8A7F76}
\definecolor{popline}{HTML}{EADFD2}
\definecolor{popgray}{HTML}{8A7F76}
\colorlet{popgrey}{popgray}
% Alias tolérants (le modèle nomme parfois les couleurs différemment)
\colorlet{popwhite}{white}
\colorlet{popblack}{popink}
\colorlet{popred}{poppink}
\colorlet{popyellow}{popgold}
% Teintes claires (fonds de tableaux/encadrés) — le modèle nomme parfois
% les couleurs avec un suffixe L (ex. popblueL) sans les avoir définies.
\colorlet{poporangeL}{poporange!20}
\colorlet{popdarkL}{popdark!20}
\colorlet{poppurpleL}{poppurple!20}
\colorlet{popgreenL}{popgreen!20}
\colorlet{poppinkL}{poppink!20}
\colorlet{popblueL}{popblue!20}
\colorlet{popgoldL}{popgold!20}
\colorlet{popredL}{popred!20}
\colorlet{popgrayL}{popgray!20}
% Teintes claires pour fonds de tableaux / zones
\colorlet{popblueL}{popblue!10!white}
\colorlet{poporangeL}{poporange!14!white}
\colorlet{popgreenL}{popgreen!12!white}
\colorlet{poppurpleL}{poppurple!12!white}
\colorlet{poppinkL}{poppink!12!white}
\colorlet{popgoldL}{popgold!14!white}

\pagecolor{popcream}

% ============================================================
% Typographie
% ============================================================
\defaultfontfeatures{Ligatures=TeX}
\setmainfont{PlusJakartaSans-Medium.ttf}[Path=../../../fonts/,BoldFont=PlusJakartaSans-Bold.ttf]
% Symboles, API et emojis absents de la police principale : repli sur DejaVu Sans (caractères actifs)
\newfontface\symfont{DejaVuSans.ttf}[Path=../../../fonts/]
\makeatletter
\newcommand\symactive[1]{\catcode#1=13 \begingroup\lccode`\~=#1 \lowercase{\endgroup\def~}{{\symfont\char#1}}}
\newcommand\symblank[1]{\catcode#1=13 \begingroup\lccode`\~=#1 \lowercase{\endgroup\def~}{}}
\newcommand\symhyph[1]{\catcode#1=13 \begingroup\lccode`\~=#1 \lowercase{\endgroup\def~}{\mbox{-}}}
\newcount\symc
\newcommand\symrange[2]{\symc=#1 \@whilenum\symc<#2 \do{\expandafter\symactive\expandafter{\the\symc}\advance\symc by 1 }}
\newcommand\symblankrange[2]{\symc=#1 \@whilenum\symc<#2 \do{\expandafter\symblank\expandafter{\the\symc}\advance\symc by 1 }}
\makeatother
\symrange{"0250}{"02FF}\symactive{"2713}\symactive{"2714}\symactive{"2717}\symactive{"2718}\symactive{"2610}\symactive{"2611}\symactive{"2612}
\symrange{"2600}{"27BF}
\catcode"2705=13 \begingroup\lccode`\~="2705 \lowercase{\endgroup\def~}{{\symfont\char"2713}}\symblank{"26BD}\symblank{"26A1}
\symrange{"2460}{"257F}\symactive{"223C}\symactive{"2248}\symactive{"2260}
\symactive{"01F8}\symactive{"01F9}\symactive{"1E3E}\symactive{"1E3F}
\symhyph{"2011}\symhyph{"2E1A}\symblankrange{"1F300}{"1FAFF}\symblank{"FE0F}
% Caractères chinois : WenQuanYi Zen Hei (sous-ensemble GB2312 + pinyin), gras/italique simulés
\setCJKmainfont{WenQuanYiZenHei-subset.ttf}[Path=../../../fonts/,AutoFakeBold=3,AutoFakeSlant=0.2]
\xeCJKsetup{CJKspace=false}
% Pinyin : voyelles à caron (ǎ ǐ ǒ ǔ ǖ ǘ ǚ ǜ) absentes de la police principale -> caractères actifs
\newfontface\pyfont{WenQuanYiZenHei-subset.ttf}[Path=../../../fonts/]
\newcommand\pyactive[1]{\catcode#1=13 \begingroup\lccode`\~=#1 \lowercase{\endgroup\def~}{{\pyfont\char#1}}}
\pyactive{"01CD}\pyactive{"01CE}\pyactive{"01CF}\pyactive{"01D0}\pyactive{"01D1}\pyactive{"01D2}\pyactive{"01D3}\pyactive{"01D4}\pyactive{"01D5}\pyactive{"01D6}\pyactive{"01D7}\pyactive{"01D8}\pyactive{"01D9}\pyactive{"01DA}\pyactive{"01DB}\pyactive{"01DC}
% Maths en sans-serif (Fira Math) avec chiffres et lettres droites de Plus
% Jakarta, pour que les formules se fondent dans le texte.
\usepackage[math-style=ISO,bold-style=ISO]{unicode-math}
\setmathfont{FiraMath-Regular.otf}[Path=../../../fonts/]
\setmathfont{PlusJakartaSans-Medium.ttf}[Path=../../../fonts/,range={up/num,up/{Latin,latin},\mathup}]
\usepackage{icomma} % virgule décimale sans espace en mode math
% Caractères mathématiques Unicode absents des polices de texte (Plus Jakarta,
% Archivo Black) : rendus par leur équivalent mathématique, en texte comme en
% formule — sinon ils s'affichent en carré vide (ex. « Arithmétique dans ℤ »).
\usepackage{newunicodechar}
\newunicodechar{ℝ}{\ensuremath{\mathbb{R}}}
\newunicodechar{ℤ}{\ensuremath{\mathbb{Z}}}
\newunicodechar{ℂ}{\ensuremath{\mathbb{C}}}
\newunicodechar{ℕ}{\ensuremath{\mathbb{N}}}
\newunicodechar{ℚ}{\ensuremath{\mathbb{Q}}}
\newunicodechar{𝔻}{\ensuremath{\mathbb{D}}}
\newunicodechar{α}{\ensuremath{\alpha}}
\newunicodechar{β}{\ensuremath{\beta}}
\newunicodechar{γ}{\ensuremath{\gamma}}
\newunicodechar{δ}{\ensuremath{\delta}}
\newunicodechar{ε}{\ensuremath{\varepsilon}}
\newunicodechar{θ}{\ensuremath{\theta}}
\newunicodechar{λ}{\ensuremath{\lambda}}
\newunicodechar{μ}{\ensuremath{\mu}}
\newunicodechar{σ}{\ensuremath{\sigma}}
\newunicodechar{τ}{\ensuremath{\tau}}
\newunicodechar{φ}{\ensuremath{\varphi}}
\newunicodechar{ψ}{\ensuremath{\psi}}
\newunicodechar{ω}{\ensuremath{\omega}}
\newunicodechar{Σ}{\ensuremath{\Sigma}}
% majuscules produites par \MakeUppercase dans les titres (α-aminés -> Α)
\newunicodechar{Α}{\ensuremath{\alpha}}
\newunicodechar{Β}{\ensuremath{\beta}}
\newunicodechar{Γ}{\ensuremath{\Gamma}}
\newunicodechar{Δ}{\ensuremath{\Delta}}
\newunicodechar{Ε}{\ensuremath{\varepsilon}}
\newunicodechar{Θ}{\ensuremath{\Theta}}
\newunicodechar{Λ}{\ensuremath{\Lambda}}
\newunicodechar{Μ}{\ensuremath{\mu}}
\newunicodechar{Τ}{\ensuremath{\tau}}
\newunicodechar{Φ}{\ensuremath{\Phi}}
\newunicodechar{Ψ}{\ensuremath{\Psi}}
\newunicodechar{Ω}{\ensuremath{\Omega}}
\newunicodechar{↦}{\ensuremath{\mapsto}}
\newunicodechar{∈}{\ensuremath{\in}}
\newunicodechar{∉}{\ensuremath{\notin}}
\newunicodechar{∧}{\ensuremath{\wedge}}
\newunicodechar{⇔}{\ensuremath{\Leftrightarrow}}
\newunicodechar{⟺}{\ensuremath{\Longleftrightarrow}}
\newunicodechar{⇒}{\ensuremath{\Rightarrow}}
\newunicodechar{⟹}{\ensuremath{\Longrightarrow}}
\newunicodechar{∘}{\ensuremath{\circ}}
\newunicodechar{∪}{\ensuremath{\cup}}
\newunicodechar{∩}{\ensuremath{\cap}}
\newunicodechar{≡}{\ensuremath{\equiv}}
\newunicodechar{⊂}{\ensuremath{\subset}}
\newunicodechar{⊥}{\ensuremath{\perp}}
\newunicodechar{∃}{\ensuremath{\exists}}
\newunicodechar{∀}{\ensuremath{\forall}}
\newunicodechar{∛}{\ensuremath{\sqrt[3]{\,}}}
\newunicodechar{∜}{\ensuremath{\sqrt[4]{\,}}}
\newunicodechar{⃗}{\ensuremath{{}^{\to}}}
\newunicodechar{ⁿ}{\ifmmode^{n}\else\textsuperscript{\ensuremath{n}}\fi}
\newunicodechar{ˣ}{\ifmmode^{x}\else\textsuperscript{\ensuremath{x}}\fi}
\newunicodechar{ᵇ}{\ifmmode^{b}\else\textsuperscript{\ensuremath{b}}\fi}
\newunicodechar{ʳ}{\ifmmode^{r}\else\textsuperscript{\ensuremath{r}}\fi}
\newunicodechar{ᵘ}{\ifmmode^{u}\else\textsuperscript{\ensuremath{u}}\fi}
\newunicodechar{ʸ}{\ifmmode^{y}\else\textsuperscript{\ensuremath{y}}\fi}
\newunicodechar{ᵃ}{\ifmmode^{a}\else\textsuperscript{\ensuremath{a}}\fi}
\newunicodechar{ᵉ}{\ifmmode^{e}\else\textsuperscript{\ensuremath{e}}\fi}
\newunicodechar{ᵏ}{\ifmmode^{k}\else\textsuperscript{\ensuremath{k}}\fi}
\newunicodechar{ᵖ}{\ifmmode^{p}\else\textsuperscript{\ensuremath{p}}\fi}
\newunicodechar{ᴺ}{\ifmmode^{N}\else\textsuperscript{\ensuremath{N}}\fi}
\newunicodechar{ᶜ}{\ifmmode^{c}\else\textsuperscript{\ensuremath{c}}\fi}
\newunicodechar{ʰ}{\ifmmode^{h}\else\textsuperscript{\ensuremath{h}}\fi}
\newunicodechar{ᵠ}{\ifmmode^{q}\else\textsuperscript{\ensuremath{q}}\fi}
\newunicodechar{ₙ}{\ifmmode_{n}\else\textsubscript{\ensuremath{n}}\fi}
\newunicodechar{ₐ}{\ifmmode_{a}\else\textsubscript{\ensuremath{a}}\fi}
\newunicodechar{ᵢ}{\ifmmode_{i}\else\textsubscript{\ensuremath{i}}\fi}
\newunicodechar{ᵣ}{\ifmmode_{r}\else\textsubscript{\ensuremath{r}}\fi}
\newunicodechar{ₖ}{\ifmmode_{k}\else\textsubscript{\ensuremath{k}}\fi}
\newunicodechar{ᵦ}{\ifmmode_{\beta}\else\textsubscript{\ensuremath{\beta}}\fi}
\newunicodechar{ₓ}{\ifmmode_{x}\else\textsubscript{\ensuremath{x}}\fi}
\newfontfamily\popdisplay{ArchivoBlack-Regular.ttf}[Path=../../../fonts/]
\newfontfamily\poplogo{SpaceGrotesk-Bold.ttf}[Path=../../../fonts/]
\newfontfamily\popsemi{PlusJakartaSans-SemiBold.ttf}[Path=../../../fonts/]
\newfontfamily\popxbold{PlusJakartaSans-ExtraBold.ttf}[Path=../../../fonts/]
\newfontfamily\popxboldls{PlusJakartaSans-ExtraBold.ttf}[Path=../../../fonts/,LetterSpace=3.0]

\setlist[enumerate]{leftmargin=1.7em,itemsep=5pt,topsep=4pt}
\setlist[itemize]{leftmargin=1.7em,itemsep=4pt,topsep=4pt,label={\color{poporange}\textbullet}}
\setlength{\parindent}{0pt}
\setlength{\parskip}{5pt}
\renewcommand{\arraystretch}{1.25}

% pgfplots : style Pop par défaut
\pgfplotsset{
  every axis/.append style={axis line style={popink,line width=1pt},tick style={popink},
    grid style={popline,line width=0.5pt},label style={font=\small},tick label style={font=\footnotesize}},
  popcurve/.style={poporange,line width=1.8pt,no markers,smooth},
}
% Même style disponible dans un \draw TikZ simple (hors pgfplots)
\tikzset{popcurve/.style={poporange,line width=1.8pt}}
\tikzset{popgrid/.style={popline,very thin}}
\colorlet{popcurve}{poporange} % le nom du style est parfois utilisé comme couleur

% ============================================================
% Pill / squiggle
% ============================================================
\newcommand{\poppill}[3]{%
  \tikz[baseline=-0.55ex]\node[draw=popink,line width=1.2pt,rounded corners=2.6pt,%
    fill=#2,inner xsep=6.5pt,inner ysep=3pt]{%
    \popxboldls\fontsize{7}{7}\selectfont\color{#3}\MakeUppercase{#1}};%
}
\newcommand{\popsquiggle}[1]{%
  \tikz[baseline=-0.4ex]{
    \draw[#1,line width=2.6pt,line cap=round]
      plot [smooth] coordinates {(0,0.09) (0.34,0.01) (0.68,0.21) (1.02,0.01) (1.36,0.10)};
  }%
}
% Mot-clé surligné (marqueur orange sous le texte)
\newcommand{\cle}[1]{{\popxbold\color{popdark}#1}}

% ============================================================
% En-tête / pied
% ============================================================
\pagestyle{fancy}
\fancyhf{}
\renewcommand{\headrulewidth}{0pt}
\renewcommand{\footrulewidth}{0pt}
\lhead{\raisebox{-0.15ex}{\poplogo\fontsize{13}{13}\selectfont\color{popink}OnBuch\color{poporange}+}}
\rhead{\poppill{\DOCMATIERE\ \textbullet\ \DOCNIVEAU\ \textbullet\ Leçon \DOCLECON}{white}{popink}}
\lfoot{\popxbold\fontsize{7.6}{7.6}\selectfont\color{popmuted}\MakeUppercase{Cours signé OnBuch+}\ \textbullet\ {\popsemi\DOCTITRE}}
\rfoot{\tikz[baseline=-0.6ex]\node[rounded corners=8.5pt,fill=popink,minimum height=17pt,minimum width=17pt,inner xsep=5pt,inner sep=0]{\popxbold\fontsize{8}{8}\selectfont\color{popcream}\thepage\,/\,\pageref*{LastPage}};}

% ============================================================
% Titres
% ============================================================
\newcommand{\chaptitre}[2]{%
  \begin{tcolorbox}[enhanced,colback=poporange,colframe=popink,boxrule=2.6pt,arc=11pt,
      drop shadow=popink,left=15pt,right=15pt,top=12pt,bottom=11pt,before skip=3pt,after skip=18pt]
    \poppill{#2}{popink}{popcream}\par\vspace{9pt}
    {\raggedright\hyphenpenalty=10000\exhyphenpenalty=10000\popdisplay\fontsize{22}{24}\selectfont\color{popcream}\MakeUppercase{#1}\par}
  \end{tcolorbox}
}
% Section numérotée : pastille orange avec le numéro + titre Archivo
\newcounter{popsec}
\newcommand{\coursec}[1]{%
  \stepcounter{popsec}\setcounter{popsub}{0}%
  \par\vspace{16pt}\noindent
  \begin{minipage}{\linewidth}
  \tikz[baseline=-0.7ex]\node[draw=popink,line width=1.6pt,fill=poporange,rounded corners=4pt,minimum size=22pt,inner sep=0]{\popdisplay\fontsize{12}{12}\selectfont\color{popcream}\thepopsec};%
  \hspace{8pt}{\popdisplay\fontsize{15}{17}\selectfont\color{popink}\MakeUppercase{#1}}\par
  \vspace{4pt}\noindent{\color{poporange}\rule{36pt}{3.6pt}}
  \end{minipage}\par\nopagebreak\vspace{8pt}
}
\newcounter{popsub}
\newcommand{\courssub}[1]{%
  \stepcounter{popsub}%
  \par\vspace{9pt}\noindent{\popxbold\fontsize{11.5}{13}\selectfont\color{popink}{\color{poporange}\thepopsec.\thepopsub}\hspace{6pt}#1}\par\nopagebreak\vspace{4pt}
}

% ============================================================
% Boîtes pédagogiques — même recette : fond blanc, bordure épaisse colorée,
% ombre dure encre, badge plein. Titre optionnel [#1].
% ============================================================
\tcbset{popbase/.style={breakable,enhanced,colback=white,boxrule=2.3pt,arc=9pt,
  shadow={3pt}{-3pt}{0pt}{popink},left=14pt,right=14pt,top=11pt,bottom=11pt,before skip=11pt,after skip=15pt,
  fontupper=\popsemi\fontsize{10.6}{14.5}\selectfont\color{popink},
  fontlower=\popsemi\fontsize{10.6}{14.5}\selectfont\color{popink}}}
% Coche dessinée (le glyphe ✓ n'existe pas dans les polices du document)
\newcommand{\popcheck}[1]{\tikz[baseline=-0.5ex]\draw[#1,line width=1.6pt,line cap=round,line join=round](-0.13,0.0)--(-0.04,-0.09)--(0.13,0.1);}
\providecommand{\coche}{\popcheck{popgreen}}
\providecommand{\resultat}[1]{\par\smallskip\noindent\fcolorbox{popgreen}{popgreenL}{\parbox{\dimexpr\linewidth-2\fboxsep-2\fboxrule\relax}{\textbf{Résultat.} #1}}\par\smallskip}
\newcommand{\popboxhead}[3]{\poppill{#1}{#2}{white}\ifx\relax#3\relax\else\hspace{6pt}{\popxbold\fontsize{10.6}{12}\selectfont\color{popink}#3}\fi\par\vspace{7pt}}

\newtcolorbox{aretenir}[1][]{popbase,colframe=popgreen,before upper={\popboxhead{À retenir}{popgreen}{#1}}}
% Texte en chinois (document de lecture, dialogue, extrait) : cadre
% d'étude. Un titre optionnel (source, auteur) s'affiche dans l'en-tête.
\newtcolorbox{textebox}[1][]{popbase,colframe=popblue,before upper={\popboxhead{文本 · Texte}{popblue}{#1}},after upper={}}
% Marques ✓ / ✗ (forme correcte / fautive) souvent utilisées par les modèles
\providecommand{\cross}{\textcolor{poppink}{\ensuremath{\times}}}
\providecommand{\xmark}{\cross}\providecommand{\crossmark}{\cross}\providecommand{\cmark}{\ensuremath{\checkmark}}
% Ligne de réponse / justification à compléter (inventée par les modèles)
\providecommand{\justification}[1]{\par\noindent\textit{Justification~:} \dotfill\par}
% \textipa (package tipa non chargé) : la police ne couvre pas l'API -> simple passage
\providecommand{\textipa}[1]{#1}
% Soulignement (ulem non chargé) : \uline, \ul
\providecommand{\uline}[1]{\underline{#1}}\providecommand{\ul}[1]{\underline{#1}}
% \glossaire{\item ...} inventée par les modèles : liste de glossaire
\providecommand{\glossaire}[1]{\begin{itemize}#1\end{itemize}}
% \alert (beamer) inventée par les modèles : texte mis en évidence
\providecommand{\alert}[1]{\textbf{\textcolor{poppink}{#1}}}
% variantes ASCII d'ordinaux écrites en macro
\providecommand{\iere}{\textsuperscript{re}}\providecommand{\ieme}{\textsuperscript{e}}\providecommand{\ere}{\textsuperscript{re}}\providecommand{\eme}{\textsuperscript{e}}\providecommand{\er}{\textsuperscript{er}}\providecommand{\ie}{\textsuperscript{e}}
% Ordinaux français écrits en macro par les modèles : 1\ière, 3\ième
\newcommand{\ière}{\textsuperscript{re}}\newcommand{\ième}{\textsuperscript{e}}
% \exemple (inventée par les modèles) : étiquette « Exemple : »
\providecommand{\exemple}{\textbf{Exemple~:}\ }
% \square utilisable aussi hors mode math (cases à cocher)
\AtBeginDocument{\let\squareMath\square\renewcommand{\square}{\ensuremath{\squareMath}}}
% Mot ou phrase en chinois à l'intérieur d'un paragraphe français
\newcommand{\zh}[1]{\ifmmode\text{#1}\else{#1}\fi}
\let\eng\zh \let\esp\zh \let\all\zh % alias tolérés
% flèches d'intonation inventées par les modèles
\providecommand{\textsearrow}{}\renewcommand{\textsearrow}{\ensuremath{\searrow}}\providecommand{\textnearrow}{}\renewcommand{\textnearrow}{\ensuremath{\nearrow}}
% \fr{..} : mot français (inventé par les modèles)
\providecommand{\fr}[1]{\textit{#1}}
% zhbox : encadré de texte chinois inventé par les modèles
\newenvironment{zhbox}{\begin{quote}}{\end{quote}}
% \courssubsub : titre de niveau 3 inventé par les modèles
\providecommand{\courssubsub}[1]{\par\medskip\noindent\textbf{#1}\par\smallskip}
% \zhbf : texte chinois en gras inventé par les modèles
\providecommand{\zhbf}[1]{\textbf{#1}}
% \vs : « versus » inventé par les modèles
\providecommand{\vs}{\textit{vs}\ }
% \blank : case à compléter inventée par les modèles
\providecommand{\blank}[1][]{\underline{\hspace{1.6em}}}
\newtcolorbox{exemplebox}[1][]{popbase,colframe=poporange,before upper={\popboxhead{Exemple}{poporange}{#1}}}
\newtcolorbox{definition}[1][]{popbase,colframe=popblue,before upper={\popboxhead{Définition}{popblue}{#1}}}
\newtcolorbox{propriete}[1][]{popbase,colframe=poppurple,before upper={\popboxhead{Propriété}{poppurple}{#1}}}
\newtcolorbox{methode}[1][]{popbase,colframe=popgold,before upper={\popboxhead{Méthode}{popgold}{#1}}}
\newtcolorbox{attention}[1][]{popbase,colframe=poppink,before upper={\popboxhead{Attention}{poppink}{#1}}}
\newtcolorbox{experience}[1][]{popbase,colframe=popblue,colback=popblueL,before upper={\popboxhead{Expérience}{popblue}{#1}}}
% « Le savais-tu ? » : carte violette pleine (contexte, culture, Cameroun)
\newtcolorbox{savaistu}[1][]{popbase,colback=poppurple,colframe=popink,
  fontupper=\popsemi\fontsize{10.4}{14.2}\selectfont\color{white},
  before upper={\poppill{Le savais-tu\,?}{white}{poppurple}\ifx\relax#1\relax\else\hspace{6pt}{\popxbold\color{white}#1}\fi\par\vspace{7pt}}}
% Exercice résolu : énoncé en haut, \tcblower puis la solution sur fond crème
\newcounter{popexo}\newcounter{popexr}
\newtcolorbox{exoresolu}[1][]{popbase,colframe=poporange,colbacklower=popcream,
  segmentation style={popink,line width=1.2pt,dashed},
  before upper={\stepcounter{popexr}\popboxhead{Exercice résolu \thepopexr}{poporange}{#1}},
  before lower={\poppill{Solution}{popink}{popcream}\par\vspace{6pt}}}
% Exercice (non résolu en place) — la correction est dans la partie Corrigés
\newtcolorbox{exercice}[1][]{popbase,colframe=popink,
  before upper={\stepcounter{popexo}\popboxhead{Exercice \thepopexo}{popink}{#1}}}
% Corrigé
\newtcolorbox{corrige}[1][]{popbase,colframe=popgreen,colback=popgreenL,
  before upper={\popboxhead{Corrigé}{popgreen}{#1}}}
% Objectifs / prérequis / bilan
\newtcolorbox{objectifs}{popbase,colframe=popink,colback=poporangeL,
  before upper={\popboxhead{Objectifs de la leçon}{popink}{}}}
\newtcolorbox{prerequis}{popbase,colframe=popink,colback=popblueL,
  before upper={\popboxhead{Prérequis}{popblue}{}}}
\newtcolorbox{bilan}{popbase,colframe=popink,colback=white,boxrule=2.8pt,
  before upper={\popboxhead{Fiche bilan}{poporange}{L'essentiel à connaître pour l'examen}}}
% Figure encadrée avec légende
\newtcolorbox{popfigure}[1][]{enhanced,breakable=false,colback=white,colframe=popline,boxrule=1.4pt,arc=8pt,
  left=10pt,right=10pt,top=10pt,bottom=8pt,before skip=10pt,after skip=12pt,halign=center,
  lower separated=false,halign lower=center,
  fontlower=\popsemi\fontsize{9}{11.5}\selectfont\color{popmuted},
  #1}
\newcommand{\legende}[1]{\par\vspace{6pt}{\popsemi\fontsize{9}{11.5}\selectfont\color{popmuted}#1}}

% ============================================================
% Signature OnBuch+
% ============================================================
\newcommand{\poplogoinline}[1]{{\poplogo\fontsize{#1}{#1}\selectfont\color{popink}OnBuch\color{poporange}+}}
\newcommand{\signatureonbuch}{%
  \begin{tcolorbox}[enhanced,colback=popink,colframe=popink,boxrule=2.2pt,arc=10pt,
      drop shadow=poporange,left=14pt,right=14pt,top=10pt,bottom=10pt,before skip=0pt,after skip=0pt]
    \tikz[baseline=-0.7ex]\node[circle,fill=popgreen,draw=popcream,line width=1.3pt,minimum size=22pt,inner sep=0]{\popcheck{white}};%
    \hspace{9pt}\begin{minipage}[c]{0.62\linewidth}
      {\popxbold\fontsize{12}{13}\selectfont\color{popcream}Signé\ \textbullet\ {\poplogo OnBuch\color{poporange}+}}\par\vspace{2pt}
      {\raggedright\popsemi\fontsize{8}{10.5}\selectfont\color{popline}\MakeUppercase{Équipe pédagogique OnBuch+}\\\MakeUppercase{Conforme au programme officiel MINESEC}\par}
    \end{minipage}\hfill
    {\poplogo\fontsize{22}{22}\selectfont\color{popcream}OnBuch\color{poporange}+}
  \end{tcolorbox}%
}

% ============================================================
% Page de garde
% #1 titre · #2 matière · #3 niveau · #4 module · #5 n° leçon · #6 durée ·
% #7 description / objectifs (texte ou itemize)
% ============================================================
\newcommand{\pagegarde}[7]{%
  \thispagestyle{empty}
  \newgeometry{a4paper,margin=1.7cm,top=1.6cm,bottom=1.4cm}%
  \begin{tikzpicture}[remember picture,overlay]
    \fill[poporange!18] ([xshift=-3.2cm,yshift=-3.6cm]current page.north east) circle (4.6cm);
    \fill[poppurple!14] ([xshift=2.4cm,yshift=7.6cm]current page.south west) circle (3.8cm);
  \end{tikzpicture}%
  {\poplogo\fontsize{30}{30}\selectfont\color{popink}OnBuch\color{poporange}+}\hfill
  \raisebox{0.6ex}{\poppill{Cours signé \textbullet\ Premium}{popink}{popcream}}\par
  \vspace{1pt}\popsquiggle{poppurple}\par
  \vspace{8pt}
  \begin{tcolorbox}[enhanced,colback=poporange,colframe=popink,boxrule=2.8pt,arc=13pt,
      drop shadow=popink,left=18pt,right=18pt,top=12pt,bottom=13pt,after skip=10pt]
    \poppill{#2 \textbullet\ #3}{popink}{popcream}\hspace{5pt}\poppill{#4}{white}{popink}\par\vspace{8pt}
    {\popdisplay\fontsize{14}{14}\selectfont\color{popink}LEÇON #5}\par\vspace{4pt}
    {\raggedright\hyphenpenalty=10000\exhyphenpenalty=10000\popdisplay\fontsize{25}{27}\selectfont\color{popcream}\MakeUppercase{#1}\par}
    \vspace{8pt}
    {\popxbold\fontsize{9.5}{9.5}\selectfont\color{popink}\MakeUppercase{Durée conseillée : #6}}
  \end{tcolorbox}
  \begin{tcolorbox}[enhanced,colback=white,colframe=popink,boxrule=2.2pt,arc=10pt,
      drop shadow=popink,left=16pt,right=16pt,top=10pt,bottom=10pt,after skip=8pt]
    \poppill{Dans cette leçon}{poppurple}{white}\par\vspace{5pt}
    \popsemi\fontsize{9.3}{12.6}\selectfont\color{popink}#7
  \end{tcolorbox}
  \noindent\begin{minipage}[t]{0.487\linewidth}
    \begin{tcolorbox}[enhanced,colback=popgreen,colframe=popink,boxrule=2.2pt,arc=10pt,
        drop shadow=popink,left=11pt,right=11pt,top=10pt,bottom=10pt,height=3.1cm,valign=center]
      \tcbox[colback=white,colframe=popink,boxrule=1.3pt,arc=5pt,boxsep=0pt,nobeforeafter,
        left=3pt,right=3pt,top=3pt,bottom=3pt,valign=center]{\qrcode[height=1.9cm]{https://play.google.com/store/apps/details?id=com.luvvix.onbuch&hl=fr}}%
      \hspace{8pt}\begin{minipage}[c]{0.5\linewidth}
        {\popdisplay\fontsize{11}{12.5}\selectfont\color{white}\MakeUppercase{Télécharge OnBuch}}\par\vspace{4pt}
        {\raggedright\popsemi\fontsize{8.4}{11}\selectfont\color{white}Gratuit \textbullet\ Google Play\\Tuteur IA, annales, concours, résultats\par}
      \end{minipage}
    \end{tcolorbox}
  \end{minipage}\hfill
  \begin{minipage}[t]{0.487\linewidth}
    \begin{tcolorbox}[enhanced,colback=poppurple,colframe=popink,boxrule=2.2pt,arc=10pt,
        drop shadow=popink,left=11pt,right=11pt,top=10pt,bottom=10pt,height=3.1cm,valign=center]
      \tikz[baseline=-0.7ex]\node[circle,fill=white,draw=popink,line width=1.3pt,minimum size=1.6cm,inner sep=0]{\popdisplay\fontsize{24}{24}\selectfont\color{poppurple}+};%
      \hspace{8pt}\begin{minipage}[c]{0.6\linewidth}
        {\popdisplay\fontsize{11}{12.5}\selectfont\color{white}\MakeUppercase{Passe à OnBuch+}}\par\vspace{4pt}
        {\raggedright\popsemi\fontsize{8.4}{11}\selectfont\color{white}Tous les cours signés, Tuteur IA illimité, fiches PDF — onglet OnBuch+ de l'app\par}
      \end{minipage}
    \end{tcolorbox}
  \end{minipage}
  \vspace{8pt}
  \popcontenu
  \vfill
  \signatureonbuch
  \restoregeometry
  \newpage
}

% Rangée « Ce cours contient » (page de garde)
\newcommand{\poptile}[3]{% #1 couleur #2 gros texte #3 légende
  \begin{tcolorbox}[enhanced,colback=#1,colframe=popink,boxrule=2pt,arc=9pt,shadow={2.5pt}{-2.5pt}{0pt}{popink},
      left=6pt,right=6pt,top=9pt,bottom=9pt,halign=center,valign=center,height=1.75cm,width=\linewidth,nobeforeafter]
    {\popdisplay\fontsize{12.5}{13}\selectfont\color{white}\MakeUppercase{#2}}\par\vspace{3pt}
    {\centering\popsemi\fontsize{7.8}{9.5}\selectfont\color{white}#3\par}
  \end{tcolorbox}}
\newcommand{\popcontenu}{%
  \par\noindent{\popxboldls\fontsize{7.5}{7.5}\selectfont\color{popmuted}CE COURS SIGNÉ CONTIENT}\par\vspace{7pt}
  \noindent\begin{minipage}[t]{0.235\linewidth}\poptile{popblue}{Cours}{complet, illustré, pas à pas}\end{minipage}\hfill
  \begin{minipage}[t]{0.235\linewidth}\poptile{poporange}{Méthodes}{et exercices résolus}\end{minipage}\hfill
  \begin{minipage}[t]{0.235\linewidth}\poptile{poppink}{Exercices}{type examen + corrigés}\end{minipage}\hfill
  \begin{minipage}[t]{0.235\linewidth}\poptile{popgreen}{Fiche bilan}{l'essentiel à retenir}\end{minipage}\par}

% Sommaire
\newtcolorbox{sommaire}{enhanced,colback=white,colframe=popink,boxrule=2.2pt,arc=10pt,
  drop shadow=popink,left=18pt,right=18pt,top=13pt,bottom=13pt,before skip=6pt,after skip=20pt,
  before upper={\poppill{Sommaire}{poppurple}{white}\par\vspace{9pt}\popsemi\fontsize{10.6}{14.5}\selectfont\color{popink}}}

% Signature de fin de cours — poussée tout en bas de la dernière page
\newcommand{\findecours}{%
  \par\vfill\nopagebreak
  \begin{center}\popsquiggle{poporange}\end{center}\vspace{4pt}
  \signatureonbuch
}
\AtBeginDocument{\def\llbracket{\mathopen{[\mkern-3.5mu[}}\def\rrbracket{\mathclose{]\mkern-3.5mu]}}} % intervalles d'entiers (glyphe ⟦ absent de la police)
% Glyphes absents de Fira Math : redéfinis à partir de glyphes disponibles
\AtBeginDocument{%
  \def\setminus{\mathbin{\backslash}}\let\smallsetminus\setminus
  \def\vdots{\mathord{\vbox{\baselineskip3.2pt\lineskiplimit0pt\kern4pt\hbox{.}\hbox{.}\hbox{.}}}}%
  \def\ddots{\mathinner{\mkern1mu\raise6.5pt\hbox{.}\mkern2mu\raise3.6pt\hbox{.}\mkern2mu\raise0.7pt\hbox{.}\mkern1mu}}%
  \def\triangle{\mathord{\scalebox{0.9}{$\symup{\Delta}$}}}%
  \def\nabla{\mathord{\raisebox{\depth}{\scalebox{1}[-1]{$\symup{\Delta}$}}}}%
  \def\checkmark{\popcheck{popdark}}%
  \def\star{\mathbin{\ast}}\def\bigstar{\mathord{\ast}}%
  \def\diamond{\mathbin{\scalebox{0.6}{$\Diamond$}}}%
  \def\bigcup{\mathop{\mathchoice{\raisebox{-0.3em}{\scalebox{1.7}{$\cup$}}}{\scalebox{1.25}{$\cup$}}{\cup}{\cup}}\displaylimits}%
  \def\bigcap{\mathop{\mathchoice{\raisebox{-0.3em}{\scalebox{1.7}{$\cap$}}}{\scalebox{1.25}{$\cap$}}{\cap}{\cap}}\displaylimits}%
  \def\lesseqqgtr{\mathrel{\lessgtr}}\def\bigoplus{\mathop{\oplus}\displaylimits}%
}
\providecommand{\ssi}{si et seulement si} % abréviation fréquente
\AtBeginDocument{\providecommand{\wideparen}{\overparen}} % arc de cercle

\begin{document}

\pagegarde{Leçon 37 — Éléments socioculturels : les outils TICs de marques chinoises et camerounaises}{Chinois}{Tle A}{Module 5 : Mass médias et télécommunication}{ZH37}{2 h}{Cette leçon socioculturelle de Terminale A explore les outils TIC de marques chinoises (Huawei, Xiaomi, WeChat, TikTok) et camerounaises (mobile money, startups locales, médias en ligne). Elle fournit le vocabulaire spécialisé en caractères et pinyin, des faits comparés Cameroun/Chine et des activités de réflexion pour exprimer un point de vue simple en chinois.
\vspace{4pt}
\begin{multicols}{2}\small
\begin{itemize}
\item À la fin de cette leçon, je sais nommer en chinois les principaux outils TIC (téléphone, applications, réseaux sociaux, paiement mobile).
\item Je sais citer des marques chinoises présentes au Cameroun (Huawei, Xiaomi, Tecno, TikTok) et des réalités numériques camerounaises.
\item Je sais comparer factuellement les usages numériques du Cameroun et de la Chine.
\item Je sais lire et comprendre un texte simple en chinois sur les TIC.
\item Je sais employer la structure de comparaison 比 pour opposer deux pratiques.
\item Je sais exprimer un point de vue simple avec 我觉得 et 因为.
\end{itemize}\end{multicols}}

\begin{sommaire}
\begin{multicols}{2}
\begin{enumerate}
\item Découverte : les TIC au quotidien
\item Vocabulaire clé du thème
\item Clés et écriture des caractères
\item Texte support et compréhension
\item Comparer et donner son avis
\item Mini-exposé et réflexion
\item Activité d'intégration
\item Méthodes et exercices résolus
\item Exercices
\item Corrigés des exercices
\item Fiche bilan
\end{enumerate}
\end{multicols}
\end{sommaire}

\begin{objectifs}
\begin{itemize}
\item À la fin de cette leçon, je sais nommer en chinois les principaux outils TIC (téléphone, applications, réseaux sociaux, paiement mobile).
\item Je sais citer des marques chinoises présentes au Cameroun (Huawei, Xiaomi, Tecno, TikTok) et des réalités numériques camerounaises.
\item Je sais comparer factuellement les usages numériques du Cameroun et de la Chine.
\item Je sais lire et comprendre un texte simple en chinois sur les TIC.
\item Je sais employer la structure de comparaison 比 pour opposer deux pratiques.
\item Je sais exprimer un point de vue simple avec 我觉得 et 因为.
\item Je sais écrire correctement les caractères clés du thème (电, 话, 网, 机).
\item Je sais préparer et présenter un mini-exposé sur les TIC dans mon pays.
\end{itemize}
\end{objectifs}

\begin{prerequis}
\begin{itemize}
\item Vocabulaire de base de la vie quotidienne (maison, école, famille) acquis en Première
\item Structure sujet-verbe-objet et phrase avec 是
\item Nombres et expressions de quantité (很多, 一些)
\item Verbes modaux simples : 会, 可以, 想
\item Notions sur les médias vues dans le Module 5 (leçons précédentes)
\end{itemize}
\end{prerequis}

\newpage

\coursec{Découverte : les TIC au quotidien}

\courssub{Situation d'accroche : un smartphone dans chaque poche}

Regardez autour de vous dans la salle de classe, à Yaoundé comme à Douala : presque chaque élève tient en main un smartphone Tecno, Infinix, Huawei ou Xiaomi. Ces marques chinoises dominent le marché camerounais, tandis que des services locaux comme le mobile money de MTN (\zh{MTN Mobile Money}) et d'Orange (\zh{Orange Money}) rythment la vie quotidienne : on paie le taxi, on envoie de l'argent à la famille au village, on achète du crédit téléphonique. En Chine, c'est un autre monde numérique : on discute et on paie avec WeChat (\zh{微信}), on achète sur Taobao (\zh{淘宝}), on regarde de courtes vidéos sur Douyin (\zh{抖音}). Au Cameroun, on envoie de l'argent par téléphone et on suit les informations sur Facebook et WhatsApp.

\textbf{Problématique :} comment parler de tout cela en chinois ? Quels mots utiliser pour nommer les technologies, les marques, les applications ? Et comment comparer, de façon factuelle et respectueuse, les habitudes numériques du Cameroun et de la Chine ? Cette leçon nous donne les mots et les faits pour comparer les deux mondes numériques.

\courssub{Observation : un premier texte}

Avant toute définition, lisons ce court texte. Il présente, en deux phrases, le cœur de notre leçon.

\begin{textebox}[Deux mondes numériques]
在中国，很多人用微信聊天和付钱。在喀麦隆，很多人用手机付钱。\\
\emph{Zài Zhōngguó, hěn duō rén yòng Wēixìn liáotiān hé fùqián. Zài Kāmàilóng, hěn duō rén yòng shǒujī fùqián.}\\
« En Chine, beaucoup de gens utilisent WeChat pour discuter et payer. Au Cameroun, beaucoup de gens paient par téléphone. »
\end{textebox}

Observons ce texte :
\begin{itemize}
\item \zh{在} (\emph{zài}) + lieu : « en, à ». Structure : \cle{在 + lieu + ，sujet + verbe}. Exemple : \zh{在喀麦隆} (\emph{zài Kāmàilóng}) « au Cameroun ».
\item \zh{用} (\emph{yòng}) : « utiliser, se servir de ». Structure : \cle{Sujet + 用 + outil + verbe}. Exemple : \zh{用手机付钱} (\emph{yòng shǒujī fùqián}) « payer avec le téléphone ».
\item \zh{付钱} (\emph{fùqián}) : « payer » (littéralement « verser de l'argent »). C'est un verbe dit « séparable » (离合词) : \zh{付} est le verbe, \zh{钱} « argent » est son complément d'objet figé. On peut insérer un élément entre les deux : \zh{付了一百块钱} (\emph{fù le yībǎi kuài qián}) « payer 100 yuans ».
\item \zh{聊天} (\emph{liáotiān}) : « bavarder, discuter » (littéralement « causer du ciel »). Verbe séparable aussi : \zh{聊} (verbe) + \zh{天} (complément). Exemple : \zh{聊了半天天} (\emph{liáo le bàntiān tiān}) « avoir bavardé une demi-journée ».
\end{itemize}

\courssub{Qu'est-ce que les TIC ?}

\begin{definition}[Les TIC en chinois]
Les \cle{TIC} (technologies de l'information et de la communication) se disent en chinois \zh{信息技术} (\emph{xìnxī jìshù}) : \zh{信息} « information » + \zh{技术} « technologie ». C'est l'ensemble des technologies utilisées pour communiquer et traiter l'information : téléphone, internet, ordinateur, applications. Dans la vie courante, on dit aussi plus simplement \zh{网络技术} (\emph{wǎngluò jìshù}, « technologies du réseau ») ou \zh{网络} (\emph{wǎngluò}, « internet, le réseau ») et l'action \zh{上网} (\emph{shàngwǎng}, « aller sur internet, se connecter »).
\end{definition}

\begin{attention}[Ne pas traduire « TIC » mot à mot]
Les francophones veulent souvent traduire « technologies de l'information et de la communication » littéralement, ce qui donne une phrase lourde et artificielle. En chinois, on dit sobrement \zh{信息技术} dans un contexte scolaire ou officiel, et \zh{网络} ou \zh{上网} (\emph{shàngwǎng}) dans la conversation. Retiens aussi : on ne dit pas \zh{*用TIC}, on dit \zh{用网络} « utiliser internet » ou \zh{用手机} « utiliser le téléphone ».
\end{attention}

\courssub{Inventaire de classe : quels outils utilisons-nous ?}

Activité de découverte : chaque élève répond à voix haute à la question \zh{你用什么手机？} (\emph{Nǐ yòng shénme shǒujī ?} « Quel téléphone utilises-tu ? ») et \zh{你用什么软件？} (\emph{Nǐ yòng shénme ruǎnjiàn ?} « Quelles applications utilises-tu ? »). Voici les réponses-types :

\begin{exemplebox}[Réponses possibles]
\zh{我用华为手机。} \emph{Wǒ yòng Huáwéi shǒujī.} « J'utilise un téléphone Huawei. »\\
\zh{我用WhatsApp跟朋友聊天。} \emph{Wǒ yòng WhatsApp gēn péngyou liáotiān.} « J'utilise WhatsApp pour discuter avec mes amis. »\\
\zh{我用手机付钱。} \emph{Wǒ yòng shǒujī fùqián.} « Je paie avec mon téléphone. »\\
\zh{我在Facebook上看新闻。} \emph{Wǒ zài Facebook shàng kàn xīnwén.} « Je lis les informations sur Facebook. »
\end{exemplebox}

Note la structure de la dernière phrase : \cle{Sujet + 在 + lieu/application + 上 + verbe}. \zh{在……上} signifie « sur » (une plateforme, un site) : \zh{在网上} (\emph{zài wǎng shàng}) « sur internet », \zh{在Facebook上} (\emph{zài Facebook shàng}) « sur Facebook ».

\coursec{Clés et écriture des caractères}

\courssub{Radicaux (clés) et ordre des traits des mots-clés}

Pour écrire correctement le vocabulaire de cette leçon, identifions les clés (radicaux) et l'ordre des traits principaux. Rappel : haut → bas, gauche → droite, horizontal avant vertical, trait central avant les ailes.

\begin{center}
\begin{tabularx}{\linewidth}{|X|X|X|X|X|}
\hline
\rowcolor{popblueL} Caractère & Pinyin & Clé (radical) & Nombre de traits & Ordre des traits (résumé) \\
\hline
信 & xìn & \zh{亻} (ren, homme) & 9 & \zh{亻} (2) → \zh{言} (7 : 点, 横, 横, 横, 竖, 撇, 捺) \\
\hline
息 & xī & \zh{心} (xin, cœur) & 10 & \zh{自} (6) → \zh{心} (4 : 点, 点, 竖, 点) \\
\hline
技 & jì & \zh{扌} (shou, main) & 7 & \zh{扌} (3) → \zh{支} (4 : 横, 撇, 竖, 点) \\
\hline
术 & shù & \zh{术} (shu, propre) & 8 & 横, 竖, 横, 撇, 点, 撇, 横, 竖弯钩 \\
\hline
手 & shǒu & \zh{手} (shou, main) & 4 & 横, 竖钩, 提, 点 \\
\hline
机 & jī & \zh{木} (mu, bois) & 6 & \zh{木} (4) → \zh{幂} (2 : 撇, 点) \\
\hline
软 & ruǎn & \zh{车} (che, véhicule) & 8 & \zh{车} (4) → \zh{冖} (2) → \zh{小} (3) → 点 (attention : \zh{车} s'écrit différemment en bas de \zh{软}) \\
\hline
件 & jiàn & \zh{亻} (ren, homme) & 6 & \zh{亻} (2) → \zh{牛} (4 : 横, 竖, 撇, 捺) \\
\hline
付 & fù & \zh{亻} (ren, homme) & 5 & \zh{亻} (2) → \zh{寸} (3 : 横, 撇, 捺) \\
\hline
钱 & qián & \zh{钅} (jin, or/métal) & 10 & \zh{钅} (5) → \zh{戈} (5 : 横, 横, 撇, 横, 竖弯钩) \\
\hline
微 & wēi & \zh{彳} (chi, pas) & 13 & \zh{彳} (3) → \zh{⺌} (5) → \zh{戈} (5) \\
\hline
信 & xìn & (voir ci-dessus) & 9 & (voir ci-dessus) \\
\hline
支 & zhī & \zh{十} (shi, dix) & 4 & 横, 竖, 撇, 点 \\
\hline
付 & fù & (voir ci-dessus) & 5 & (voir ci-dessus) \\
\hline
宝 & bǎo & \zh{宀} (mian, toit) & 8 & \zh{宀} (3) → \zh{玉} (5 : 横, 竖, 横, 横, 点) \\
\hline
\end{tabularx}
\end{center}

\begin{methode}[Écrire les caractères composés : repérer la clé phonétique]
\begin{enumerate}
\item Identifie la \cle{clé de sens} (souvent à gauche ou en haut) : \zh{扌} pour \zh{技} (action de la main), \zh{钅} pour \zh{钱} (métal/monnaie), \zh{宀} pour \zh{宝} (toit/abri).
\item Identifie la \cle{clé phonétique} (souvent à droite ou en bas) qui donne une indication de prononciation : \zh{支} (\emph{zhī}) dans \zh{技} (\emph{jì}) et \zh{支} (\emph{zhī}) ; \zh{戈} (\emph{gē}) dans \zh{微} (\emph{wēi}) — attention, l'indication n'est pas toujours parfaite.
\item Écris la clé de sens, puis la clé phonétique, en respectant l'ordre des traits global (gauche → droite, haut → bas).
\item Entraîne-toi à écrire chaque caractère 5 à 10 fois sur papier quadrillé (田字格) en visualisant la structure.
\end{enumerate}
\end{methode}

\coursec{Vocabulaire clé du thème}

\courssub{Liste de référence (HSK 2--4)}

\begin{center}
\begin{tabularx}{\linewidth}{|X|X|X|X|}
\hline
\rowcolor{popblueL} Caractères & Pinyin & Nature & Traduction \\
\hline
信息技术 & xìnxī jìshù & nom & technologies de l'information (TIC) \\
\hline
网络 & wǎngluò & nom & internet, le réseau \\
\hline
上网 & shàngwǎng & verbe & aller sur internet, se connecter \\
\hline
手机 & shǒujī & nom & téléphone portable \\
\hline
软件 & ruǎnjiàn & nom & application, logiciel \\
\hline
应用 & yìngyòng & nom/verbe & application (app) / appliquer \\
\hline
华为 & Huáwéi & nom propre & Huawei (marque chinoise) \\
\hline
小米 & Xiǎomǐ & nom propre & Xiaomi (litt. « petit millet ») \\
\hline
传音 & Chuányīn & nom propre & Transsion (maison mère de Tecno, Infinix) \\
\hline
中兴 & Zhōngxīng & nom propre & ZTE (marque chinoise) \\
\hline
微信 & Wēixìn & nom propre & WeChat (application chinoise tout-en-un) \\
\hline
支付宝 & Zhīfùbǎo & nom propre & Alipay (service de paiement Alibaba) \\
\hline
抖音 & Dǒuyīn & nom propre & Douyin (version chinoise de TikTok) \\
\hline
淘宝 & Táobǎo & nom propre & Taobao (site d'achats en ligne) \\
\hline
付钱 & fùqián & verbe (sep.) & payer \\
\hline
手机支付 & shǒujī zhīfù & nom & paiement mobile (mobile money) \\
\hline
聊天 & liáotiān & verbe (sep.) & discuter, bavarder \\
\hline
新闻 & xīnwén & nom & informations, nouvelles \\
\hline
社交媒体 & shèjiāo méitǐ & nom & réseaux sociaux, médias sociaux \\
\hline
短视频 & duǎn shìpín & nom & vidéo courte \\
\hline
直播 & zhíbō & nom/verbe & direct (streaming) / diffuser en direct \\
\hline
扫码 & sǎomǎ & verbe & scanner un code QR \\
\hline
二维码 & èrwéimǎ & nom & code QR \\
\hline
流量 & liúliàng & nom & données mobiles (forfait internet) \\
\hline
信号 & xìnhào & nom & signal (réseau) \\
\hline
充值 & chōngzhí & verbe & recharger (crédit, forfait) \\
\hline
转账 & zhuǎnzhàng & verbe & virer de l'argent, faire un virement \\
\hline
\end{tabularx}
\end{center}

\begin{savaistu}[Tecno et Infinix sont chinoises !]
Tecno et Infinix, les marques de téléphones les plus populaires au Cameroun, appartiennent au groupe chinois Transsion (\zh{传音}, \emph{Chuányīn}), coté à la bourse de Shenzhen. Ce groupe a conçu ses téléphones spécialement pour le marché africain : double carte SIM, batterie longue durée, appareil photo adapté à toutes les carnations de peau, prix accessibles. C'est un bel exemple de coopération économique entre la Chine et l'Afrique : des produits chinois pensés pour les besoins camerounais.
\end{savaistu}

\courssub{Carte mentale : le monde des TIC}

\begin{popfigure}
\begin{tikzpicture}[mindmap, grow cyclic, every node/.style={concept, align=center, font=\small},
root concept/.append style={concept color=popblue, font=\small\bfseries},
level 1/.append style={level distance=3.2cm, sibling angle=90},
level 2/.append style={level distance=2.4cm, sibling angle=45, font=\scriptsize}]
\node[root concept, align=center]{TIC\\信息技术};
\node[concept color=poporangeL, align=center] at (90:3.2){Téléphones\\手机}
  child { node[concept color=poporange!40, align=center]{华为 Huawei\\小米 Xiaomi} }
  child { node[concept color=poporange!40, align=center]{传音 Transsion\\(Tecno, Infinix)} };
\node[concept color=popgreenL, align=center] at (0:3.2){Applications\\软件/应用}
  child { node[concept color=popgreen!40, align=center]{Cameroun :\\WhatsApp, Facebook} }
  child { node[concept color=popgreen!40, align=center]{Chine :\\微信 WeChat\\抖音 Douyin} };
\node[concept color=popgoldL, align=center] at (270:3.2){Paiement\\付钱/支付}
  child { node[concept color=popgold!40, align=center]{Cameroun :\\手机支付\\MTN/Orange Money} }
  child { node[concept color=popgold!40, align=center]{Chine :\\微信支付\\支付宝 Alipay} };
\node[concept color=poppurpleL, align=center] at (180:3.2){Médias \& Achats\\媒体/购物}
  child { node[concept color=poppurple!40, align=center]{Cameroun :\\Facebook Marketplace\\Radio/TV en ligne} }
  child { node[concept color=poppurple!40, align=center]{Chine :\\淘宝 Taobao\\京东 JD.com} };
\end{tikzpicture}
\legende{Carte mentale du vocabulaire des TIC : quatre branches (téléphones, applications, paiement, médias/achats) avec les exemples du Cameroun et de la Chine.}
\end{popfigure}

\coursec{Texte support et compréhension}

\courssub{Texte authentique : deux écosystèmes numériques}

Lis le texte suivant. Il utilise le vocabulaire du tableau et les structures \zh{用} et \zh{在……上}.

\begin{textebox}[Le quotidien connecté : Yaoundé -- Pékin]
喀麦隆和中国都进入了数字时代。在喀麦隆的大街小巷，你能看到传音（Tecno, Infinix）、华为和小米的手机店。喀麦隆人喜欢用WhatsApp和Facebook跟家人朋友聊天，看新闻。付钱时，大家习惯用手机支付，比如MTN Mobile Money或Orange Money，不需要银行卡。\\
\emph{Kāmàilóng hé Zhōngguó dōu jìnrù le shùzì shídài. Zài Kāmàilóng de dàjiēxiǎoxiàng, nǐ néng kàndào Chuányīn (Tecno, Infinix), Huáwéi hé Xiǎomǐ de shǒujī diàn. Kāmàilóng rén xǐhuan yòng WhatsApp hé Facebook gēn jiārén péngyou liáotiān, kàn xīnwén. Fùqián shí, dàjiā xíguàn yòng shǒujī zhīfù, bǐrú MTN Mobile Money huò Orange Money, bù xūyào yínhángkǎ.}\\
\\
在中国，微信不仅是聊天软件，还是“钱包”。中国人用微信支付、支付宝买菜、坐地铁、交水电费。购物上淘宝或京东，看短视频刷抖音。不用现金，只带手机就能生活。\\
\emph{Zài Zhōngguó, Wēixìn bùjǐn shì liáotiān ruǎnjiàn, hái shì “qiánbāo”. Zhōngguó rén yòng Wēixìn zhīfù, Zhīfùbǎo mǎicài, zuò dìtiě, jiāo shuǐdiànfèi. Gòuwù shàng Táobǎo huò Jīngdōng, kàn duǎn shìpín shuā Dǒuyīn. Bù yòng xiànjīn, zhǐ dài shǒujī jiù néng shēnghuó.}\\
\\
« Le Cameroun et la Chine sont tous deux entrés dans l'ère numérique. Dans les rues du Cameroun, on voit des boutiques de téléphones Transsion (Tecno, Infinix), Huawei et Xiaomi. Les Camerounais aiment utiliser WhatsApp et Facebook pour discuter avec leur famille et leurs amis, lire les infos. Pour payer, tout le monde a l'habitude d'utiliser le paiement mobile, comme MTN Mobile Money ou Orange Money, pas besoin de carte bancaire. \\
En Chine, WeChat n'est pas seulement un logiciel de discussion, c'est aussi un "portefeuille". Les Chinois utilisent WeChat Pay, Alipay pour acheter des légumes, prendre le métro, payer l'eau et l'électricité. Pour faire les courses, on va sur Taobao ou JD.com, pour regarder des vidéos courtes on scrolle Douyin. Pas de liquide, avec juste le téléphone on peut vivre. »
\end{textebox}

\courssub{Compréhension de l'écrit}

\begin{exercice}[Questions sur le texte]
\begin{enumerate}
\item Selon le texte, quelles marques de téléphones voit-on dans les rues du Cameroun ? Cite-les en chinois.
\item Quelles sont les deux applications principales pour discuter au Cameroun ? Et en Chine ?
\item Au Cameroun, comment paie-t-on généralement ? Donne le terme chinois et deux exemples de services.
\item En Chine, WeChat est décrit comme un « portefeuille » (\zh{钱包}). Quels sont trois exemples de paiements quotidiens cités dans le texte ?
\item Traduction : \zh{不用现金，只带手机就能生活。}
\item Vrai ou Faux ? Justifie par une phrase du texte.
\begin{itemize}
\item[a)] Les Chinois utilisent encore beaucoup l'argent liquide (\zh{现金}).
\item[b)] Au Cameroun, il faut obligatoirement une carte bancaire pour payer par téléphone.
\end{itemize}
\end{enumerate}
\end{exercice}

\begin{corrige}[Exercice : Questions sur le texte]
\begin{enumerate}
\item \zh{传音（Tecno, Infinix）、华为和小米} (Transsion, Huawei, Xiaomi).
\item Cameroun : \zh{WhatsApp和Facebook} ; Chine : \zh{微信} (WeChat).
\item \zh{手机支付} (paiement mobile) : \zh{MTN Mobile Money} et \zh{Orange Money}.
\item Acheter des légumes (\zh{买菜}), prendre le métro (\zh{坐地铁}), payer l'eau et l'électricité (\zh{交水电费}).
\item « Sans utiliser d'argent liquide, en apportant seulement son téléphone on peut vivre. »
\item[a)] Faux. Texte : \zh{不用现金} (on n'utilise pas d'argent liquide). [b)] Faux. Texte : \zh{不需要银行卡} (pas besoin de carte bancaire).
\end{enumerate}
\end{corrige}

\coursec{Comparer et donner son avis}

\courssub{Tableau comparatif : usages Cameroun / Chine}

Chaque pays a ses outils préférés, liés à son histoire, à son infrastructure bancaire et à ses besoins. Voici un tableau factuel pour réviser le vocabulaire et préparer l'expression orale.

\begin{center}
\begin{tabularx}{\linewidth}{|X|X|X|}
\hline
\rowcolor{popblueL} Domaine & Cameroun 喀麦隆 & Chine 中国 \\
\hline
Téléphones courants & Tecno, Infinix, Huawei, Xiaomi & Huawei, Xiaomi, Oppo, Vivo, Honor \\
\hline
Discussion (messagerie) & WhatsApp, Facebook Messenger & 微信 WeChat \\
\hline
Paiement mobile & 手机支付 mobile money (MTN, Orange) & 微信支付 WeChat Pay, 支付宝 Alipay \\
\hline
Achats en ligne & Facebook Marketplace, Jumia, sites web & 淘宝 Taobao, 京东 JD.com, 拼多多 Pinduoduo \\
\hline
Vidéos courtes & TikTok, Instagram Reels & 抖音 Douyin, 快手 Kuaishou \\
\hline
Informations & Radio/TV en ligne, Facebook, WhatsApp & 微信公众号, 今日头条, 网易新闻 \\
\hline
Code QR & Utilisé pour paiement, menus & Omniprésent : paiement, ajout ami, location vélo, métro \\
\hline
\end{tabularx}
\end{center}

\courssub{Structures grammaticales pour comparer (niveau HSK 3)}

Pour exprimer la similitude, la différence ou la préférence, utilise ces structures maîtrisées au programme de Terminale.

\begin{propriete}[Structures de comparaison et d'opinion]
\begin{enumerate}
\item \textbf{Aussi / Également} : \cle{Sujet + 也 + Verbe} \\
\zh{喀麦隆人也用手机付钱。} (\emph{Kāmàilóng rén yě yòng shǒujī fùqián.}) « Les Camerounais paient aussi par téléphone. »
\item \textbf{Pas pareil / Différent} : \cle{A 和/跟 B 不一样} \\
\zh{中国的软件和喀麦隆的软件不一样。} (\emph{Zhōngguó de ruǎnjiàn hé Kāmàilóng de ruǎnjiàn bù yíyàng.}) « Les applis chinoises et camerounaises ne sont pas pareilles. »
\item \textbf{Comparatif "plus que"} : \cle{A 比 B + Adj} (sans \zh{更}) \\
\zh{中国的移动支付比喀麦隆普及。} (\emph{Zhōngguó de yídòng zhīfù bǐ Kāmàilóng pǔjí.}) « Le paiement mobile est plus répandu en Chine qu'au Cameroun. »
\item \textbf{Préférence} : \cle{Sujet + 比较 喜欢/喜欢……多了} \\
\zh{我比较喜欢用微信，因为方便。} (\emph{Wǒ bǐjiào xǐhuan yòng Wēixìn, yīnwèi fāngbiàn.}) « Je préfère utiliser WeChat car c'est pratique. »
\item \textbf{Parce que / Donc} : \cle{因为……，所以……} \\
\zh{因为没有银行卡，所以很多人用手机支付。} (\emph{Yīnwèi méiyǒu yínhángkǎ, suǒyǐ hěn duō rén yòng shǒujī zhīfù.})
\end{enumerate}
\end{propriete}

\begin{exemplebox}[Phrases modèles pour l'oral]
\zh{我觉得喀麦隆的手机支付很方便，因为不需要银行卡。} \\
\emph{Wǒ juéde Kāmàilóng de shǒujī zhīfù hěn fāngbiàn, yīnwèi bù xūyào yínhángkǎ.} \\
« Je trouve que le mobile money camerounais est pratique, car pas besoin de carte bancaire. »\\

\zh{中国的微信功能太多了：聊天、付钱、看新闻、买票，什么都能做。} \\
\emph{Zhōngguó de Wēixìn gōngnéng tài duō le: liáotiān, fùqián, kàn xīnwén, mǎi piào, shénme dōu néng zuò.} \\
« WeChat en Chine a trop de fonctions : discuter, payer, lire les infos, acheter des billets, il peut tout faire. »\\

\zh{虽然品牌一样，但喀麦隆人用WhatsApp，中国人用微信，习惯不一样。} \\
\emph{Suīrán pǐnpái yíyàng, dàn Kāmàilóng rén yòng WhatsApp, Zhōngguó rén yòng Wēixìn, xíguàn bù yíyàng.} \\
« Bien que les marques soient les mêmes, les Camerounais utilisent WhatsApp, les Chinois WeChat, les habitudes ne sont pas les mêmes. »
\end{exemplebox}

\begin{attention}[Erreurs fréquentes des francophones -- Comparaison]
\begin{itemize}
\item \textbf{Oubli de \zh{也} avant le verbe} : \zh{*喀麦隆人用手机付钱也} $\rightarrow$ \zh{喀麦隆人\underline{也}用手机付钱。}
\item \textbf{Mauvaise place de \zh{比}} : \zh{*中国移动支付普及比喀麦隆} $\rightarrow$ \zh{中国的移动支付\underline{比}喀麦隆\underline{普及}。} (Le terme comparé doit suivre \zh{比}).
\item \textbf{Confusion \zh{不一样} vs \zh{不同}} : \zh{不一样} (bù yíyàng) est oral, courant ; \zh{不同} (bù tóng) est plus formel, écrit. Les deux marchent ici.
\item \textbf{Traduction littérale de "préférer"} : Ne dis pas \zh{*我喜欢A比B}. Dis \zh{我比较喜欢A} ou \zh{我喜欢A多了}.
\item \textbf{Marques en caractères} : Écris \zh{华为}, \zh{小米}, \zh{微信}, \zh{支付宝}, \zh{抖音}, \zh{淘宝} en caractères. Garde WhatsApp, Facebook, TikTok, MTN, Orange en lettres latines (usage standard dans les textes chinois contemporains).
\end{itemize}
\end{attention}

\coursec{Mini-exposé et réflexion}

\courssub{Méthode : Préparer un mini-exposé comparatif (2--3 minutes)}

\begin{methode}[Étapes pour réussir ton exposé oral]
\begin{enumerate}
\item \textbf{Choisis un angle précis} : ne dis pas « les TIC » en général. Exemples : « Le paiement mobile : Cameroun vs Chine », « Mes applications préférées », « Les marques de téléphone à Yaoundé ».
\item \textbf{Structure ton discours} (utilise les connecteurs logiques) :\\
\zh{首先} (Premièrement), \zh{其次} (Deuxièmement), \zh{最后} (Enfin) / \zh{总之} (En résumé).\\
\zh{一方面……，另一方面……} (D'un côté…, de l'autre côté…).
\item \textbf{Utilise le vocabulaire de la leçon} : au moins 5 mots du tableau (ex: \zh{手机支付}, \zh{微信}, \zh{扫码}, \zh{二维码}, \zh{传音}, \zh{方便}, \zh{普及}).
\item \textbf{Apporte un fait chiffré ou une observation personnelle} : « Dans mon quartier, 8 sur 10 ont un Tecno » / « En 2023, le taux de pénétration du mobile money au Cameroun est... » (donnée approximative acceptée).
\item \textbf{Conclus par ton avis} : \zh{我觉得……} / \zh{我认为……} / \zh{对我来说，……更适合。}
\item \textbf{Répète à voix haute} 3 fois en chronométrant. Vise 2 minutes. Note les mots qui bloquent (tons, vocabulaire) et révise-les.
\end{enumerate}
\end{methode}

\courssub{Sujets proposés pour le mini-exposé (choisis-en un)}

\begin{enumerate}
\item \zh{我的手机和我常用的应用} (Mon téléphone et mes applis habituelles).
\item \zh{喀麦隆和中国的手机支付：异同点} (Le paiement mobile Cameroun/Chine : similitudes et différences).
\item \zh{为什么微信在中国“无所不能”，而在喀麦隆用WhatsApp？} (Pourquoi WeChat est "tout-puissant" en Chine alors qu'on utilise WhatsApp au Cameroun ?).
\item \zh{中国手机品牌在喀麦隆的成功：以传音为例} (Le succès des marques chinoises au Cameroun : cas Transsion).
\item \zh{短视频时代：抖音 vs TikTok，年轻人的选择} (Ère de la vidéo courte : Douyin vs TikTok, le choix des jeunes).
\end{enumerate}

\begin{experience}[Activité de classe : Speed-dating des exposés]
\begin{enumerate}
\item Les élèves forment deux cercles concentriques (ou deux rangées face à face).
\item Chaque élève a 2 minutes pour présenter son sujet à son vis-à-vis en chinois uniquement.
\item Signal sonore : le cercle extérieur tourne d'une place. Nouveau partenaire.
\item Répéter 4 à 5 fois. L'enseignant circule, note les points forts (vocabulaire, structures \zh{比}, \zh{因为……所以……}, prononciation des tons) et les erreurs récurrentes pour une correction collective finale.
\item Débriefing : 3 élèves volontaires refont leur exposé devant la classe ; la classe note 2 points forts et 1 amélioration.
\end{enumerate}
\end{experience}

\begin{exoresolu}[Traduction thème/version : comparer les écosystèmes]
Énoncé : Traduis en chinois le paragraphe suivant.
« Au Cameroun et en Chine, les gens utilisent tous des smartphones. Mais les applications sont différentes. Au Cameroun, on utilise WhatsApp pour discuter et le mobile money pour payer. En Chine, on utilise WeChat pour tout faire : discuter, payer, faire les courses. Je trouve que WeChat est plus pratique, mais le mobile money camerounais est aussi très utile pour ceux qui n'ont pas de compte en banque. »
\tcblower
Solution détaillée :
\begin{enumerate}
\item \zh{喀麦隆和中国，人们都用智能手机。} (\emph{Kāmàilóng hé Zhōngguó, rénmen dōu yòng zhìnéng shǒujī.}) — \zh{智能手机} (smartphone), \zh{都} (tous) placé avant le verbe.
\item \zh{但是，软件不一样。} (\emph{Dànshì, ruǎnjiàn bù yíyàng.})
\item \zh{在喀麦隆，人们用WhatsApp聊天，用手机支付付钱。} (\emph{Zài Kāmàilóng, rénmen yòng WhatsApp liáotiān, yòng shǒujī zhīfù fùqián.})
\item \zh{在中国，人们用微信做什么都可以：聊天、付钱、买东西。} (\emph{Zài Zhōngguó, rénmen yòng Wēixìn zuò shénme dōu kěyǐ: liáotiān, fùqián, mǎi dōngxi.}) — Structure \zh{什么都可以} « tout est possible / on peut tout faire ».
\item \zh{我觉得微信比较方便，但是喀麦隆的手机支付对没有银行账户的人也很有用。} (\emph{Wǒ juéde Wēixìn bǐjiào fāngbiàn, dànshì Kāmàilóng de shǒujī zhīfù duì méiyǒu yínháng zhànghù de rén yě hěn yǒuyòng.}) — \zh{账户} (compte), \zh{有用} (utile).
\end{enumerate}
\end{exoresolu}

\begin{exercice}[Entraînement complet : écrit, vocabulaire, grammaire]
\begin{enumerate}
\item \textbf{Écriture des caractères} : Écris 5 fois chacun les caractères suivants en respectant l'ordre des traits : \zh{信}, \zh{技}, \zh{付}, \zh{钱}, \zh{微}, \zh{支}, \zh{宝}, \zh{扫}, \zh{码}, \zh{直}.
\item \textbf{Complète les phrases} avec le mot juste du vocabulaire (\zh{软件}, \zh{应用}, \zh{扫码}, \zh{二维码}, \zh{充值}, \zh{转账}, \zh{流量}, \zh{信号}) :
\begin{itemize}
\item[a)] 我手机没\_\_\_\_\_\_了，上不了网。 (plus de data)
\item[b)] 请用微信\_\_\_\_\_\_付款。 (scanner le code)
\item[c)] 这个\_\_\_\_\_\_很方便，可以买火车票。 (cette appli)
\item[d)] 我要给妈妈\_\_\_\_\_\_500元。 (virer de l'argent)
\item[e)] 这里\_\_\_\_\_\_不好，打不通电话。 (le signal)
\end{itemize}
\item \textbf{Transforme les phrases} en utilisant la structure \zh{比} :
\begin{itemize}
\item[a)] 中国的移动支付很普及。喀麦隆的移动支付也普及，但中国更普及。 $\rightarrow$ \zh{中国的移动支付\_\_\_\_\_\_。}
\item[b)] 微信的功能多。WhatsApp的功能少。 $\rightarrow$ \zh{微信的功能\_\_\_\_\_\_。}
\end{itemize}
\item \textbf{Expression écrite} : Rédige un court message (60--80 caractères) à un ami chinois pour lui expliquer comment tu paies tes courses au marché à Douala. Utilise \zh{手机支付}, \zh{扫码}, \zh{方便}, \zh{不用现金}.
\item \textbf{Traduction version} : Traduis en français :
\zh{现在的年轻人，不管在喀麦隆还是在中国，都离不开手机。每天醒来第一件事就是看微信或WhatsApp。网络已经变成了生活的一部分。}
\end{enumerate}
\end{exercice}

\begin{corrige}[Exercice : Entraînement complet]
\begin{enumerate}
\item \textit{Exercice d'écriture manuscrite à faire sur cahier. Vérifie l'ordre des traits avec le tableau de la section "Clés et écriture".}
\item \textbf{Complétion :}
[a)] \zh{流量} (liúliàng) — \zh{我手机没流量了，上不了网。}
[b)] \zh{扫码} (sǎomǎ) — \zh{请用微信扫码付款。}
[c)] \zh{应用} (yìngyòng) ou \zh{软件} (ruǎnjiàn) — \zh{这个应用很方便，可以买火车票。}
[d)] \zh{转账} (zhuǎnzhàng) — \zh{我要给妈妈转账500元。}
[e)] \zh{信号} (xìnhào) — \zh{这里信号不好，打不通电话。}
\item \textbf{Structure \zh{比} :}
[a)] \zh{中国的移动支付比喀麦隆普及。} (ou \zh{比喀麦隆的普及})
[b)] \zh{微信的功能比WhatsApp多。}
\item \textbf{Expression écrite (modèle) :}
\zh{我在杜阿拉买菜都用手机支付，扫一下二维码就行，很方便，不用带现金。} (49 caractères)
\item \textbf{Traduction :}
« Les jeunes d'aujourd'hui, qu'ils soient au Cameroun ou en Chine, ne peuvent pas se passer de leur téléphone. La première chose qu'ils font en se réveillant chaque jour, c'est regarder WeChat ou WhatsApp. Internet est déjà devenu une partie de la vie. »
\end{enumerate}
\end{corrige}

\begin{aretenir}[L'essentiel de la leçon 37 -- TIC : Cameroun / Chine]
\begin{itemize}
\item \textbf{Vocabulaire clé} : \zh{信息技术} (TIC), \zh{网络/上网} (internet), \zh{手机} (smartphone), \zh{软件/应用} (appli), \zh{微信/支付宝/抖音/淘宝} (apps chinoises), \zh{手机支付/扫码/二维码} (paiement mobile), \zh{传音} (Transsion/Tecno/Infinix).
\item \textbf{Structures grammaticales} :
\begin{itemize}
\item \cle{Sujet + 用 + outil + verbe} (\zh{用手机付钱}).
\item \cle{在 + plateforme + 上 + verbe} (\zh{在微信上聊天}).
\item \cle{Sujet + 也 + Verbe} (similitude).
\item \cle{A 比 B + Adj} (comparatif).
\item \cle{因为……，所以……} (cause/conséquence).
\end{itemize}
\item \textbf{Faits socioculturels} :
\begin{itemize}
\item Marques de téléphones : Transsion (Tecno, Infinix), Huawei, Xiaomi dominent le marché camerounais.
\item Paiement : Mobile Money (MTN, Orange) au Cameroun (inclusion financière sans banque) vs WeChat Pay / Alipay en Chine (écosystème sans cash).
\item Applications : WhatsApp/Facebook (Cameroun) vs WeChat « super-app » / Douyin (Chine).
\item Achats : Facebook Marketplace / Jumia (Cameroun) vs Taobao / JD.com / Pinduoduo (Chine).
\end{itemize}
\item \textbf{Compétences visées} : Comparer factuellement, employer le lexique TIC, exprimer un avis simple en chinois (oral/écrit), lire un texte authentique, écrire les caractères clés.
\end{itemize}
\end{aretenir}

\coursec{Vocabulaire clé du thème}

Dans la section précédente, nous avons découvert à quel point les TIC (technologies de l'information et de la communication) sont présentes dans notre quotidien : au marché de Douala, on paie par téléphone ; à Yaoundé, les élèves cherchent des informations sur internet ; en Chine, on commande presque tout avec une seule application. Pour parler de tout cela en chinois, il nous faut maintenant un vocabulaire précis et bien organisé. C'est l'objectif de cette section : apprendre les mots clés du thème, les comprendre, savoir les prononcer avec les bons tons et les employer dans des phrases correctes.

\courssub{Observer d'abord : une phrase pour entrer dans le lexique}

Avant toute liste de mots, observons une phrase très simple, du type de celles que l'on entend chaque jour au Cameroun :

\begin{textebox}[Phrase d'observation]
我的手机是中国品牌，又便宜又好用。\\
Wǒ de shǒujī shì Zhōngguó pǐnpái, yòu piányi yòu hǎoyòng.\\
« Mon téléphone est une marque chinoise, à la fois bon marché et agréable à utiliser. »
\end{textebox}

Cette seule phrase contient déjà quatre mots clés de notre leçon : \zh{手机} (téléphone portable), \zh{品牌} (marque), \zh{便宜} (bon marché) et la structure \zh{又……又……} (à la fois... et...). Remarquez aussi l'ordre des mots : \zh{我的手机} « mon téléphone » — le possessif \zh{我的} se place avant le nom, comme en français, mais avec la particule \zh{的} entre les deux. Gardons cette phrase en tête : nous y reviendrons.

\courssub{Le glossaire du thème}

Voici le vocabulaire officiel de la leçon, à connaître activement (savoir lire, prononcer, écrire et employer).

\begin{center}
\begin{tabularx}{\linewidth}{|X|X|X|X|}
\hline
\rowcolor{popblueL} Caractères & Pinyin & Nature & Traduction \\
\hline
手机 & shǒujī & n. & téléphone portable \\
\hline
电脑 & diànnǎo & n. & ordinateur \\
\hline
网络 & wǎngluò & n. & réseau, internet \\
\hline
上网 & shàngwǎng & v. & aller sur internet \\
\hline
应用程序 & yìngyòng chéngxù & n. & application \\
\hline
软件 & ruǎnjiàn & n. & logiciel \\
\hline
微信 & Wēixìn & n. propre & WeChat \\
\hline
支付宝 & Zhīfùbǎo & n. propre & Alipay \\
\hline
抖音 & Dǒuyīn & n. propre & Douyin / TikTok \\
\hline
手机支付 & shǒujī zhīfù & n. & paiement mobile \\
\hline
付钱 & fùqián & v. & payer \\
\hline
品牌 & pǐnpái & n. & marque \\
\hline
公司 & gōngsī & n. & entreprise, société \\
\hline
社交媒体 & shèjiāo méitǐ & n. & réseaux sociaux \\
\hline
消息 & xiāoxi & n. & message, information \\
\hline
方便 & fāngbiàn & adj. & pratique \\
\hline
便宜 & piányi & adj. & bon marché \\
\hline
流行 & liúxíng & adj. & populaire, à la mode \\
\hline
\end{tabularx}
\end{center}

\begin{definition}[Les marques chinoises à connaître]
\zh{华为} Huáwéi (Huawei) et \zh{小米} Xiǎomǐ (Xiaomi) sont deux grandes entreprises chinoises de téléphones et d'appareils électroniques, très présentes sur le marché camerounais. \zh{小米} se traduit littéralement par « petit millet » (une céréale) : un nom simple et facile à retenir pour une marque mondiale !
\end{definition}

\courssub{Comprendre les mots de l'intérieur : les familles de caractères}

Le chinois n'est pas une langue de mots isolés : la plupart des mots sont composés de deux caractères qui apportent chacun un sens. Comprendre ces « briques » permet de deviner le sens de mots nouveaux. Observons :

\begin{itemize}
\item \zh{电} diàn = électricité → \zh{电话} diànhuà « parole électrique » = téléphone ; \zh{电脑} diànnǎo « cerveau électrique » = ordinateur ; \zh{电视} diànshì « vision électrique » = télévision.
\item \zh{网} wǎng = filet, réseau → \zh{网络} wǎngluò = réseau, internet ; \zh{上网} shàngwǎng « monter sur le réseau » = aller sur internet.
\item \zh{付} fù = payer → \zh{付钱} fùqián « payer de l'argent » ; \zh{支付} zhīfù = payer, paiement ; \zh{支付宝} Zhīfùbǎo « trésor de paiement » = Alipay.
\item \zh{信} xìn = message, lettre → \zh{微信} Wēixìn « micro-message » = WeChat ; \zh{信息} xìnxī = information.
\end{itemize}

\begin{methode}[Mémoriser par réseaux de mots]
Pour retenir durablement ce vocabulaire, n'apprenez pas les mots un par un dans le désordre. Procédez ainsi :
\begin{enumerate}
\item Choisissez un caractère-clé (par exemple \zh{电} diàn, « électricité »).
\item Listez tous les mots de la leçon qui le contiennent : \zh{电脑}, \zh{电话}, \zh{电视}.
\item Fabriquez une phrase avec chaque mot : \zh{我用电脑学习。} Wǒ yòng diànnǎo xuéxí. « J'étudie avec l'ordinateur. »
\item Recommencez avec \zh{网} (\zh{网络}, \zh{上网}), puis \zh{付} (\zh{付钱}, \zh{支付}, \zh{支付宝}).
\item Le lendemain, testez-vous : cachez le pinyin et la traduction, lisez les caractères seuls.
\end{enumerate}
Apprendre par famille est trois fois plus efficace qu'apprendre par liste alphabétique : le cerveau retient les réseaux, pas les listes.
\end{methode}

\begin{popfigure}
\begin{center}
\begin{tikzpicture}[
  mindmap,
  concept color=popblue,
  every node/.style={concept, align=center},
  level 1 concept/.append style={sibling angle=90, level distance=3.2cm},
  level 2 concept/.append style={sibling angle=50, level distance=2.8cm},
  grow cyclic
]
\node[concept color=poporange, font=\Large\bfseries, align=center]{手机\\ shǒujī}
  child[concept color=popblue] { node[align=center]{应用\\ yìngyòng}
    child { node[align=center]{软件\\ ruǎnjiàn} }
    child { node[align=center]{微信\\ Wēixìn} }
  }
  child[concept color=popgreen] { node[align=center]{支付\\ zhīfù}
    child { node[align=center]{付钱\\ fùqián} }
    child { node[align=center]{支付宝\\ Zhīfùbǎo} }
  }
  child[concept color=poppurple] { node[align=center]{社交\\ shèjiāo}
    child { node[align=center]{社交媒体\\ shèjiāo méitǐ} }
    child { node[align=center]{抖音\\ Dǒuyīn} }
  }
  child[concept color=poppink] { node[align=center]{信息\\ xìnxī}
    child { node[align=center]{消息\\ xiāoxi} }
    child { node[align=center]{上网\\ shàngwǎng} }
  };
\end{tikzpicture}
\end{center}
\legende{Carte mentale lexicale : le smartphone (\zh{手机}) au centre de quatre familles de mots — applications, paiement, réseaux sociaux, informations.}
\end{popfigure}

\courssub{Employer le vocabulaire : exemples et structures}

Un mot n'est vraiment appris que s'il est employé dans une phrase. Observons les structures les plus utiles de la leçon.

\begin{exemplebox}[Exemples traduits à étudier]
\zh{我每天上网两个小时。}\\
\emph{Wǒ měitiān shàngwǎng liǎng ge xiǎoshí.}\\
« Je vais sur internet deux heures par jour. »\\[4pt]
\zh{这个应用很方便。}\\
\emph{Zhège yìngyòng hěn fāngbiàn.}\\
« Cette application est très pratique. »\\[4pt]
\zh{我用手机付钱。}\\
\emph{Wǒ yòng shǒujī fùqián.}\\
« Je paie avec mon téléphone. »\\[4pt]
\zh{华为是一个很有名的中国公司。}\\
\emph{Huáwéi shì yí ge hěn yǒumíng de Zhōngguó gōngsī.}\\
« Huawei est une entreprise chinoise très connue. »\\[4pt]
\zh{现在手机支付在喀麦隆越来越流行。}\\
\emph{Xiànzài shǒujī zhīfù zài Kāmàilóng yuè lái yuè liúxíng.}\\
« Aujourd'hui, le paiement mobile est de plus en plus populaire au Cameroun. »\\[4pt]
\zh{她用微信给朋友发消息。}\\
\emph{Tā yòng Wēixìn gěi péngyou fā xiāoxi.}\\
« Elle envoie des messages à ses amis avec WeChat. »
\end{exemplebox}

\begin{propriete}[Trois structures à retenir]
\begin{enumerate}
\item \textbf{Sujet + 用 + outil + verbe} : \zh{用} yòng signifie « utiliser, avec ». \zh{我用电脑学习。} « J'étudie avec l'ordinateur. » L'outil se place \emph{avant} le verbe, contrairement au français (« j'étudie \emph{avec} l'ordinateur »).
\item \textbf{又 + adj. + 又 + adj.} : « à la fois... et... ». \zh{这个手机又便宜又好用。} « Ce téléphone est à la fois bon marché et agréable à utiliser. » Réservée aux adjectifs ; ne pas l'utiliser avec des verbes d'action.
\item \textbf{越来越 + adjectif} : « de plus en plus... ». \zh{越来越流行} « de plus en plus populaire ». Ne jamais mettre \zh{很} après \zh{越来越}.
\end{enumerate}
\end{propriete}

\begin{attention}[Pièges fréquents des francophones]
\begin{itemize}
\item \textbf{Prononciation de 便宜} : ce mot se prononce \emph{piányi} (deuxième ton + ton neutre), et non \emph{biànyí}. C'est l'erreur de prononciation la plus fréquente chez les francophones : le caractère \zh{便} se lit \emph{biàn} dans \zh{方便} fāngbiàn, mais \emph{pián} dans \zh{便宜} piányi. Retenez : « fāng\textbf{biàn} » mais « \textbf{pián}yi ».
\item \textbf{上网 est un verbe séparable} : on dit \zh{上网两个小时} « aller sur internet pendant deux heures », la durée se place entre les deux parties du verbe ou après : \zh{我每天上网两个小时。} Ne dites pas « shàngwǎng le internet ».
\item \textbf{消息 est invariable} : pas de pluriel en chinois. \zh{一个消息} « un message », \zh{很多消息} « beaucoup de messages ».
\item \textbf{公司 vs 品牌} : \zh{公司} est l'entreprise (Huawei est une \zh{公司}), \zh{品牌} est la marque (Huawei est aussi une \zh{品牌}). Une entreprise peut avoir plusieurs marques.
\end{itemize}
\end{attention}

\begin{savaistu}[微信, le « micro-message » devenu géant]
\zh{微信} Wēixìn signifie littéralement « micro-message » (\zh{微} wēi = micro, minuscule ; \zh{信} xìn = message). Lancée en 2011 par l'entreprise Tencent, cette application est devenue bien plus qu'une messagerie : en Chine, on s'en sert pour discuter, payer, commander un taxi, prendre rendez-vous chez le médecin et même payer ses factures. De même, \zh{支付宝} Zhīfùbǎo signifie « trésor (\zh{宝}) de paiement (\zh{支付}) ». Au Cameroun, ce sont les services de paiement mobile des opérateurs locaux qui jouent un rôle comparable : deux pays, deux solutions, un même besoin — payer facilement avec son téléphone.
\end{savaistu}

\courssub{Tableau comparatif : dire la même chose en français et en chinois}

\begin{center}
\begin{tabularx}{\linewidth}{|X|X|X|}
\hline
\rowcolor{popblueL} Structure & Exemple en chinois + pinyin & Traduction \\
\hline
用 + outil + verbe & 我用手机付钱。Wǒ yòng shǒujī fùqián. & Je paie avec mon téléphone. \\
\hline
又 + adj. + 又 + adj. & 这个应用又便宜又方便。Zhège yìngyòng yòu piányi yòu fāngbiàn. & Cette application est à la fois bon marché et pratique. \\
\hline
越来越 + adj. & 微信越来越流行。Wēixìn yuè lái yuè liúxíng. & WeChat est de plus en plus populaire. \\
\hline
在 + lieu + verbe & 我在网上学习汉语。Wǒ zài wǎngshàng xuéxí Hànyǔ. & J'apprends le chinois sur internet. \\
\hline
给 + personne + 发消息 & 他给我发消息。Tā gěi wǒ fā xiāoxi. & Il m'envoie un message. \\
\hline
\end{tabularx}
\end{center}

\courssub{Exercice résolu et entraînement}

\begin{exoresolu}[Traduire en chinois]
Énoncé — Traduisez en chinois, avec le pinyin : « Mon téléphone est une marque chinoise. Il est à la fois bon marché et pratique. Je paie avec mon téléphone tous les jours. »
\tcblower
Solution détaillée — On procède phrase par phrase.\\
1. « Mon téléphone » = \zh{我的手机} ; « est une marque chinoise » = \zh{是中国品牌} (pas d'article en chinois). → \zh{我的手机是中国品牌。} Wǒ de shǒujī shì Zhōngguó pǐnpái.\\
2. « À la fois bon marché et pratique » → structure \zh{又……又……} : \zh{又便宜又方便}. → \zh{它又便宜又方便。} Tā yòu piányi yòu fāngbiàn.\\
3. « Je paie avec mon téléphone » → structure \zh{用 + outil + verbe} : \zh{我用手机付钱} ; « tous les jours » = \zh{每天}, placé après le sujet. → \zh{我每天用手机付钱。} Wǒ měitiān yòng shǒujī fùqián.\\
Vérification : tons de \zh{便宜} (piányi), place de \zh{每天} avant le verbe, particule \zh{的} dans \zh{我的手机}.
\end{exoresolu}

\begin{exercice}[À vous de jouer]
\begin{enumerate}
\item Traduisez en français : \zh{小米的手机在喀麦隆很流行。} (Xiǎomǐ de shǒujī zài Kāmàilóng hěn liúxíng.)
\item Complétez avec le bon mot (\zh{上网}, \zh{付钱}, \zh{消息}, \zh{品牌}) : \zh{我用手机给朋友发\_\_\_\_\_\_。}
\item Traduisez en chinois : « Cette entreprise est très connue, ses applications sont très pratiques. »
\item Classez ces mots en deux familles (\zh{电} et \zh{网}) : \zh{电脑}, \zh{网络}, \zh{电话}, \zh{上网}, \zh{电视}.
\end{enumerate}
\end{exercice}

\begin{corrige}[Exercice « À vous de jouer »]
1. « Les téléphones Xiaomi sont très populaires au Cameroun. »\\
2. \zh{消息} : \zh{我用手机给朋友发消息。} « J'envoie des messages à mes amis avec mon téléphone. »\\
3. \zh{这个公司很有名，它的应用很方便。} Zhège gōngsī hěn yǒumíng, tā de yìngyòng hěn fāngbiàn.\\
4. Famille \zh{电} : \zh{电脑}, \zh{电话}, \zh{电视}. Famille \zh{网} : \zh{网络}, \zh{上网}.
\end{corrige}

\begin{aretenir}[L'essentiel du vocabulaire]
\begin{itemize}
\item Les mots clés : \zh{手机} téléphone portable, \zh{电脑} ordinateur, \zh{网络} internet, \zh{应用} application, \zh{品牌} marque, \zh{公司} entreprise, \zh{付钱} payer, \zh{方便} pratique, \zh{便宜} bon marché, \zh{流行} populaire.
\item \zh{便宜} se prononce \emph{piányi}, jamais \emph{biànyí}.
\item Apprenez par familles de caractères : \zh{电}, \zh{网}, \zh{付}, \zh{信}.
\item Trois structures utiles : \zh{用 + outil + verbe}, \zh{又……又……}, \zh{越来越 + adjectif}.
\end{itemize}
\end{aretenir}

Ce vocabulaire est désormais votre boîte à outils. Dans la prochaine section, nous étudierons les clés (radicaux) et l'écriture des caractères les plus importants de cette liste, trait par trait.

\coursec{Clés et écriture des caractères}

Dans la section précédente, nous avons découvert le vocabulaire clé des outils TIC : \zh{手机} (téléphone portable), \zh{电脑} (ordinateur), \zh{网络} (réseau), \zh{付钱} (payer). Vous avez remarqué que certains caractères reviennent dans plusieurs mots : \zh{电} dans \zh{电话} et \zh{电脑}, \zh{机} dans \zh{手机} et \zh{飞机}. Cette section va vous apprendre à \emph{comprendre} ces caractères au lieu de les apprendre par cœur : grâce aux \cle{clés} (部首 \emph{bùshǒu}) et à l'\cle{ordre des traits}, écrire le chinois devient un travail logique, presque mathématique.

\courssub{Qu'est-ce qu'une clé (部首) ?}

\begin{definition}[La clé — 部首 bùshǒu]
Une \cle{clé} (部首 \emph{bùshǒu}, littéralement « tête de section ») est un \textbf{composant graphique} qui sert à \textbf{classer} le caractère dans les dictionnaires et qui \textbf{suggère son sens}. Exemples : la clé 讠(\zh{言} simplifié, « parole ») se trouve dans \zh{话} (parole), \zh{说} (parler), \zh{语} (langue) ; la clé 木 (« bois ») se trouve dans \zh{机} (machine), \zh{树} (arbre), \zh{桌} (table).
\end{definition}

Observez d'abord ces mots avant de lire la règle :

\begin{exemplebox}[Observation : la clé donne une piste de sens]
\zh{说话} \emph{shuōhuà} « parler » ; \zh{电话} \emph{diànhuà} « téléphone » ; \zh{汉语} \emph{Hànyǔ} « langue chinoise » ; \zh{法语} \emph{Fǎyǔ} « langue française ».\par
Que remarquez-vous ? Tous ces mots contiennent des caractères avec la clé 讠, et tous ont un rapport avec la \textbf{parole} ou la \textbf{langue}. La clé fonctionne comme une étiquette de sens.
\end{exemplebox}

\begin{propriete}[Règle de lecture d'un caractère composé]
La plupart des caractères chinois se décomposent en deux parties :
\begin{itemize}
\item la \cle{clé} (souvent à gauche ou en haut) qui indique le \textbf{domaine de sens} ;
\item la \cle{phonétique} (souvent à droite) qui suggère la \textbf{prononciation}.
\end{itemize}
Exemple : \zh{话} = 讠(clé : parole) + 舌 \emph{shé} (phonétique : langue, qui donne le son proche de \emph{huà}). De même \zh{活} = 氵(clé : eau) + 舌 → \emph{huó} « vivre ». Même phonétique 舌, clés différentes, sens différents.
\end{propriete}

\begin{popfigure}
\centering
\begin{tikzpicture}[mindmap, grow cyclic, every node/.style={concept, align=center}, root concept/.append style={concept color=popblue, font=\Large\bfseries}, level 1/.append style={level distance=3.2cm, sibling angle=72}, level 1 concept/.append style={concept color=popblueL, font=\normalsize}]
\node[align=center]{讠\\ \footnotesize clé de la parole}
  child { node[align=center]{话 huà\\ parole} }
  child { node[align=center]{说 shuō\\ parler} }
  child { node[align=center]{语 yǔ\\ langue} }
  child { node[align=center]{读 dú\\ lire} }
  child { node[align=center]{谢 xiè\\ remercier} };
\end{tikzpicture}
\legende{Carte mentale : la famille de la clé 讠(parole). Tous ces verbes et noms ont un rapport avec le langage.}
\end{popfigure}

\courssub{Les cinq caractères de la leçon, analysés un par un}

\textbf{1. 电 diàn — « électricité »}\par
Clé : 田/曰 (le caractère est classé sous la clé 曰 « dire » dans certains dictionnaires, mais graphiquement on y voit le champ 田 traversé par un éclair 乚). Nombre de traits : \textbf{5}. Le caractère représente à l'origine un éclair dans le ciel au-dessus d'un champ — d'où le sens « électricité ». Famille de mots : \zh{电话} \emph{diànhuà} « téléphone » (électricité + parole), \zh{电脑} \emph{diànnǎo} « ordinateur » (électricité + cerveau), \zh{电视} \emph{diànshì} « télévision » (électricité + regarder).

\textbf{2. 话 huà — « parole »}\par
Clé : 讠(parole) à gauche ; phonétique : 舌 \emph{shé} « langue » à droite. Nombre de traits : \textbf{8}. Famille : \zh{说话} \emph{shuōhuà} « parler », \zh{电话} \emph{diànhuà} « téléphone », \zh{汉语} \emph{Hànyǔ} « le chinois » (littéralement « langue des Han »).

\textbf{3. 网 wǎng — « filet, réseau »}\par
Clé : 冂 (l'« enceinte » extérieure). Nombre de traits : \textbf{6}. Le caractère dessine un filet : un cadre extérieur avec deux croisillons à l'intérieur. Famille : \zh{网络} \emph{wǎngluò} « réseau », \zh{上网} \emph{shàngwǎng} « aller sur Internet ». Au Cameroun comme en Chine, on dit \zh{我每天晚上上网。} \emph{Wǒ měitiān wǎnshang shàngwǎng.} « Je vais sur Internet tous les soirs. »

\textbf{4. 机 jī — « machine »}\par
Clé : 木 (bois) à gauche — les premières machines étaient en bois ! Phonétique : 几 \emph{jī/jǐ} à droite. Nombre de traits : \textbf{6}. Famille : \zh{手机} \emph{shǒujī} « téléphone portable », \zh{飞机} \emph{fēijī} « avion » (littéralement « machine qui vole »), \zh{机器} \emph{jīqì} « machine ».

\begin{exemplebox}[Le mot 手机 décortiqué]
\zh{手机} \emph{shǒujī} = 手 \emph{shǒu} « main » + 机 \emph{jī} « machine » → littéralement « \textbf{machine à main} » : le téléphone portable !\par
De même : \zh{电脑} \emph{diànnǎo} = 电 « électricité » + 脑 \emph{nǎo} « cerveau » → « \textbf{cerveau électrique} » : l'ordinateur.\par
Le chinois construit ses mots modernes avec des blocs de sens très concrets — c'est une excellente nouvelle pour l'apprenant : chaque mot nouveau se devine souvent à partir de ses composants.
\end{exemplebox}

\textbf{5. 付 fù — « payer »}\par
Clé : 亻(l'homme, variante de 人) à gauche ; à droite 寸 \emph{cùn} « pouce », qui évoque la main qui mesure et donne. Nombre de traits : \textbf{5}. Famille : \zh{付钱} \emph{fùqián} « payer (de l'argent) », \zh{支付} \emph{zhīfù} « payer, régler » — le mot que l'on voit dans \zh{支付宝} \emph{Zhīfùbǎo} « Alipay ». Exemple : \zh{我用手机付钱。} \emph{Wǒ yòng shǒujī fùqián.} « Je paie avec mon téléphone. »

\begin{center}
\begin{tabularx}{\linewidth}{|X|X|X|X|}
\hline
\rowcolor{popblueL} Caractère & Clé & Traits & Famille de mots \\
\hline
电 diàn (électricité) & 田/曰 & 5 & 电话, 电脑, 电视 \\
\hline
话 huà (parole) & 讠+ 舌 & 8 & 说话, 电话, 汉语 \\
\hline
网 wǎng (filet/réseau) & 冂 & 6 & 网络, 上网 \\
\hline
机 jī (machine) & 木 + 几 & 6 & 手机, 飞机, 机器 \\
\hline
付 fù (payer) & 亻+ 寸 & 5 & 付钱, 支付 \\
\hline
\end{tabularx}
\end{center}

\courssub{L'ordre des traits : les règles d'or}

\begin{methode}[Comment écrire un caractère chinois correctement]
\begin{enumerate}
\item \textbf{De haut en bas} : on commence par le trait le plus haut.
\item \textbf{De gauche à droite} : la partie gauche avant la partie droite (ex. \zh{话} : d'abord 讠, ensuite 舌).
\item \textbf{L'extérieur avant l'intérieur} : on trace d'abord le cadre, puis le contenu, et on « ferme » à la fin (ex. \zh{网} : le cadre 冂 d'abord, les croisillons ensuite).
\item \textbf{L'horizontal avant le vertical} quand deux traits se croisent.
\item \textbf{Le centre avant les côtés} pour les caractères symétriques.
\end{enumerate}
Pourquoi respecter cet ordre ? Parce qu'il rend l'écriture plus rapide, plus régulière, et qu'il est indispensable pour utiliser les dictionnaires par nombre de traits et les claviers d'écriture manuscrite sur téléphone.
\end{methode}

\textbf{Démonstration au tableau : 电 (5 traits)}\par
1) 丨 le vertical gauche du champ ; 2) (trait brisé) horizontal-vertical qui forme le haut et la droite du champ ; 3) 一 l'horizontal médian ; 4) 一... non : le vertical central qui traverse ; 5) 乚 le grand crochet final qui sort du champ — c'est « l'éclair ». Retenez : \emph{on dessine le champ 田, puis l'éclair 乚 sort en bas à droite}.

\textbf{Démonstration au tableau : 网 (6 traits)}\par
1) 丨 vertical gauche du cadre ; 2) (trait brisé) horizontal-vertical-crochet qui forme le haut et la droite ; 3) ノ premier tombant intérieur ; 4) 丶 premier point ; 5) ノ deuxième tombant ; 6) 丶 deuxième point. Règle appliquée : \textbf{l'extérieur avant l'intérieur}.

\textbf{Démonstration au tableau : 话 (8 traits)}\par
D'abord la clé 讠(2 traits : 丶 puis ㇊), puis 舌 (6 traits : ノ, 一, 一, 丨, (trait brisé), 一). Règle appliquée : \textbf{la gauche avant la droite}.

\begin{popfigure}
\centering
\begin{tikzpicture}[font=\Large, node distance=0.35cm]
% 电 : ordre des traits
\node[align=center] (d1) at (0,0) {\zh{电}};
\node[below=0.1cm of d1, font=\footnotesize, align=center] {丨 → (trait brisé) → 一\\ → 丨 → 乚};
\draw[->, thick, popblue] (1.1,0) -- (1.9,0);
\node[align=center, text=popblue] (d2) at (2.9,0) {5 traits\\ \footnotesize de haut en bas};
% 网 : ordre des traits
\node[align=center] (w1) at (5.4,0) {\zh{网}};
\node[below=0.1cm of w1, font=\footnotesize, align=center] {丨 → (trait brisé) → ノ\\ → 丶 → ノ → 丶};
\draw[->, thick, poporange] (6.5,0) -- (7.3,0);
\node[align=center, text=poporange] (w2) at (8.4,0) {6 traits\\ \footnotesize extérieur\\ \footnotesize avant intérieur};
\end{tikzpicture}
\legende{Ordre des traits de 电 et 网 : on construit le cadre d'abord, on remplit ensuite.}
\end{popfigure}

\courssub{Pièges graphiques : ne pas confondre}

\begin{attention}[话 huà ≠ 活 huó — et 机 jī ≠ 饥 jī]
\begin{itemize}
\item \zh{话} \emph{huà} « parole » : clé 讠(parole) + 舌. \zh{活} \emph{huó} « vivre » : clé 氵(eau) + 舌. \textbf{Même phonétique 舌, clés différentes} : la parole appartient au langage, la vie (comme l'eau qui coule) appartient au domaine de l'eau. Comparez : \zh{我说汉语。} \emph{Wǒ shuō Hànyǔ.} « Je parle chinois. » / \zh{他住在雅温得，生活很好。} \emph{Tā zhù zài Yǎwēndé, shēnghuó hěn hǎo.} « Il vit à Yaoundé, sa vie est agréable. »
\item \zh{机} \emph{jī} « machine » : clé 木 (bois). \zh{饥} \emph{jī} « avoir faim » : clé 饣(nourriture). Même prononciation, sens opposés ! \zh{手机} « portable » n'a rien à voir avec la faim.
\end{itemize}
Conseil de méthode : quand deux caractères se ressemblent, \textbf{regardez la clé} — elle trahit toujours le sens.
\end{attention}

\begin{exoresolu}[Identifier clés et familles]
Énoncé. Pour chaque caractère, indiquez la clé, le nombre de traits et un mot de sa famille : 电, 话, 网, 机, 付. Puis traduisez : \zh{我用电脑上网。}
\tcblower
Solution détaillée.
\begin{itemize}
\item 电 : clé 田/曰, 5 traits ; famille : 电话 « téléphone », 电脑 « ordinateur ».
\item 话 : clé 讠, 8 traits ; famille : 说话 « parler », 汉语 « chinois (langue) ».
\item 网 : clé 冂, 6 traits ; famille : 网络 « réseau », 上网 « aller sur Internet ».
\item 机 : clé 木, 6 traits ; famille : 手机 « portable », 飞机 « avion ».
\item 付 : clé 亻, 5 traits ; famille : 付钱 « payer », 支付 « régler ».
\end{itemize}
Traduction : \zh{我用电脑上网。} \emph{Wǒ yòng diànnǎo shàngwǎng.} « Je vais sur Internet avec mon ordinateur. » Structure : Sujet 我 + 用 « avec » + instrument 电脑 + verbe 上网. Remarquez que \zh{上网} est un verbe séparable : on peut dire \zh{上了两个小时的网} « surfer deux heures ».
\end{exoresolu}

\begin{exercice}[À vous d'écrire]
\begin{enumerate}
\item Écrivez 电, 话, 网 cinq fois chacun en respectant l'ordre des traits vu au tableau.
\item Classez ces caractères selon leur clé : 说, 机, 语, 付, 活, 读, 饥, 他.
\item Traduisez en français : \zh{我妈妈用手机付钱。} ; \zh{喀麦隆的学生喜欢上网。}
\end{enumerate}
\end{exercice}

\begin{corrige}[Exercice « À vous d'écrire »]
1) Vérification en classe au tableau : 电 (丨, (trait brisé), 一, 丨, 乚) ; 话 (讠puis 舌) ; 网 (cadre 冂 puis ノ丶ノ丶).\par
2) Clé 讠(parole) : 说, 语, 读 ; clé 木 : 机 ; clé 亻(homme) : 付, 他 ; clé 氵(eau) : 活 ; clé 饣(nourriture) : 饥.\par
3) \zh{我妈妈用手机付钱。} \emph{Wǒ māma yòng shǒujī fùqián.} « Ma maman paie avec son téléphone portable. » ; \zh{喀麦隆的学生喜欢上网。} \emph{Kāmàilóng de xuésheng xǐhuan shàngwǎng.} « Les élèves du Cameroun aiment aller sur Internet. »
\end{corrige}

\begin{aretenir}[L'essentiel de la section]
\begin{itemize}
\item Une \cle{clé} (部首) classe le caractère et suggère son sens ; la phonétique suggère le son.
\item 电 (5 traits, électricité), 话 (8 traits, clé 讠, parole), 网 (6 traits, clé 冂, réseau), 机 (6 traits, clé 木, machine), 付 (5 traits, clé 亻, payer).
\item Ordre des traits : haut → bas, gauche → droite, extérieur → intérieur.
\item Ne jamais confondre 话 \emph{huà} (parole) et 活 \emph{huó} (vivre), ni 机 \emph{jī} (machine) et 饥 \emph{jī} (faim) : la clé fait toute la différence.
\end{itemize}
\end{aretenir}

\coursec{Texte support et compréhension}

Dans la section précédente, nous avons appris à écrire les caractères clés du thème des TIC : \zh{手机} (shǒujī, téléphone portable), \zh{品牌} (pǐnpái, marque), \zh{应用} (yìngyòng, application) et \zh{方便} (fāngbiàn, pratique), en respectant l'ordre des traits et en identifiant leurs clés : le radical \zh{扌} (main) dans \zh{手} (de \zh{手机}), le radical \zh{口} (bouche) dans \zh{品} (de \zh{品牌}), le radical \zh{广} (abri) dans \zh{应} (de \zh{应用}) et le radical \zh{辶} (marche) dans \zh{方} (de \zh{方便}). Il est temps maintenant de remettre ces mots dans leur élément naturel : un texte authentique. Lire un texte en chinois, c'est vérifier que les caractères que l'on sait écrire, on sait aussi les reconnaître et les comprendre dans une phrase réelle.

\courssub{Présentation du texte}

Le texte que nous allons lire s'intitule \zh{手机和我们的生活} (\emph{Shǒujī hé wǒmen de shēnghuó}), « Le téléphone portable et notre vie ». Il compare, de façon simple et factuelle, l'usage des outils TIC au Cameroun et en Chine : les marques de téléphones, les applications de paiement et de messagerie. Avant de lire, observez le titre : \zh{和} (hé) signifie « et », \zh{的} (de) est le mot-outil de la possession et de la relation. Le titre suit donc le schéma : A + 和 + B + 的 + nom, « A et notre B ».

\begin{textebox}[手机和我们的生活 — Shǒujī hé wǒmen de shēnghuó]
现在，手机对我们的生活很重要。\\
\emph{Xiànzài, shǒujī duì wǒmen de shēnghuó hěn zhòngyào.}\\
« Aujourd'hui, le téléphone portable est très important dans notre vie. »\\[4pt]
在喀麦隆，很多人用中国品牌的手机，比如华为和小米，因为这些手机又便宜又好。\\
\emph{Zài Kāmàilóng, hěn duō rén yòng Zhōngguó pǐnpái de shǒujī, bǐrú Huáwéi hé Xiǎomǐ, yīnwèi zhèxiē shǒujī yòu piányi yòu hǎo.}\\
« Au Cameroun, beaucoup de gens utilisent des téléphones de marques chinoises, comme Huawei et Xiaomi, parce que ces téléphones sont à la fois bon marché et de bonne qualité. »\\[4pt]
我们也用手机付钱、发消息、看新闻。\\
\emph{Wǒmen yě yòng shǒujī fùqián, fā xiāoxi, kàn xīnwén.}\\
« Nous utilisons aussi le téléphone pour payer, envoyer des messages et lire les informations. »\\[4pt]
在中国，人们常常用微信聊天，用支付宝付钱。\\
\emph{Zài Zhōngguó, rénmen chángcháng yòng Wēixìn liáotiān, yòng Zhīfùbǎo fùqián.}\\
« En Chine, les gens discutent souvent avec WeChat et paient avec Alipay. »\\[4pt]
虽然两个国家用的应用不一样，但是手机都让生活更方便。\\
\emph{Suīrán liǎng ge guójiā yòng de yìngyòng bù yíyàng, dànshì shǒujī dōu ràng shēnghuó gèng fāngbiàn.}\\
« Bien que les applications utilisées dans les deux pays soient différentes, le téléphone rend partout la vie plus pratique. »
\end{textebox}

\begin{popfigure}
\begin{center}
\begin{tabularx}{\linewidth}{|X|X|X|X|}
\hline
\rowcolor{popblueL} \textbf{Caractères} & \textbf{Pinyin} & \textbf{Nature} & \textbf{Traduction française} \\
\hline
现在，手机对我们的生活很重要。 & Xiànzài, shǒujī duì wǒmen de shēnghuó hěn zhòngyào. & Phrase & Aujourd'hui, le téléphone est très important dans notre vie. \\
\hline
很多人用中国品牌的手机。 & Hěn duō rén yòng Zhōngguó pǐnpái de shǒujī. & Phrase & Beaucoup de gens utilisent des téléphones de marques chinoises. \\
\hline
这些手机又便宜又好。 & Zhèxiē shǒujī yòu piányi yòu hǎo. & Phrase & Ces téléphones sont à la fois bon marché et bons. \\
\hline
我们也用手机付钱。 & Wǒmen yě yòng shǒujī fùqián. & Phrase & Nous utilisons aussi le téléphone pour payer. \\
\hline
人们常常用微信聊天。 & Rénmen chángcháng yòng Wēixìn liáotiān. & Phrase & Les gens discutent souvent avec WeChat. \\
\hline
手机让生活更方便。 & Shǒujī ràng shēnghuó gèng fāngbiàn. & Phrase & Le téléphone rend la vie plus pratique. \\
\hline
\end{tabularx}
\end{center}
\legende{Le texte découpé en phrases : caractères, pinyin, nature, traduction}
\end{popfigure}

\courssub{Méthode de lecture : comment aborder un texte chinois}

Face à un texte en caractères, le réflexe du mauvais lecteur est de paniquer devant le premier mot inconnu. Le bon lecteur, lui, procède par étapes.

\begin{methode}[Lire un texte chinois en quatre étapes]
\begin{enumerate}
\item \textbf{Lecture globale (skimming).} Lisez le texte une première fois sans dictionnaire. Repérez les mots que vous connaissez déjà : ici \zh{手机}, \zh{中国}, \zh{用}, \zh{很}, \zh{好}. Demandez-vous simplement : « De quoi parle ce texte ? »
\item \textbf{Deviner par le contexte.} Pour chaque mot nouveau, observez son entourage. Dans \zh{又便宜又好}, le mot \zh{便宜} est encadré par \zh{又……又……} « à la fois… et… » et associé à \zh{好} « bon » : on devine une qualité positive (bon marché).
\item \textbf{Lecture détaillée.} Relevez les noms propres (marques, pays, applications) et les mots-outils grammaticaux (\zh{因为}, \zh{虽然……但是……}) qui donnent la logique du texte.
\item \textbf{Vérification.} Relisez la traduction ou le glossaire, puis relisez le texte chinois à voix haute en respectant les tons.
\end{enumerate}
\end{methode}

\courssub{Lecture 1 : compréhension globale}

Après une première lecture, vous devez pouvoir répondre à la question : \textbf{de quoi parle le texte ?}

Le texte parle de l'importance du téléphone portable (\zh{手机}) dans la vie quotidienne, au Cameroun et en Chine. L'idée principale est contenue dans la première et la dernière phrase : le téléphone est important et rend la vie plus pratique (\zh{更方便}) dans les deux pays, même si les applications utilisées sont différentes.

\begin{exercice}[Vrai ou faux ?]
Lisez le texte et dites si les phrases suivantes sont vraies (对, duì) ou fausses (不对, bú duì).
\begin{enumerate}
\item \zh{手机对我们的生活不重要。} (Le téléphone n'est pas important dans notre vie.)
\item \zh{在喀麦隆，很多人用华为和小米的手机。} (Au Cameroun, beaucoup utilisent des téléphones Huawei et Xiaomi.)
\item \zh{在中国，人们用支付宝聊天。} (En Chine, les gens discutent avec Alipay.)
\item \zh{两个国家用的应用一样。} (Les deux pays utilisent les mêmes applications.)
\end{enumerate}
\end{exercice}

\begin{corrige}[Vrai ou faux]
\begin{enumerate}
\item \textbf{不对} (faux) : le texte dit \zh{手机对我们的生活很重要}, le téléphone est \emph{très important}.
\item \textbf{对} (vrai) : \zh{很多人用中国品牌的手机，比如华为和小米}.
\item \textbf{不对} (faux) : on paie avec Alipay (\zh{用支付宝付钱}), on discute avec WeChat (\zh{用微信聊天}). Attention à ne pas intervertir les fonctions des deux applications.
\item \textbf{不对} (faux) : \zh{两个国家用的应用不一样}, les applications sont \emph{différentes} (\zh{不一样}).
\end{enumerate}
\end{corrige}

\courssub{Lecture 2 : relever les marques et les applications}

La deuxième lecture est une lecture de repérage : il s'agit de retrouver dans le texte tous les noms propres et le vocabulaire clé pour les classer. C'est une compétence essentielle : en chinois, les noms propres étrangers sont transcrits phonétiquement en caractères, et il faut apprendre à les reconnaître.

\begin{center}
\begin{tabularx}{\linewidth}{|X|X|X|X|}
\hline
\rowcolor{popblueL} Caractères & Pinyin & Nature & Traduction \\
\hline
华为 & Huáwéi & nom propre & Huawei (marque chinoise) \\
\hline
小米 & Xiǎomǐ & nom propre & Xiaomi (marque chinoise) \\
\hline
微信 & Wēixìn & nom propre & WeChat (application de messagerie) \\
\hline
支付宝 & Zhīfùbǎo & nom propre & Alipay (application de paiement) \\
\hline
喀麦隆 & Kāmàilóng & nom propre & le Cameroun \\
\hline
品牌 & pǐnpái & nom & marque \\
\hline
应用 & yìngyòng & nom & application \\
\hline
消息 & xiāoxi & nom & message, nouvelle \\
\hline
新闻 & xīnwén & nom & informations, actualités \\
\hline
付钱 & fùqián & verbe & payer \\
\hline
聊天 & liáotiān & verbe & bavarder, discuter \\
\hline
方便 & fāngbiàn & adjectif & pratique, commode \\
\hline
\end{tabularx}
\end{center}

\begin{popfigure}
\begin{center}
\begin{tikzpicture}[mindmap, grow cyclic, every node/.style={concept, align=center, font=\small},
concept color=popblueL, text=popdark,
level 1/.append style={level distance=2.9cm, sibling angle=72},
root concept/.append style={concept color=popblue, text=white, font=\bfseries}]
\node[root concept, align=center]{手机\\ shǒujī\\ TIC}
child { node[concept color=poporangeL, align=center]{品牌\\ pǐnpái\\ 华为, 小米} }
child { node[concept color=popgreenL, align=center]{付钱\\ fùqián\\ 支付宝} }
child { node[concept color=poppinkL, align=center]{聊天\\ liáotiān\\ 微信} }
child { node[concept color=popgoldL, align=center]{看新闻\\ kàn xīnwén\\ informations} }
child { node[concept color=poppurpleL, align=center]{发消息\\ fā xiāoxi\\ messages} };
\end{tikzpicture}
\end{center}
\legende{Carte mentale lexicale : les usages du téléphone dans le texte}
\end{popfigure}

Remarquez la logique des transcriptions chinoises : \zh{微信} (Wēixìn) signifie littéralement « micro-message », c'est le nom chinois d'origine de WeChat ; \zh{支付宝} (Zhīfùbǎo) signifie « trésor de paiement » (\zh{支付} payer + \zh{宝} trésor). Les marques Huawei \zh{华为} et Xiaomi \zh{小米} (« petit millet », une céréale !) gardent leur nom chinois d'origine dans le monde entier.

\courssub{Grammaire en contexte : deux structures du texte}

\textbf{1. La concession avec \zh{虽然……但是……} (suīrán… dànshì…), « bien que… mais… ».}

Observons d'abord la dernière phrase du texte : \zh{虽然两个国家用的应用不一样，但是手机都让生活更方便。} On y trouve deux idées opposées : les applications sont différentes (concession), mais le résultat est le même (idée principale).

\begin{propriete}[La structure 虽然……但是……]
Structure : \cle{虽然} + proposition concédée, \cle{但是} + proposition principale.\\
Emploi : exprimer une opposition, comme « bien que… (mais)… » en français.\\
Différence majeure avec le français : en français, on dit « bien que… » \emph{ou} « mais… », jamais les deux ensemble. En chinois, les deux mots se combinent naturellement dans la même phrase, et c'est même la forme la plus courante. On peut omettre \zh{但是}, mais garder les deux est toujours correct.
\end{propriete}

\begin{exemplebox}[虽然……但是…… en action]
\zh{虽然应用不一样，但是都很方便。}\\
\emph{Suīrán yìngyòng bù yíyàng, dànshì dōu hěn fāngbiàn.}\\
« Bien que les applications soient différentes, elles sont toutes pratiques. »\\[4pt]
\zh{虽然华为手机不贵，但是质量很好。}\\
\emph{Suīrán Huáwéi shǒujī bú guì, dànshì zhìliàng hěn hǎo.}\\
« Bien que les téléphones Huawei ne soient pas chers, leur qualité est très bonne. »\\[4pt]
\zh{虽然我不会说汉语，但是我喜欢中国文化。}\\
\emph{Suīrán wǒ bú huì shuō Hànyǔ, dànshì wǒ xǐhuan Zhōngguó wénhuà.}\\
« Bien que je ne sache pas parler chinois, j'aime la culture chinoise. »
\end{exemplebox}

\textbf{2. La préposition \zh{对} (duì) : « pour, envers ».}

Observons la première phrase : \zh{手机对我们的生活很重要。} Littéralement : « téléphone / pour / notre vie / très important ».

\begin{propriete}[La structure 对……很 + adjectif]
Structure : A + \cle{对} + B + 很 + adjectif.\\
Emploi : exprimer l'attitude ou l'effet de A envers B : « A est important pour B », « A est bon/mauvais pour B ».\\
Exemple : \zh{运动对身体很好。} \emph{Yùndòng duì shēntǐ hěn hǎo.} « Le sport est bon pour la santé. »
\end{propriete}

\begin{attention}[Ne dites jamais 是很重要对 !]
Les francophones traduisent mot à mot « le téléphone \emph{est très important pour} notre vie » et produisent la phrase fautive \zh{手机是很重要对我们的生活}. C'est impossible en chinois. La préposition \zh{对} et son complément se placent \textbf{avant} l'adjectif : \zh{手机对我们的生活很重要。} Retenez l'ordre chinois : Sujet + 对 + complément + 很 + adjectif.
\end{attention}

\courssub{Questions ouvertes et traduction guidée}

\begin{exercice}[Questions de compréhension en chinois simple]
Répondez en phrases complètes, en chinois.
\begin{enumerate}
\item \zh{在喀麦隆，很多人用什么品牌的手机？} (Quelles marques de téléphones utilise-t-on au Cameroun ?)
\item \zh{在中国，人们用什么付钱？} (Avec quoi paie-t-on en Chine ?)
\item \zh{你用手机做什么？} (Que fais-tu avec ton téléphone ?)
\end{enumerate}
\end{exercice}

\begin{corrige}[Questions de compréhension]
\begin{enumerate}
\item \zh{在喀麦隆，很多人用中国品牌的手机，比如华为和小米。} (\emph{Zài Kāmàilóng, hěn duō rén yòng Zhōngguó pǐnpái de shǒujī, bǐrú Huáwéi hé Xiǎomǐ.} « Au Cameroun, beaucoup utilisent des marques chinoises, comme Huawei et Xiaomi. »)
\item \zh{在中国，人们用支付宝付钱。} (\emph{Zài Zhōngguó, rénmen yòng Zhīfùbǎo fùqián.} « En Chine, on paie avec Alipay. »)
\item Réponse personnelle, modèle : \zh{我用手机发消息、看新闻。} (\emph{Wǒ yòng shǒujī fā xiāoxi, kàn xīnwén.} « J'envoie des messages et je lis les informations avec mon téléphone. »)
\end{enumerate}
\end{corrige}

\begin{methode}[Traduire une phrase du chinois vers le français]
\begin{enumerate}
\item Découpez la phrase en groupes de sens : sujet / verbe / compléments.
\item Identifiez les mots-outils (\zh{对}, \zh{因为}, \zh{虽然……但是……}) : ils donnent la structure logique.
\item Traduisez d'abord littéralement, puis reformulez en français naturel en replaçant les mots dans l'ordre français.
\end{enumerate}
\end{methode}

\begin{exoresolu}[Traduction guidée de deux phrases clés]
Énoncé. Traduisez en français : (1) \zh{因为这些手机又便宜又好。} (2) \zh{虽然两个国家用的应用不一样，但是手机都让生活更方便。}
\tcblower
Solution détaillée.\\
(1) Découpage : \zh{因为} (parce que) + \zh{这些手机} (ces téléphones) + \zh{又……又……} (à la fois… et…) + \zh{便宜} (bon marché) + \zh{好} (bon). Traduction littérale : « parce que ces téléphones à la fois bon marché et bons ». Traduction naturelle : « parce que ces téléphones sont à la fois bon marché et de bonne qualité ».\\
(2) Découpage : \zh{虽然} (bien que) + \zh{两个国家用的应用} (les applications utilisées dans les deux pays ; notez \zh{用的} : « que l'on utilise ») + \zh{不一样} (ne sont pas pareilles) + \zh{但是} (mais) + \zh{手机} (le téléphone) + \zh{都} (tous, dans les deux cas) + \zh{让} (rendre, faire que) + \zh{生活} (la vie) + \zh{更方便} (plus pratique). Traduction naturelle : « Bien que les applications utilisées dans les deux pays soient différentes, le téléphone rend partout la vie plus pratique. »
\end{exoresolu}

\begin{aretenir}[L'essentiel de la section]
\begin{itemize}
\item Lire un texte chinois se fait en étapes : mots connus d'abord, mots nouveaux devinés par le contexte ensuite.
\item \zh{虽然……但是……} : en chinois, « bien que » et « mais » coexistent dans la même phrase, contrairement au français.
\item \zh{对……很重要} : la préposition \zh{对} et son complément précèdent l'adjectif ; jamais l'ordre français « est important pour ».
\item Le vocabulaire TIC du texte : \zh{品牌} marque, \zh{应用} application, \zh{付钱} payer, \zh{聊天} discuter, \zh{方便} pratique.
\end{itemize}
\end{aretenir}

\coursec{Comparer et donner son avis}

Dans la section précédente, nous avons lu et compris un texte sur les outils TIC chinois et camerounais. Nous avons appris que les élèves de Yaoundé utilisent WhatsApp et le mobile money, tandis que leurs camarades de Pékin utilisent WeChat et Alipay. Maintenant, une question naturelle se pose : comment dire en chinois que « WeChat a plus de fonctions que WhatsApp », que « le paiement mobile est plus pratique que l'argent liquide » ? Comment dire « je pense que... » et « parce que... donc... » ? C'est exactement l'objet de cette section : apprendre à \cle{comparer} et à \cle{donner son avis} en chinois, avec respect et sans hiérarchie entre les pays.

\courssub{Observer d'abord : deux élèves comparent leurs applications}

Avant toute règle, observons. Voici un mini-dialogue entre deux élèves de Terminale, Amina (Yaoundé) et Lǐ Míng (Pékin), qui échangent par visioconférence dans le cadre d'un partenariat scolaire.

\begin{textebox}[Dialogue — 你最喜欢什么应用？]
Amina : 李明，你最喜欢什么应用？\\
(Lǐ Míng, nǐ zuì xǐhuan shénme yìngyòng ?)\\
« Lǐ Míng, quelle application préfères-tu ? »\\
李明 : 我觉得微信比WhatsApp功能多，因为微信可以聊天，也可以付款。\\
(Wǒ juéde Wēixìn bǐ WhatsApp gōngnéng duō, yīnwèi Wēixìn kěyǐ liáotiān, yě kěyǐ fùkuǎn.)\\
« Je pense que WeChat a plus de fonctions que WhatsApp, parce qu'avec WeChat on peut discuter et aussi payer. »\\
Amina : 有意思！可是在喀麦隆，WhatsApp比微信常用。\\
(Yǒu yìsi! Kěshì zài Kāmàilóng, WhatsApp bǐ Wēixìn chángyòng.)\\
« Intéressant ! Mais au Cameroun, WhatsApp est plus utilisé que WeChat. »\\
李明 : 手机支付呢？\\
(Shǒujī zhīfù ne?)\\
« Et le paiement mobile ? »\\
Amina : 我觉得手机支付很方便，因为不用带现金，所以市场、学校和出租车都能用。\\
(Wǒ juéde shǒujī zhīfù hěn fāngbiàn, yīnwèi bú yòng dài xiànjīn, suǒyǐ shìchǎng, xuéxiào hé chūzūchē dōu néng yòng.)\\
« Je trouve le paiement mobile très pratique, parce qu'on n'a pas besoin d'argent liquide, donc on peut l'utiliser au marché, à l'école et dans les taxis. »\\
李明 : 中国也一样！支付宝比现金方便多了。\\
(Zhōngguó yě yíyàng! Zhīfùbǎo bǐ xiànjīn fāngbiàn duō le.)\\
« En Chine c'est pareil ! Alipay est bien plus pratique que l'argent liquide. »
\end{textebox}

Relevons les phrases importantes du dialogue :

\begin{itemize}
\item \zh{微信比WhatsApp功能多。} (Wēixìn bǐ WhatsApp gōngnéng duō.) « WeChat a plus de fonctions que WhatsApp. »
\item \zh{WhatsApp比微信常用。} (WhatsApp bǐ Wēixìn chángyòng.) « WhatsApp est plus utilisé que WeChat. »
\item \zh{我觉得手机支付很方便。} (Wǒ juéde shǒujī zhīfù hěn fāngbiàn.) « Je trouve le paiement mobile très pratique. »
\item \zh{因为不用带现金，所以市场、学校和出租车都能用。} (Yīnwèi bú yòng dài xiànjīn, suǒyǐ shìchǎng, xuéxiào hé chūzūchē dōu néng yòng.) « Parce qu'on n'a pas besoin d'argent liquide, donc on peut l'utiliser au marché, à l'école et dans les taxis. »
\end{itemize}

Que remarquez-vous ? Le mot \zh{比} (bǐ) se place \emph{entre} les deux éléments comparés, et l'adjectif vient \emph{après}. C'est l'inverse du français, où l'adjectif précède le « que » : « WeChat est \emph{plus pratique que} WhatsApp » se dit mot à mot « WeChat \zh{比} WhatsApp pratique ».

\courssub{La comparaison avec 比 : A 比 B + adjectif}

\begin{propriete}[La structure de comparaison avec 比]
Pour dire « A est plus ... que B », le chinois utilise la structure :

\textbf{A + 比 + B + adjectif}

\begin{itemize}
\item \zh{比} (bǐ) est une préposition qui signifie « comparé à » ; elle introduit le terme de comparaison B.
\item L'adjectif se place \emph{après} B, sans « plus » : la supériorité est déjà contenue dans la structure.
\item \textbf{Exception capitale} : dans cette structure, l'adjectif n'est \emph{jamais} précédé de \zh{很} (hěn). On ne dit pas \zh{比……很贵}.
\item Pour renforcer, on peut ajouter \zh{得多} (de duō, « beaucoup plus »), \zh{多了} (duō le) ou \zh{一点儿} (yìdiǎnr, « un peu plus ») \emph{après} l'adjectif : \zh{支付宝比现金方便得多。} « Alipay est beaucoup plus pratique que l'argent liquide. »
\end{itemize}
\end{propriete}

\begin{popfigure}
\begin{tikzpicture}[node distance=0.5cm]
\node[draw=popblue, fill=popblueL, rounded corners, minimum width=2.6cm, minimum height=1cm, align=center] (A) {\textbf{A}\\华为手机};
\node[draw=poporange, fill=poporangeL, rounded corners, minimum width=1.4cm, minimum height=1cm, align=center, right=of A] (bi) {\textbf{比}\\bǐ};
\node[draw=popblue, fill=popblueL, rounded corners, minimum width=2.6cm, minimum height=1cm, align=center, right=of bi] (B) {\textbf{B}\\苹果手机};
\node[draw=popgreen, fill=popgreenL, rounded corners, minimum width=2.4cm, minimum height=1cm, align=center, right=of B] (adj) {\textbf{adjectif}\\便宜};
\draw[-{Stealth}, thick, popdark] (A) -- (bi);
\draw[-{Stealth}, thick, popdark] (bi) -- (B);
\draw[-{Stealth}, thick, popdark] (B) -- (adj);
\node[below=0.7cm of bi, align=center, text=popmuted] (trad) {Huáwéi shǒujī bǐ píngguǒ shǒujī piányi.\\« Le téléphone Huawei est moins cher que l'iPhone. »};
\draw[dashed, popmuted] (bi.south) -- (trad.north);
\end{tikzpicture}
\legende{Schéma de la structure A + 比 + B + adjectif, avec un exemple annoté.}
\end{popfigure}

\begin{exemplebox}[Comparaisons sur le thème des TIC]
\begin{itemize}
\item \zh{华为手机比苹果手机便宜。} (Huáwéi shǒujī bǐ píngguǒ shǒujī piányi.) « Le téléphone Huawei est moins cher que l'iPhone. »
\item \zh{微信比WhatsApp功能多。} (Wēixìn bǐ WhatsApp gōngnéng duō.) « WeChat a plus de fonctions que WhatsApp. »
\item \zh{网络比以前快多了。} (Wǎngluò bǐ yǐqián kuài duō le.) « Internet est beaucoup plus rapide qu'avant. »
\item \zh{电脑比手机大。} (Diànnǎo bǐ shǒujī dà.) « L'ordinateur est plus grand que le téléphone. »
\item \zh{在喀麦隆，WhatsApp比微信常用。} (Zài Kāmàilóng, WhatsApp bǐ Wēixìn chángyòng.) « Au Cameroun, WhatsApp est plus utilisé que WeChat. »
\end{itemize}
\end{exemplebox}

\begin{attention}[L'erreur n° 1 des francophones : l'ordre des mots]
En français, on dit « le téléphone Huawei est \emph{moins cher que} l'iPhone » : l'adjectif vient \emph{avant} le « que ». Beaucoup d'élèves transposent cet ordre et produisent :

\zh{华为手机很便宜比苹果手机。} — \textbf{FAUX !}

Deux erreurs se cachent ici : (1) l'adjectif est placé avant \zh{比}, alors qu'il doit venir \emph{après} B ; (2) l'adjectif est précédé de \zh{很}, interdit dans la structure de comparaison. La phrase correcte est :

\zh{华为手机比苹果手机便宜。} (Huáwéi shǒujī bǐ píngguǒ shǒujī piányi.)

\textbf{Astuce} : mémorisez le schéma comme une formule mathématique : \cle{A 比 B + adj.}, jamais autrement.
\end{attention}

\courssub{La comparaison négative : 没有 et 不如}

Comment dire « A n'est pas aussi ... que B » ? Le chinois offre deux structures, toutes deux très fréquentes.

\begin{propriete}[Nier la comparaison : 没有 et 不如]
\begin{enumerate}
\item \textbf{A + 没有 + B + adjectif} : « A n'est pas aussi ... que B » (constat neutre).\\
\zh{我的手机没有你的手机新。} (Wǒ de shǒujī méiyǒu nǐ de shǒujī xīn.) « Mon téléphone n'est pas aussi récent que le tien. »
\item \textbf{A + 不如 + B + adjectif} : « A n'égale pas B en ... », « A est moins ... que B » (légèrement évaluatif). L'adjectif peut même être sous-entendu.\\
\zh{我的汉语不如你的汉语好。} (Wǒ de Hànyǔ bùrú nǐ de Hànyǔ hǎo.) « Mon chinois est moins bon que le tien. »\\
\zh{这个功能不如那个功能。} (Zhège gōngnéng bùrú nàge gōngnéng.) « Cette fonction n'égale pas celle-là. »
\end{enumerate}
\textbf{Remarque} : avec \zh{没有} et \zh{不如}, l'adjectif n'est pas précédé de \zh{很} non plus.
\end{propriete}

\begin{exemplebox}[Comparaisons négatives]
\begin{itemize}
\item \zh{这个功能没有那个功能好用。} (Zhège gōngnéng méiyǒu nàge gōngnéng hǎoyòng.) « Cette fonction n'est pas aussi pratique que celle-là. »
\item \zh{我们学校的网络不如中国学校的网络快。} (Wǒmen xuéxiào de wǎngluò bùrú Zhōngguó xuéxiào de wǎngluò kuài.) « Le réseau de notre école est moins rapide que celui des écoles chinoises. »
\item \zh{发短信没有打电话快。} (Fā duǎnxìn méiyǒu dǎ diànhuà kuài.) « Envoyer un SMS n'est pas aussi rapide que téléphoner. »
\end{itemize}
\end{exemplebox}

\courssub{Exprimer son opinion : 我觉得 et 我认为}

Pour donner son avis, deux verbes sont indispensables :

\begin{definition}[我觉得 et 我认为]
\begin{itemize}
\item \zh{我觉得} (wǒ juéde) : « je pense que », « je trouve que » — le plus courant à l'oral, ton personnel et spontané.
\item \zh{我认为} (wǒ rènwéi) : « je considère que », « j'estime que » — plus formel, utilisé dans les exposés et les débats.
\end{itemize}
Les deux introduisent une proposition complète : \zh{我觉得 / 我认为 + Sujet + Verbe}.
\end{definition}

\begin{exemplebox}[Donner son avis sur les TIC]
\begin{itemize}
\item \zh{我觉得手机支付很方便。} (Wǒ juéde shǒujī zhīfù hěn fāngbiàn.) « Je trouve le paiement mobile très pratique. »
\item \zh{我认为学习汉语很有用。} (Wǒ rènwéi xuéxí Hànyǔ hěn yǒuyòng.) « Je considère qu'apprendre le chinois est très utile. »
\item \zh{我觉得中国的应用很有意思。} (Wǒ juéde Zhōngguó de yìngyòng hěn yǒu yìsi.) « Je trouve les applications chinoises très intéressantes. »
\end{itemize}
\textbf{Attention} : après \zh{觉得}, l'adjectif seul est en général précédé de \zh{很} (phrase descriptive normale) : \zh{很方便}. C'est seulement dans la structure avec \zh{比} que \zh{很} est interdit.
\end{exemplebox}

\courssub{Justifier : 因为……所以……}

Une opinion sans justification n'est qu'une impression. Le chinois, comme le français, articule cause et conséquence, mais avec une particularité remarquable : il emploie souvent \emph{les deux} connecteurs ensemble.

\begin{propriete}[La paire 因为……所以……]
\textbf{因为 + cause，所以 + conséquence} : « parce que ..., donc ... ».

\zh{因为不用带现金，所以手机支付很方便。} (Yīnwèi bú yòng dài xiànjīn, suǒyǐ shǒujī zhīfù hěn fāngbiàn.) « Parce qu'on n'a pas besoin d'argent liquide, le paiement mobile est donc très pratique. »

\textbf{Différence avec le français} : en français, « parce que » et « donc » ne s'emploient jamais ensemble dans la même phrase ; en chinois, la paire \zh{因为……所以……} est non seulement correcte mais très naturelle. On peut toutefois n'en garder qu'un seul : \zh{不用带现金，所以手机支付很方便。}
\end{propriete}

\begin{popfigure}
\begin{tikzpicture}[node distance=1.1cm and 1.6cm]
\node[draw=poppurple, fill=poppurpleL, rounded corners, minimum width=4.2cm, minimum height=1.1cm, align=center] (op) {\textbf{Opinion}\\我觉得手机支付很方便。};
\node[draw=poporange, fill=poporangeL, rounded corners, minimum width=4.2cm, minimum height=1.1cm, align=center, below left=of op] (raison) {\textbf{Raison}\\因为不用带现金。};
\node[draw=popgreen, fill=popgreenL, rounded corners, minimum width=4.2cm, minimum height=1.1cm, align=center, below right=of op] (concl) {\textbf{Conclusion}\\所以人人都喜欢用。};
\draw[-{Stealth}, very thick, popdark] (op.south west) ++(0.3,0) -- (raison.north);
\draw[-{Stealth}, very thick, popdark] (raison.east) -- (concl.west);
\draw[-{Stealth}, very thick, popdark, dashed] (op.south east) ++(-0.3,0) -- (concl.north);
\end{tikzpicture}
\legende{Arbre de justification : l'opinion (我觉得) s'appuie sur une raison (因为) qui mène à une conclusion (所以).}
\end{popfigure}

\begin{exemplebox}[Opinions justifiées — le modèle à imiter]
\begin{itemize}
\item \zh{我觉得手机支付很方便，因为不用带现金。} (Wǒ juéde shǒujī zhīfù hěn fāngbiàn, yīnwèi bú yòng dài xiànjīn.) « Je trouve le paiement mobile pratique parce qu'on n'a pas besoin d'argent liquide. »
\item \zh{我认为WhatsApp很好用，因为喀麦隆的每个人都用它。} (Wǒ rènwéi WhatsApp hěn hǎoyòng, yīnwèi Kāmàilóng de měi ge rén dōu yòng tā.) « Je pense que WhatsApp est très pratique parce qu'au Cameroun tout le monde l'utilise. »
\item \zh{因为网络越来越快，所以我们可以上中文课。} (Yīnwèi wǎngluò yuè lái yuè kuài, suǒyǐ wǒmen kěyǐ shàng Zhōngwén kè.) « Parce qu'Internet est de plus en plus rapide, nous pouvons donc suivre des cours de chinois (en ligne). »
\end{itemize}
\end{exemplebox}

\courssub{Comparer avec respect : Cameroun et Chine}

Comparer n'est pas classer. Dire \zh{微信比WhatsApp功能多} ne signifie pas que WeChat est « meilleur » en absolu : au Cameroun, WhatsApp répond parfaitement aux besoins des familles, des commerçants et des élèves. Les usages dépendent des habitudes, des réseaux disponibles et des besoins de chaque pays. En chinois comme en français, on peut marquer ce respect par des formules de nuance :

\begin{itemize}
\item \zh{各有各的好处。} (Gè yǒu gè de hǎochu.) « Chacun a ses avantages. »
\item \zh{两个国家的情况不一样。} (Liǎng ge guójiā de qíngkuàng bù yíyàng.) « Les situations des deux pays sont différentes. »
\item \zh{在中国……，在喀麦隆……} (Zài Zhōngguó..., zài Kāmàilóng...) « En Chine..., au Cameroun... » — structure parallèle idéale pour comparer sans juger.
\end{itemize}

\begin{savaistu}[Le mobile money : une invention partagée]
Au Cameroun, le mobile money (Orange Money, MTN Mobile Money) permet de payer les frais de scolarité, les achats au marché de Mokolo ou les factures d'électricité, sans avoir de compte bancaire : il suffit d'un téléphone, même simple. Ce modèle a transformé la vie quotidienne de millions de Camerounais. La Chine connaît une révolution comparable avec Alipay (支付宝) et WeChat Pay (微信支付) : dans les villes chinoises, on paie le taxi, le restaurant et même le vendeur de rue en scannant un code QR. Deux pays, deux technologies, une même idée : le téléphone devient le porte-monnaie. Ni l'un ni l'autre n'a « copié » : chacun a répondu à ses propres besoins.
\end{savaistu}

\courssub{Tableau récapitulatif des structures}

\begin{center}
\begin{tabularx}{\linewidth}{|X|X|X|}
\hline
\rowcolor{popblueL} \textbf{Structure} & \textbf{Exemple (caractères + pinyin)} & \textbf{Traduction} \\
\hline
A 比 B + adj. & 微信比WhatsApp功能多。Wēixìn bǐ WhatsApp gōngnéng duō. & WeChat a plus de fonctions que WhatsApp. \\
\hline
A 比 B + adj. + 得多 & 支付宝比现金方便得多。Zhīfùbǎo bǐ xiànjīn fāngbiàn de duō. & Alipay est beaucoup plus pratique que l'argent liquide. \\
\hline
A 没有 B + adj. & 我的手机没有你的新。Wǒ de shǒujī méiyǒu nǐ de xīn. & Mon téléphone n'est pas aussi récent que le tien. \\
\hline
A 不如 B + adj. & 我的汉语不如你的好。Wǒ de Hànyǔ bùrú nǐ de hǎo. & Mon chinois est moins bon que le tien. \\
\hline
我觉得 / 我认为 + phrase & 我觉得手机支付很方便。Wǒ juéde shǒujī zhīfù hěn fāngbiàn. & Je trouve le paiement mobile très pratique. \\
\hline
因为……所以…… & 因为不用带现金，所以很方便。Yīnwèi bú yòng dài xiànjīn, suǒyǐ hěn fāngbiàn. & Parce qu'on n'a pas besoin de liquide, c'est donc pratique. \\
\hline
\end{tabularx}
\end{center}

\begin{center}
\begin{tabularx}{\linewidth}{|X|X|X|X|}
\hline
\rowcolor{popblueL} Caractères & Pinyin & Nature & Traduction \\
\hline
比较 & bǐjiào & verbe / adverbe & comparer ; assez, relativement \\
\hline
功能 & gōngnéng & nom & fonction, fonctionnalité \\
\hline
方便 & fāngbiàn & adjectif & pratique, commode \\
\hline
便宜 & piányi & adjectif & bon marché, pas cher \\
\hline
现金 & xiànjīn & nom & argent liquide, espèces \\
\hline
手机支付 & shǒujī zhīfù & nom & paiement mobile \\
\hline
支付宝 & Zhīfùbǎo & nom propre & Alipay \\
\hline
微信 & Wēixìn & nom propre & WeChat \\
\hline
好处 & hǎochu & nom & avantage, bénéfice \\
\hline
情况 & qíngkuàng & nom & situation, circonstances \\
\hline
越来越 & yuè lái yuè & adverbe & de plus en plus \\
\hline
常用 & chángyòng & adjectif & couramment utilisé \\
\hline
好用 & hǎoyòng & adjectif & pratique, facile à utiliser \\
\hline
合作 & hézuò & verbe / nom & coopérer / coopération \\
\hline
\end{tabularx}
\end{center}

\courssub{Exercice résolu et entraînement}

\begin{exoresolu}[Traduire et comparer]
Énoncé. Traduisez en chinois : « Je pense que le paiement mobile est plus pratique que l'argent liquide, parce qu'on peut payer l'école et le marché avec le téléphone. »
\tcblower
Solution détaillée, étape par étape :
\begin{enumerate}
\item \textbf{Opinion} : « je pense que » → \zh{我觉得} (wǒ juéde).
\item \textbf{Comparaison} : « plus pratique que » → structure A \zh{比} B + adjectif : \zh{手机支付比现金方便} (shǒujī zhīfù bǐ xiànjīn fāngbiàn). Attention : pas de \zh{很} avant \zh{方便} !
\item \textbf{Justification} : « parce que » → \zh{因为} (yīnwèi) ; « payer l'école et le marché avec le téléphone » → \zh{可以用手机付学费，也可以在市场买东西} (kěyǐ yòng shǒujī fù xuéfèi, yě kěyǐ zài shìchǎng mǎi dōngxi).
\item \textbf{Assemblage} : \zh{我觉得手机支付比现金方便，因为可以用手机付学费，也可以在市场买东西。}\\
(Wǒ juéde shǒujī zhīfù bǐ xiànjīn fāngbiàn, yīnwèi kěyǐ yòng shǒujī fù xuéfèi, yě kěyǐ zài shìchǎng mǎi dōngxi.)
\end{enumerate}
Vérification : la structure \zh{比} est bien en position A 比 B + adj., sans \zh{很} ; la cause est introduite par \zh{因为}. La phrase est correcte et naturelle.
\end{exoresolu}

\begin{exercice}[À vous de comparer et de donner votre avis]
\begin{enumerate}
\item Traduisez en chinois : « Mon téléphone n'est pas aussi grand que ton téléphone. » (utilisez \zh{没有})
\item Traduisez : « Je pense qu'apprendre le chinois est utile, parce que la Chine et le Cameroun travaillent beaucoup ensemble. » (utilisez \zh{我认为} et \zh{因为……所以……} si possible)
\item Corrigez la phrase fautive : \zh{支付宝很方便比现金。}
\item Complétez avec \zh{比}, \zh{没有} ou \zh{不如} : 我们城市的网络\_\_\_\_北京快。(Le réseau de ma ville est moins rapide que celui de Pékin.)
\item Production libre : rédigez deux phrases avec \zh{我觉得……因为……} pour comparer une application que vous utilisez avec une application chinoise.
\end{enumerate}
\end{exercice}

\begin{corrige}[Exercice « À vous de comparer »]
\begin{enumerate}
\item \zh{我的手机没有你的手机大。} (Wǒ de shǒujī méiyǒu nǐ de shǒujī dà.) — structure A 没有 B + adj., sans \zh{很}.
\item \zh{我认为学习汉语很有用，因为中国和喀麦隆经常合作。} (Wǒ rènwéi xuéxí Hànyǔ hěn yǒuyòng, yīnwèi Zhōngguó hé Kāmàilóng jīngcháng hézuò.) — ici \zh{很} est correct car il n'y a pas de \zh{比}.
\item \zh{支付宝比现金方便。} (Zhīfùbǎo bǐ xiànjīn fāngbiàn.) — on supprime \zh{很} et on place l'adjectif après B.
\item \zh{没有} (ou \zh{不如}) : \zh{我们城市的网络没有北京快。} (Wǒmen chéngshì de wǎngluò méiyǒu Běijīng kuài.)
\item Réponses libres ; modèle : \zh{我觉得WhatsApp很方便，因为我的家人和朋友都用它。} (Wǒ juéde WhatsApp hěn fāngbiàn, yīnwèi wǒ de jiārén hé péngyou dōu yòng tā.)
\end{enumerate}
\end{corrige}

\begin{aretenir}[L'essentiel de la section]
\begin{itemize}
\item Comparaison : \cle{A 比 B + adjectif} ; jamais de \zh{很} devant l'adjectif ; renforcement possible avec \zh{得多} ou \zh{多了}.
\item Négation : \cle{A 没有 B + adj.} (neutre) et \cle{A 不如 B + adj.} (moins ... que).
\item Opinion : \cle{我觉得} (oral, spontané) et \cle{我认为} (formel, exposé).
\item Justification : \cle{因为……所以……}, les deux connecteurs peuvent coexister en chinois, contrairement au français.
\item Comparer, c'est décrire des différences d'usages, pas établir une hiérarchie : \zh{各有各的好处} — « chacun a ses avantages ».
\end{itemize}
\end{aretenir}

Vous disposez maintenant de tous les outils pour la dernière étape de la leçon : le mini-exposé, où vous présenterez à la classe votre comparaison personnelle entre un outil TIC camerounais et un outil chinois, en utilisant au moins une comparaison avec \zh{比}, une opinion avec \zh{我觉得} et une justification avec \zh{因为……所以……}.

\coursec{Mini-exposé et réflexion}

Dans la section précédente, nous avons appris à comparer les pratiques numériques du Cameroun et de la Chine avec la structure \zh{比} et à donner un avis simple avec \zh{我觉得}. Nous allons maintenant rassembler tous ces outils dans une activité de synthèse : le \cle{mini-exposé} oral. C'est l'aboutissement de la leçon : vous devrez présenter, comparer et donner votre point de vue sur les outils TIC chinois et camerounais, en chinois, devant vos camarades.

\courssub{Questions de réflexion guidée}

Avant de préparer l'exposé, prenons quelques minutes pour réfléchir ensemble. Voici les questions directrices de cette leçon. Lisez-les à voix haute, puis discutez-en en français d'abord, en chinois ensuite.

\begin{enumerate}
\item \zh{信息通信技术让喀麦隆和中国更近吗？} \emph{Xìnxī tōngxìn jìshù ràng Kāmàilóng hé Zhōngguó gèng jìn ma?} — Les TIC rapprochent-elles le Cameroun et la Chine ?
\item \zh{手机有什么好处？有什么坏处？} \emph{Shǒujī yǒu shénme hǎochù? Yǒu shénme huàichù?} — Quels sont les avantages et les inconvénients du téléphone portable ?
\item \zh{你最喜欢哪个品牌？为什么？} \emph{Nǐ zuì xǐhuan nǎge pǐnpái? Wèishénme?} — Quelle marque préfères-tu ? Pourquoi ?
\item \zh{没有网络，你的生活会怎么样？} \emph{Méiyǒu wǎngluò, nǐ de shēnghuó huì zěnmeyàng?} — Sans Internet, à quoi ressemblerait ta vie ?
\end{enumerate}

\begin{exemplebox}[Deux réponses modèles]
\zh{我觉得信息通信技术让两个国家更近，因为我们可以用微信和中国人聊天。}\\
\emph{Wǒ juéde xìnxī tōngxìn jìshù ràng liǎng ge guójiā gèng jìn, yīnwèi wǒmen kěyǐ yòng Wēixìn hé Zhōngguó rén liáotiān.}\\
« Je pense que les TIC rapprochent les deux pays, parce que nous pouvons discuter avec des Chinois grâce à WeChat. »

\zh{手机很方便，但是也很浪费时间。}\\
\emph{Shǒujī hěn fāngbiàn, dànshì yě hěn làngfèi shíjiān.}\\
« Le téléphone est pratique, mais il fait aussi perdre du temps. »
\end{exemplebox}

\begin{attention}[Avantages et limites : le vocabulaire de l'équilibre]
Pour parler des avantages et des limites, les francophones oublient souvent les mots de contraste. Retenez la paire \cle{但是} (dànshì, mais) et \cle{可是} (kěshì, mais, plus oral). Structure : \textbf{phrase positive + 但是/可是 + phrase négative}. Ne dites jamais \zh{手机方便但是贵} sans sujet répété si la phrase devient longue : préférez \zh{手机很方便，但是太贵了。} (Shǒujī hěn fāngbiàn, dànshì tài guì le.) — « Le téléphone est pratique, mais il est trop cher. »
\end{attention}

\courssub{Le plan du mini-exposé en trois parties}

Un bon exposé, même court, suit toujours un plan clair. Pour cette leçon, nous utilisons le plan en trois parties : \textbf{présentation} (介绍 jièshào), \textbf{comparaison} (比较 bǐjiào), \textbf{avis personnel} (我的看法 wǒ de kànfǎ). La figure suivante résume le parcours complet, de l'introduction à la conclusion.

\begin{popfigure}
\begin{center}
\begin{tikzpicture}[
  node distance=0.55cm,
  box/.style={draw=popblue, fill=popblueL, rounded corners=3pt, align=center, minimum width=4.6cm, minimum height=1cm, font=\small},
  arr/.style={-{Stealth[length=2.5mm]}, thick, poporange}
]
\node[box, align=center] (intro){Introduction\\ \zh{我要介绍……}};
\node[box, below=of intro, align=center] (comp){Comparaison Cameroun / Chine\\ \zh{第一……第二……} + \zh{比}};
\node[box, below=of comp, align=center] (avis){Mon avis\\ \zh{我的看法是……} / \zh{我觉得……}};
\node[box, below=of avis, align=center] (concl){Conclusion\\ \zh{谢谢大家！}};
\draw[arr] (intro) -- (comp);
\draw[arr] (comp) -- (avis);
\draw[arr] (avis) -- (concl);
\end{tikzpicture}
\end{center}
\legende{Organigramme du plan d'exposé : quatre étapes, trois parties de contenu.}
\end{popfigure}

\begin{methode}[Réussir son mini-exposé en cinq étapes]
\begin{enumerate}
\item \textbf{Préparer une fiche} : écrivez 6 à 8 phrases courtes, pas un long texte. Une phrase = une idée.
\item \textbf{Utiliser le vocabulaire de la leçon} : \zh{手机} (shǒujī, téléphone), \zh{品牌} (pǐnpái, marque), \zh{网络} (wǎngluò, Internet), \zh{方便} (fāngbiàn, pratique), \zh{支付} (zhīfù, paiement).
\item \textbf{Faire au moins une comparaison avec} \zh{比} : \zh{中国的手机比法国的手机便宜。} (Zhōngguó de shǒujī bǐ Fǎguó de shǒujī piányi.) — « Les téléphones chinois sont moins chers que les téléphones français. »
\item \textbf{Donner un avis avec} \zh{我觉得} ou \zh{我的看法是} : c'est ce qui rend votre exposé personnel.
\item \textbf{Parler lentement, en regardant le public}, et conclure par \zh{谢谢大家！} (Xièxie dàjiā!) — « Merci à tous ! »
\end{enumerate}
\end{methode}

\courssub{Banque de phrases pour l'exposé}

Voici les structures-outils à mémoriser. Elles fonctionnent comme des « rails » : il suffit de remplacer la fin de la phrase par vos propres idées.

\begin{center}
\begin{tabularx}{\linewidth}{|X|X|X|}
\hline
\rowcolor{popblueL} Fonction & Structure chinoise + pinyin & Traduction \\
\hline
Annoncer le sujet & \zh{我要介绍……} Wǒ yào jièshào… & Je vais présenter… \\
\hline
Premier argument & \zh{第一，……} Dì-yī, … & Premièrement, … \\
\hline
Deuxième argument & \zh{第二，……} Dì-èr, … & Deuxièmement, … \\
\hline
Comparer & A \zh{比} B + adjectif & A est plus … que B \\
\hline
Donner son avis & \zh{我的看法是……} Wǒ de kànfǎ shì… & Mon avis est que… \\
\hline
Donner son avis (oral) & \zh{我觉得……} Wǒ juéde… & Je trouve que… \\
\hline
Nuancer & ……\zh{，但是……} …, dànshì… & …, mais… \\
\hline
Conclure & \zh{谢谢大家！} Xièxie dàjiā! & Merci à tous ! \\
\hline
\end{tabularx}
\end{center}

\begin{exemplebox}[Arguments enchaînés]
\zh{第一，中国品牌很便宜；第二，手机支付很方便。}\\
\emph{Dì-yī, Zhōngguó pǐnpái hěn piányi; dì-èr, shǒujī zhīfù hěn fāngbiàn.}\\
« Premièrement, les marques chinoises sont bon marché ; deuxièmement, le paiement mobile est pratique. »

Remarquez la ponctuation chinoise : le point-virgule pleine chasse \zh{；} relie les deux arguments, comme en français.
\end{exemplebox}

\begin{propriete}[Exprimer son point de vue : 我的看法是……]
\textbf{Forme} : \zh{我的看法是} + proposition complète.\\
\textbf{Emploi} : \cle{看法} (kànfǎ, point de vue, opinion) est un \textbf{nom}. En chinois, on ne peut pas dire « j'ai un avis » avec \zh{有} dans ce contexte ; on dit littéralement « mon avis \emph{est}… », d'où l'emploi obligatoire de \zh{是}.\\
\textbf{Équivalent français} : « mon avis est que… », « selon moi… ».\\
\textbf{Variante plus légère} : \zh{我觉得} + proposition (« je trouve que… »), très fréquente à l'oral.
\end{propriete}

\begin{attention}[Piège : ne dites pas 我有看法]
Les francophones traduisent mot à mot « j'ai un avis » par \zh{我有看法} : c'est une erreur. \zh{看法} s'emploie avec \zh{是} : \zh{我的看法是……}. On peut dire \zh{我有一个问题} (Wǒ yǒu yí ge wèntí, « j'ai une question »), car \zh{问题} se « possède », mais l'opinion s'« énonce » avec \zh{是} ou se « sent » avec \zh{觉得}.
\end{attention}

\courssub{Texte support : un exposé modèle}

Lisez ce court exposé rédigé par une élève de Terminale de Yaoundé. Il suit exactement le plan en trois parties et peut vous servir de modèle.

\begin{textebox}[Mon exposé : les TIC au Cameroun et en Chine]
\zh{大家好！今天我要介绍信息通信技术。}\\
\emph{Dàjiā hǎo! Jīntiān wǒ yào jièshào xìnxī tōngxìn jìshù.}\\
« Bonjour à tous ! Aujourd'hui, je vais présenter les technologies de l'information et de la communication. »\\[4pt]
\zh{第一，在喀麦隆，很多人用中国品牌的手机，因为这些手机不太贵。第二，在中国，人们常常用手机支付，这很方便。}\\
\emph{Dì-yī, zài Kāmàilóng, hěn duō rén yòng Zhōngguó pǐnpái de shǒujī, yīnwèi zhèxiē shǒujī bú tài guì. Dì-èr, zài Zhōngguó, rénmen chángcháng yòng shǒujī zhīfù, zhè hěn fāngbiàn.}\\
« Premièrement, au Cameroun, beaucoup de gens utilisent des téléphones de marques chinoises, parce qu'ils ne sont pas trop chers. Deuxièmement, en Chine, les gens paient souvent avec leur téléphone, c'est très pratique. »\\[4pt]
\zh{我觉得中国的网络比喀麦隆的快，但是喀麦隆的年轻人也很喜欢新技术。}\\
\emph{Wǒ juéde Zhōngguó de wǎngluò bǐ Kāmàilóng de kuài, dànshì Kāmàilóng de niánqīng rén yě hěn xǐhuan xīn jìshù.}\\
« Je trouve qu'Internet en Chine est plus rapide qu'au Cameroun, mais les jeunes Camerounais aiment aussi beaucoup les nouvelles technologies. »\\[4pt]
\zh{我的看法是，手机让喀麦隆人和中国人更容易交流。不过，手机也很浪费时间，我们要小心使用。}\\
\emph{Wǒ de kànfǎ shì, shǒujī ràng Kāmàilóng rén hé Zhōngguó rén gèng róngyì jiāoliú. Búguò, shǒujī yě hěn làngfèi shíjiān, wǒmen yào xiǎoxīn shǐyòng.}\\
« Mon avis est que le téléphone permet aux Camerounais et aux Chinois de communiquer plus facilement. Cependant, le téléphone fait aussi perdre du temps, nous devons l'utiliser avec prudence. »\\[4pt]
\zh{谢谢大家！}\\
\emph{Xièxie dàjiā!}\\
« Merci à tous ! »
\end{textebox}

\textbf{Glossaire du texte}

\begin{center}
\begin{tabularx}{\linewidth}{|X|X|X|X|}
\hline
\rowcolor{popblueL} Caractères & Pinyin & Nature & Traduction \\
\hline
\zh{介绍} & jièshào & verbe & présenter, introduire \\
\hline
\zh{品牌} & pǐnpái & nom & marque (commerciale) \\
\hline
\zh{支付} & zhīfù & verbe / nom & payer ; paiement \\
\hline
\zh{交流} & jiāoliú & verbe / nom & échanger ; échange \\
\hline
\zh{浪费} & làngfèi & verbe & gaspiller, perdre \\
\hline
\zh{小心} & xiǎoxīn & adjectif / adverbe & prudent ; avec prudence \\
\hline
\zh{不过} & búguò & conjonction & cependant, néanmoins \\
\hline
\zh{年轻人} & niánqīng rén & nom & les jeunes \\
\hline
\end{tabularx}
\end{center}

\textbf{Questions de compréhension}
\begin{enumerate}
\item Pourquoi les Camerounais utilisent-ils beaucoup de téléphones chinois, selon le texte ?
\item Quelle comparaison l'élève fait-elle entre la Chine et le Cameroun ?
\item Quel est l'inconvénient du téléphone mentionné à la fin ?
\end{enumerate}

\begin{corrige}[Questions de compréhension]
\begin{enumerate}
\item Parce que ces téléphones ne sont pas trop chers (\zh{不太贵} bú tài guì).
\item Elle compare la vitesse d'Internet : le réseau chinois est plus rapide que le réseau camerounais (\zh{比……快}).
\item Le téléphone fait perdre du temps (\zh{浪费时间} làngfèi shíjiān).
\end{enumerate}
\end{corrige}

\courssub{Passage à l'oral : du binôme à la classe}

\begin{experience}[Mini-exposé en trois temps]
\begin{enumerate}
\item \textbf{Préparation individuelle (10 minutes)} : chaque élève rédige sa fiche de 6 à 8 phrases en suivant l'organigramme : introduction, comparaison avec \zh{比}, avis avec \zh{我的看法是} ou \zh{我觉得}, conclusion.
\item \textbf{Entraînement par binôme (10 minutes)} : chacun présente son exposé à son voisin, qui écoute sans interrompre puis pose une question en chinois, par exemple \zh{你为什么喜欢这个手机？} (Nǐ wèishénme xǐhuan zhège shǒujī? — Pourquoi aimes-tu ce téléphone ?).
\item \textbf{Passage devant la classe} : trois ou quatre volontaires présentent leur exposé. Les autres élèves remplissent la grille d'évaluation par les pairs ci-dessous.
\end{enumerate}
\end{experience}

\begin{exoresolu}[Construire la phrase d'ouverture d'un exposé]
Énoncé : Paul veut commencer son exposé par : « Aujourd'hui, je vais présenter les marques de téléphones chinoises. » Aidez-le à construire la phrase en chinois.
\tcblower
Solution détaillée :
\begin{enumerate}
\item « Aujourd'hui » = \zh{今天} (jīntiān). En chinois, le complément de temps se place après le sujet, avant le verbe.
\item « Je vais présenter » = \zh{我要介绍} (wǒ yào jièshào) : \zh{要} exprime ici l'intention proche du futur.
\item « Les marques de téléphones chinoises » = \zh{中国手机的牌子} ou plus simplement \zh{中国的手机品牌} (Zhōngguó de shǒujī pǐnpái) : le complément du nom précède le nom avec \zh{的}.
\item Phrase complète : \zh{今天我要介绍中国的手机品牌。} \emph{Jīntiān wǒ yào jièshào Zhōngguó de shǒujī pǐnpái.}
\end{enumerate}
Erreur à éviter : ne placez pas \zh{今天} à la fin de la phrase comme en français (« je vais présenter… aujourd'hui ») : en chinois, le temps se dit avant le verbe.
\end{exoresolu}

\courssub{Évaluation par les pairs}

Pendant les exposés, chaque élève-évaluateur coche la grille suivante (0 = absent, 1 = partiel, 2 = réussi). L'objectif n'est pas de sanctionner, mais d'aider chacun à progresser.

\begin{center}
\begin{tabularx}{\linewidth}{|X|X|X|X|}
\hline
\rowcolor{popblueL} Critère & 0 & 1 & 2 \\
\hline
Vocabulaire de la leçon (\zh{手机}, \zh{品牌}, \zh{网络}, \zh{方便}…) & non utilisé & 2-3 mots & 4 mots ou plus \\
\hline
Une comparaison avec \zh{比} & absente & présente mais fautive & correcte \\
\hline
Un avis avec \zh{我的看法是} ou \zh{我觉得} & absent & présent mais fautif & correct \\
\hline
Clarté et prononciation (tons) & difficile à suivre & compréhensible & claire \\
\hline
Plan respecté (intro → comparaison → avis → conclusion) & non & partiellement & oui \\
\hline
\end{tabularx}
\end{center}

\begin{exercice}[À vous de parler]
Préparez votre propre mini-exposé de 6 à 8 phrases sur le thème « Les outils TIC au Cameroun et en Chine ». Votre exposé doit contenir obligatoirement : une phrase d'ouverture avec \zh{我要介绍}, deux arguments avec \zh{第一} et \zh{第二}, une comparaison avec \zh{比}, un avis avec \zh{我的看法是} ou \zh{我觉得}, une limite avec \zh{但是}, et la formule de clôture \zh{谢谢大家}.
\end{exercice}

\begin{corrige}[Exercice « À vous de parler » — proposition de réponse]
\zh{大家好！今天我要介绍喀麦隆和中国的手机。第一，很多喀麦隆人用中国品牌的手机，因为它们很便宜。第二，中国人常常用手机支付，这很方便。我觉得中国的网络比喀麦隆的快。我的看法是，手机让两个国家的人更容易交流，但是我们不应该浪费太多时间。谢谢大家！}\\
\emph{Dàjiā hǎo! Jīntiān wǒ yào jièshào Kāmàilóng hé Zhōngguó de shǒujī. Dì-yī, hěn duō Kāmàilóng rén yòng Zhōngguó pǐnpái de shǒujī, yīnwèi tāmen hěn piányi. Dì-èr, Zhōngguó rén chángcháng yòng shǒujī zhīfù, zhè hěn fāngbiàn. Wǒ juéde Zhōngguó de wǎngluò bǐ Kāmàilóng de kuài. Wǒ de kànfǎ shì, shǒujī ràng liǎng ge guójiā de rén gèng róngyì jiāoliú, dànshì wǒmen bù yīnggāi làngfèi tài duō shíjiān. Xièxie dàjiā!}\\
« Bonjour à tous ! Aujourd'hui, je vais présenter les téléphones au Cameroun et en Chine. Premièrement, beaucoup de Camerounais utilisent des téléphones de marques chinoises parce qu'ils sont bon marché. Deuxièmement, les Chinois paient souvent avec leur téléphone, c'est très pratique. Je trouve qu'Internet en Chine est plus rapide qu'au Cameroun. Mon avis est que le téléphone permet aux habitants des deux pays de communiquer plus facilement, mais nous ne devrions pas perdre trop de temps. Merci à tous ! »
\end{corrige}

\begin{savaistu}[Pourquoi 看法 prend-il 是 ?]
En chinois, de nombreux noms abstraits liés à la pensée fonctionnent comme \zh{看法} : \zh{意思} (yìsi, sens), \zh{意见} (yìjiàn, opinion, suggestion). On dit \zh{我的意见是……} (wǒ de yìjiàn shì…, « mon opinion est que… ») sur le même modèle. Le verbe \zh{是} joue le rôle de panneau indicateur : il annonce que ce qui suit est le contenu de l'opinion. C'est une différence majeure avec le français, où « avoir » domine (avoir un avis, avoir une idée) : en chinois, on \emph{déclare} son avis plutôt qu'on ne le \emph{possède}.
\end{savaistu}

\begin{aretenir}[L'essentiel de la section]
\begin{itemize}
\item Un mini-exposé suit toujours le plan : introduction (\zh{我要介绍}) → comparaison (\zh{第一……第二……} + \zh{比}) → avis (\zh{我的看法是} / \zh{我觉得}) → conclusion (\zh{谢谢大家}).
\item \zh{看法} (point de vue) s'emploie avec \zh{是} : jamais \zh{我有看法}.
\item Pour nuancer : \zh{但是} / \zh{可是} / \zh{不过} + idée opposée (\zh{很方便，但是太贵了}).
\item Le complément de temps (\zh{今天}) se place avant le verbe, jamais en fin de phrase.
\item À l'oral : phrases courtes, tons soignés, regard vers le public.
\end{itemize}
\end{aretenir}

\newpage

\coursec{Activité d'intégration}

\coursec{Leçon 37 — Éléments socioculturels : les outils TICs de marques chinoises et camerounaises}

\courssub{Objectifs de la leçon}

À l'issue de cette leçon, l'élève sera capable de :

\begin{itemize}
\item mobiliser le \cle{vocabulaire des TIC} (\zh{手机、网络、品牌、软件、支付、上网、转账}…) dans une production écrite et orale ;
\item comparer des pratiques camerounaises et chinoises à l'aide des structures \cle{比} (supériorité) et \cle{没有} (égalité / infériorité) ;
\item exprimer un point de vue personnel simple et justifié avec \cle{我觉得……因为……} ;
\item rédiger un court exposé structuré (marques, usages, comparaison, avis) et le présenter en 3 minutes ;
\item identifier des faits socioculturels factuels : présence des marques chinoises (Transsion : Tecno, Itel, Infinix ; Huawei, Xiaomi, OPPO) au Cameroun et similarités / différences d'usages (mobile money, vidéo, football).
\end{itemize}

\courssub{Découverte : les TIC au quotidien}

\begin{experience}[Observation guidée — 手机在我们身边]
En classe entière, le professeur projette ou affiche trois photos : un jeune à Yaoundé payant son taxi-moto via \zh{手机支付} (Mobile Money / MoMo), un étudiant à Pékin scannant un \zh{二维码} (code QR) pour louer un vélo partagé, un élève à Buea révisant son chinois sur une application (\zh{软件}) sur son \zh{手机}.
\begin{enumerate}
\item Nommer en français les trois actions observées.
\item Deviner les mots chinois correspondants avec l'aide du tableau de vocabulaire ci-dessous.
\item Question déclencheur : « Selon vous, les marques de téléphones tenues par ces jeunes sont-elles les mêmes ? Les usages sont-ils identiques ? »
\end{enumerate}
\end{experience}

\courssub{Vocabulaire clé du thème — 重点词汇}

\begin{center}
\begin{tabularx}{\linewidth}{|X|X|X|X|}
\hline
\rowcolor{popblueL} Caractères & Pinyin & Nature & Traduction \\
\hline
手机 & shǒujī & n. & téléphone portable, smartphone \\
\hline
品牌 & pǐnpái & n. & marque (commerciale) \\
\hline
网络 & wǎngluò & n. & réseau, Internet \\
\hline
上网 & shàngwǎng & v. & aller sur Internet, se connecter \\
\hline
软件 & ruǎnjiàn & n. & logiciel, application (appli) \\
\hline
支付 & zhīfù & v./n. & paiement, payer (souvent \zh{手机支付} / \zh{移动支付}) \\
\hline
转账 & zhuǎnzhàng & v. & virer / transférer de l'argent (mobile money) \\
\hline
方便 & fāngbiàn & adj. & pratique, commode \\
\hline
便宜 & piányi & adj. & bon marché, pas cher \\
\hline
看法 & kànfǎ & n. & point de vue, opinion \\
\hline
功能 & gōngnéng & n. & fonction, fonctionnalité \\
\hline
视频 & shìpín & n. & vidéo \\
\hline
音乐 & yīnyuè & n. & musique \\
\hline
足球比赛 & zúqiú bǐsài & n. & match de football \\
\hline
青年 & qīngnián & n. & jeune, jeunesse \\
\hline
论坛 & lùntán & n. & forum \\
\hline
\end{tabularx}
\end{center}

\courssub{Clés et écriture des caractères — 汉字拆解与书写}

\begin{definition}[Analyse de caractères clés]
\textbf{品  pǐn  (marque, produit)} : trois \zh{口} (bouche/entrée) empilés → « multitude de bouches » → réputation, gamme de produits. Ordre : haut → bas, chaque \zh{口} s'écrit : verticale, horizontale tournante, horizontale.\\
\textbf{牌  pái  (panneau, marque, carte)} : clé \zh{片} (pièce de bois/plaque) + phonétique \zh{牌} (simplifié de \zh{爿} + \zh{秀}). Idée : panneau en bois indiquant la marque. Ordre : clé \zh{片} (gauche → droite : piqué, horizontale, piqué, horizontale) puis partie droite.\\
\textbf{网  wǎng  (filet, réseau, Internet)} : clé \zh{纟} (soie, forme simplifiée de \zh{糸}) + \zh{亡}. Ordre : clé \zh{纟} (point, pli, horizontale, point) puis \zh{亡} (horizontale, horizontale, piqué, point, horizontale).\\
\textbf{支  zhī  (payer, soutenir ; classif. objets longs)} : \zh{扌} (main) + \zh{枝} (branche). Ordre : \zh{扌} (horizontale, crochet vertical, point) puis \zh{枝} (bois + phonétique).\\
\textbf{转  zhuǎn  (tourner, transférer)} : clé \zh{车} (véhicule/roue) + phonétique \zh{专}. Idée : la roue tourne → transférer. Ordre : \zh{车} (horizontale, verticale centrale, horizontale, horizontale) puis \zh{专} (point, horizontale, verticale, horizontales).\\
\textbf{软  ruǎn  (doux, mou ; software = logiciel)} : clé \zh{车} + phonétique \zh{蕊} (simplifié). Ordre : \zh{车} puis le haut (herbes \zh{艹} simplifié) et le bas \zh{车} (répété).\\
\textbf{件  jiàn  (objet, affaire ; classif. objets, vêtements, problèmes)} : \zh{亻} (homme) + \zh{牛} (bœuf). Ordre : \zh{亻} (piqué, verticale) puis \zh{牛} (horizontale, verticale, horizontale, verticale crochet, horizontale).
\end{definition}

\courssub{Textes supports et compréhension — 课文与理解}

\begin{textebox}[Annonce du forum — 论坛通知]
青年论坛：手机和我们的生活\\
Qīngnián lùntán : shǒujī hé wǒmen de shēnghuó\\
« Forum des jeunes : le téléphone portable et notre vie »\\[4pt]
请每个小组介绍两个中国手机品牌和两个喀麦隆人常用的手机功能，\\
Qǐng měi ge xiǎozǔ jièshào liǎng ge Zhōngguó shǒujī pǐnpái hé liǎng ge Kāmàilóng rén chángyòng de shǒujī gōngnéng,\\
« Chaque groupe présente deux marques chinoises de téléphones et deux usages courants du téléphone au Cameroun, »\\
再比较中国和喀麦隆的情况，最后说说你们的看法。\\
Zài bǐjiào Zhōngguó hé Kāmàilóng de qíngkuàng, zuìhòu shuōshuo nǐmen de kànfǎ.\\
« puis compare la situation de la Chine et du Cameroun, et donne enfin votre avis. »\\[4pt]
时间：三分钟。语言：汉语。\\
Shíjiān : sān fēnzhōng. Yǔyán : Hànyǔ.\\
« Durée : trois minutes. Langue : le chinois. »
\end{textebox}

\begin{textebox}[Fiche d'information — 中国朋友的资料]
在中国，很多人用华为、小米和OPPO的手机。\\
Zài Zhōngguó, hěn duō rén yòng Huáwéi, Xiǎomǐ hé OPPO de shǒujī.\\
« En Chine, beaucoup de gens utilisent des téléphones Huawei, Xiaomi et OPPO. »\\
中国人用手机付钱、买东西、看视频，很方便。\\
Zhōngguó rén yòng shǒujī fùqián, mǎi dōngxi, kàn shìpín, hěn fāngbiàn.\\
« Les Chinois paient, achètent et regardent des vidéos avec leur téléphone : c'est très pratique. »\\
在喀麦隆，很多年轻人用Tecno和Itel的手机，因为不太贵。\\
Zài Kāmàilóng, hěn duō niánqīngrén yòng Tecno hé Itel de shǒujī, yīnwèi bú tài guì.\\
« Au Cameroun, beaucoup de jeunes utilisent des téléphones Tecno et Itel, parce qu'ils ne sont pas trop chers. »\\
他们也用手机转钱、听音乐、看足球比赛。\\
Tāmen yě yòng shǒujī zhuǎnqián, tīng yīnyuè, kàn zúqiú bǐsài.\\
« Ils utilisent aussi le téléphone pour transférer de l'argent, écouter de la musique et regarder des matchs de football. »
\end{textebox}

\courssub{Exercices de compréhension et de lexique}

\begin{exercice}[Compréhension de la fiche d'information]
D'après le texte \emph{Fiche d'information}, répondre en chinois (caractères + pinyin) :
\begin{enumerate}
\item 列出文中提到的三个中国手机品牌。 (Citer les 3 marques chinoises)
\item 中国人常用手机做哪三件事？ (Quelles sont les 3 actions courantes des Chinois avec leur téléphone ?)
\item 喀麦隆年轻人常用什么品牌的手机？为什么？ (Quelle marque les jeunes Camerounais ? Pourquoi ?)
\item 喀麦隆人用手机做哪两件事？ (Deux usages au Cameroun.)
\end{enumerate}
\end{exercice}

\begin{corrige}[Exercice 1]
\begin{enumerate}
\item 华为 (Huáwéi)、小米 (Xiǎomǐ)、OPPO.
\item 付钱 (fùqián)、买东西 (mǎi dōngxi)、看视频 (kàn shìpín).
\item Tecno 和 Itel，因为不太贵 (piányi / bú tài guì).
\item 转钱 (zhuǎnqián)、听音乐 (tīng yīnyuè) / 看足球比赛 (kàn zúqiú bǐsài).
\end{enumerate}
\end{corrige}

\begin{exercice}[Complétion lexicale — 词汇填空]
Compléter les phrases avec un mot du tableau de vocabulaire (caractères + pinyin).
\begin{enumerate}
\item 我每天用手机 \_\_\_\_\_\_ 看新闻。 (aller sur Internet)
\item 这个 \_\_\_\_\_\_ 很有名，大家都喜欢买。 (marque)
\item 在喀麦隆，\_\_\_\_\_\_ 很流行，不用带现金。 (mobile money / paiement mobile)
\item 我的手机 \_\_\_\_\_\_ 很多，可以拍照、录像、上网。 (fonctionnalités)
\item 我觉得这个软件很 \_\_\_\_\_\_，操作简单。 (pratique)
\end{enumerate}
\end{exercice}

\begin{corrige}[Exercice 2]
\begin{enumerate}
\item 上网 (shàngwǎng)
\item 品牌 (pǐnpái)
\item 手机支付 (shǒujī zhīfù) / 转账 (zhuǎnzhàng)
\item 功能 (gōngnéng)
\item 方便 (fāngbiàn)
\end{enumerate}
\end{corrige}

\courssub{Grammaire : Comparer et donner son avis}

\begin{propriete}[La comparaison avec 比 (supériorité)]
\textbf{Structure} : \cle{A + 比 + B + Adjectif (SV)}.\\
\textbf{Emploi} : A possède la qualité exprimée par l'adjectif à un degré supérieur à B. \textbf{Pas de \zh{很} après 比}.\\
\textbf{Exemples} :\\
\zh{中国的网络比喀麦隆的快。} \emph{Zhōngguó de wǎngluò bǐ Kāmàilóng de kuài.} (L'Internet chinois est plus rapide que celui du Cameroun.)\\
\zh{华为的手机比Tecno的贵。} \emph{Huáwéi de shǒujī bǐ Tecno de guì.} (Les téléphones Huawei sont plus chers que les Tecno.)\\
\zh{这个软件比那个好用。} \emph{Zhège ruǎnjiàn bǐ nàge hǎoyòng.} (Cette appli est plus facile à utiliser que celle-là.)
\end{propriete}

\begin{propriete}[La comparaison avec 没有 (infériorité / inégalité)]
\textbf{Structure} : \cle{A + 没有 + B + Adjectif (SV)}.\\
\textbf{Sens} : « A n'est pas aussi [adjectif] que B » (A < B). L'adjectif garde sa forme positive.\\
\textbf{Exemples} :\\
\zh{喀麦隆的网络没有中国的快。} \emph{Kāmàilóng de wǎngluò méiyǒu Zhōngguó de kuài.} (L'Internet camerounais n'est pas aussi rapide que le chinois.)\\
\zh{我的手机没有他的新。} \emph{Wǒ de shǒujī méiyǒu tā de xīn.} (Mon téléphone n'est pas aussi neuf que le sien.)\\
\zh{Tecno没有华为贵。} \emph{Tecno méiyǒu Huáwéi guì.} (Tecno n'est pas aussi cher que Huawei = Tecno est moins cher.)
\end{propriete}

\begin{propriete}[Exprimer un point de vue : 我觉得……因为……]
\textbf{Structure} : \cle{Subject + 觉得 + Opinion/Évaluation , 因为 + Raison/Fait}.\\
\textbf{Nuance} : \zh{觉得} (juéde) = « trouver que, avoir l'impression que » (subjectif). Ne pas confondre avec \zh{认为} (rènwéi) « considérer que » (plus formel) ou \zh{以为} (yǐwéi) « croire à tort que ».\\
\textbf{Exemples} :\\
\zh{我觉得手机很有用，因为可以用手机学习。} \emph{Wǒ juéde shǒujī hěn yǒuyòng, yīnwèi kěyǐ yòng shǒujī xuéxí.}\\
\zh{我觉得Tecno性价比高，因为不太贵，功能也多。} \emph{Wǒ juéde Tecno xìngjiàbǐ gāo, yīnwèi bú tài guì, gōngnéng yě duō.} (Je trouve que Tecno a un bon rapport qualité-prix, car pas trop cher et beaucoup de fonctions.)
\end{propriete}

\courssub{Attention : Erreurs fréquentes des francophones}

\begin{attention}[Pièges à éviter]
\begin{itemize}
\item \textbf{比 + 很 + Adj.} $\rightarrow$ \textbf{Interdit}. \emph{Faux} : \zh{比很贵} / \zh{比很方便}. \emph{Correct} : \zh{比贵} / \zh{比方便}.
\item \textbf{Oubli du \zh{的} (particule possessive/attributive)}. \emph{Faux} : \zh{喀麦隆手机} (ambigu). \emph{Correct} : \zh{喀麦隆的手机} (les téléphones du Cameroun / camerounais), \zh{华为的} (ceux de Huawei).
\item \textbf{Ordre du complément de lieu \zh{在}}. \emph{Faux} : \zh{很多人用Tecno 在喀麦隆}. \emph{Correct} : \zh{在喀麦隆，很多人用Tecno} (Le lieu \zh{在...} précède le verbe).
\item \textbf{Confusion \zh{付钱} / \zh{转账} / \zh{花钱}}. \zh{付钱} = payer (un achat) ; \zh{转账} = virer/transférer (mobile money, P2P) ; \zh{花钱} = dépenser de l'argent.
\item \textbf{Classificateur \zh{个} vs \zh{款} / \zh{部}}. Pour téléphone : \zh{一部手机} (formel, appareil), \zh{一个手机} (oral, courant), \zh{一款手机} (modèle de téléphone). Au HSK 3/4, \zh{部} et \zh{款} sont attendus.
\end{itemize}
\end{attention}

\courssub{Entraînement écrit : Comparer et argumenter}

\begin{exercice}[Transformation : Comparer avec 比 et 没有]
Traduire en chinois (caractères + pinyin).
\begin{enumerate}
\item Les téléphones Xiaomi sont moins chers que les iPhones. (用 没有)
\item L'Internet au Cameroun n'est pas aussi rapide qu'en Chine. (用 没有)
\item Les applications chinoises sont plus pratiques que les applications camerounaises. (用 比)
\item Je trouve que les marques chinoises sont très bonnes. (用 我觉得)
\item J'aime bien Tecno, car le rapport qualité-prix est élevé. (用 觉得……因为…… ; \zh{性价比高} = bon rapport qualité-prix)
\end{enumerate}
\end{exercice}

\begin{corrige}[Exercice 3]
\begin{enumerate}
\item \zh{小米的手机没有iPhone贵。} \emph{Xiǎomǐ de shǒujī méiyǒu iPhone guì.}
\item \zh{喀麦隆的网络没有中国的快。} \emph{Kāmàilóng de wǎngluò méiyǒu Zhōngguó de kuài.}
\item \zh{中国的软件比喀麦隆的方便。} \emph{Zhōngguó de ruǎnjiàn bǐ Kāmàilóng de fāngbiàn.}
\item \zh{我觉得中国品牌很好。} \emph{Wǒ juéde Zhōngguó pǐnpái hěn hǎo.}
\item \zh{我觉得Tecno不错，因为性价比高。} \emph{Wǒ juéde Tecno bùcuò, yīnwèi xìngjiàbǐ gāo.}
\end{enumerate}
\end{corrige}

\begin{exercice}[Traduction thème/version — 翻译练习]
\begin{enumerate}
\item \textbf{Version} : Traduire en français.
\zh{在喀麦隆，转账和看足球是年轻人最常用的手机功能。}
\item \textbf{Thème} : Traduire en chinois (caractères + pinyin).
« Au forum, notre groupe a présenté deux marques chinoises (Huawei, Xiaomi) et deux usages camerounais (mobile money, football). Nous avons comparé les prix et la vitesse d'Internet. Je pense que le téléphone change notre vie. »
\end{enumerate}
\end{exercice}

\begin{corrige}[Exercice 4]
\begin{enumerate}
\item Au Cameroun, le transfert d'argent et regarder le football sont les fonctions de téléphone les plus utilisées par les jeunes.
\item \zh{在论坛上，我们小组介绍了两个中国品牌（华为、小米）和两个喀麦隆常用的功能（转账、看足球）。我们比较了价格和网速。我觉得手机改变了我们的生活。}\\
\emph{Zài lùntán shàng, wǒmen xiǎozǔ jièshào le liǎng ge Zhōngguó pǐnpái (Huáwéi, Xiǎomǐ) hé liǎng ge Kāmàilóng chángyòng de gōngnéng (zhuǎnzhàng, kàn zúqiú). Wǒmen bǐjiào le jiàgé hé wǎngsù. Wǒ juéde shǒujī gǎibiàn le wǒmen de shēnghuó.}
\end{enumerate}
\end{corrige}

\courssub{Production orale : Mini-exposé et affiche bilingue}

\begin{methode}[Préparation de l'affiche-orale « 手机和我们的生活 » (3 min / groupe de 4)]
\begin{enumerate}
\item \textbf{Répartition des rôles} (5 min) : Membre A = Marques chinoises ; Membre B = Usages camerounais ; Membre C = Comparaison (比 / 没有) ; Membre D = Avis du groupe (我觉得…因为…).
\item \textbf{Rédaction du script} (15 min) : Chaque membre rédige ses 2--3 phrases (caractères + pinyin). Le groupe relit et corrige (vérifier : pas de \zh{很} après \zh{比}, \zh{的} présents, ordre \zh{在...}).
\item \textbf{Conception de l'affiche} (10 min) : Format A3. Titre \zh{手机和我们的生活}. 4 cases : 1. 中国品牌 (logos + noms) 2. 喀麦隆用法 (icônes : MoMo, ⚽, ) 3. 比较 (flèches \zh{比} / \zh{没有}) 4. 我们的看法 (bulle \zh{我觉得...}). Dessins simples acceptés.
\item \textbf{Répétition} (5 min) : Chronométrer. Chacun parle ~45 sec. Regarder le public, pas la feuille.
\item \textbf{Présentation \& Interaction} : Groupes passent. Public note une question par groupe (modèles : \zh{你们用什么手机？} / \zh{为什么Tecno比华为便宜？} / \zh{你们觉得手机方便吗？}).
\end{enumerate}
\end{methode}

\begin{experience}[Forum des jeunes — 青年论坛实况]
\textbf{Déroulement} :\\
1. Installation des affiches au mur / tableau.\\
2. Chaque groupe présente (3 min). Professeur note : vocabulaire (5 pts), structures \zh{比}/\zh{没有}/\zh{觉得} (5 pts), prononciation/fluidité (5 pts), originalité affiche (5 pts).\\
3. Après chaque passage : 2 questions du public en chinois. Le groupe répond brièvement.\\
4. Vote de la classe (bulletins secrets) pour la « Meilleure affiche » (critères affichés).\\
5. Affichage de l'affiche gagnante au mur de la salle pour le module « Mass médias ».
\end{experience}

\begin{exoresolu}[Affiche-orale modèle — 手机和我们的生活 (Script complet, ~10 phrases)]
Énoncé : Rédiger le script oral complet d'un groupe, couvrant les 4 rubriques : 2 marques chinoises, 2 usages camerounais, 1 comparaison (\zh{比} ou \zh{没有}), 1 avis justifié (\zh{我觉得...因为...}).
\tcblower
\textbf{Production modèle (script du groupe) :}\\[4pt]
\zh{大家好！今天我们介绍“手机和我们的生活”。}\\
\emph{Dàjiā hǎo! Jīntiān wǒmen jièshào “shǒujī hé wǒmen de shēnghuó”.}\\
« Bonjour à tous ! Aujourd'hui nous présentons “Le téléphone et notre vie”. »\\[3pt]
\zh{在中国，华为和小米是很有名的品牌。中国人常用手机付钱、买东西、看视频，非常方便。}\\
\emph{Zài Zhōngguó, Huáwéi hé Xiǎomǐ shì hěn yǒumíng de pǐnpái. Zhōngguó rén chángyòng shǒujī fùqián, mǎi dōngxi, kàn shìpín, fēicháng fāngbiàn.}\\
« En Chine, Huawei et Xiaomi sont des marques très connues. Les Chinois paient, achètent et regardent des vidéos sur mobile, c'est extrêmement pratique. »\\[3pt]
\zh{在喀麦隆，很多年轻人用Tecno和Itel的手机，因为性价比高。我们常用手机转账、听音乐和看足球比赛。}\\
\emph{Zài Kāmàilóng, hěn duō niánqīngrén yòng Tecno hé Itel de shǒujī, yīnwèi xìngjiàbǐ gāo. Wǒmen chángyòng shǒujī zhuǎnzhàng, tīng yīnyuè hé kàn zúqiú bǐsài.}\\
« Au Cameroun, beaucoup de jeunes utilisent Tecno et Itel car le rapport qualité-prix est élevé. Nous utilisons souvent le mobile money, la musique et les matchs de foot. »\\[3pt]
\zh{Tecno的手机比华为的便宜，但是中国的网络比喀麦隆的快。}\\
\emph{Tecno de shǒujī bǐ Huáwéi de piányi, dànshì Zhōngguó de wǎngluò bǐ Kāmàilóng de kuài.}\\
« Les téléphones Tecno sont moins chers que les Huawei, mais l'Internet chinois est plus rapide que celui du Cameroun. »\\[3pt]
\zh{我觉得智能手机改变了我们的生活，因为它让学习、交流和支付都变得更容易。谢谢大家！}\\
\emph{Wǒ juéde zhìnéng shǒujī gǎibiàn le wǒmen de shēnghuó, yīnwèi tā ràng xuéxí, jiāoliú hé zhīfù dōu biànde róngyì. Xièxie dàjiā!}\\
« Je trouve que le smartphone a changé notre vie, car il rend l'étude, la communication et le paiement plus faciles. Merci à tous ! »\\[6pt]
\textbf{Commentaire linguistique :}
\begin{itemize}
\item \textbf{Marques} : Noms propres inchangés ; \zh{的品牌 / 的手机} pour qualifier.
\item \textbf{Comparaison} : \zh{Tecno的手机比华为的便宜} — \zh{的} nominalisé (ellipse de \zh{手机}) ; pas de \zh{很}.
\item \textbf{Concession} : \zh{但是} équilibre le propos (prix vs infrastructure), niveau Terminale attendu.
\item \textbf{Opinion} : \zh{我觉得……因为……} + énumération verbale \zh{学习、交流、支付} (rythme binaire/ternaire avec \zh{、}).
\item \textbf{Lexique mobilisé} : 10+ mots de la leçon (品牌, 付钱, 转账, 网络, 方便, 便宜, 性价比, 功能, 视频, 足球比赛, 看法).
\end{itemize}
\end{exoresolu}

\courssub{Élément socioculturel approfondi : Le groupe Transsion}

\begin{savaistu}[Transsion Holdings : le géant chinois « invisible » en Chine, roi en Afrique]
\begin{itemize}
\item \textbf{Transsion Holdings (传音控股 Chuányīn Kònggǔ)} est une entreprise chinoise basée à Shenzhen (深圳).
\item Ses marques \textbf{Tecno (传音 TECNO)}, \textbf{Itel (爱特尔 Àitè'ěr)}, \textbf{Infinix (英菲尼克斯 Yīngfēiníkèsī)} dominent le marché africain (part de marché \#1 au Cameroun, Nigeria, Kenya, Éthiopie...).
\item \textbf{Particularité} : Ces marques sont \textbf{quasi absentes du marché chinois continental}. Elles sont conçues \emph{pour} l'Afrique : double SIM, batterie géante, appareil photo optimisé peaux foncées, système \zh{HiOS} (basé sur Android) avec apps locales préinstallées (Mobile Money, musique africaine, football).
\item \textbf{Comparaison factuelle} : Huawei / Xiaomi / OPPO / Vivo ciblent le marché chinois (haut de gamme, 5G, écosystème IoT) + export mondial. Transsion cible \emph{spécifiquement} les marchés émergents (Afrique, Asie du Sud, Amérique Latine).
\item \textbf{Leçon} : « Marque chinoise » ne signifie pas « vendue en Chine ». La stratégie d'adaptation locale (\zh{本地化策略}) est un cas d'école en économie internationale.
\end{itemize}
\end{savaistu}

\courssub{Bilan et évaluation}

\begin{aretenir}[Bilan de la leçon 37]
\begin{itemize}
\item \textbf{Vocabulaire} : Maîtriser 15 mots TIC (marques, usages, paiement, réseau, avis).
\item \textbf{Grammaire} : \cle{A 比 B Adj.} (supériorité) ; \cle{A 没有 B Adj.} (infériorité/égalité) ; \cle{我觉得……因为……} (opinion justifiée). Respecter l'ordre \zh{在 + Lieu} en tête de phrase.
\item \textbf{Socioculturel} : Connaître les marques Huawei, Xiaomi, OPPO (Chine) vs Tecno, Itel, Infinix (Transsion, Afrique). Usages communs (vidéo, musique, paiement) ; différences d'infrastructure (vitesse réseau, écosystème paiement).
\item \textbf{Compétences} : Lire une annonce + fiche info ; rédiger un script structuré ; présenter à l'oral 3 min ; poser/répondre à une question simple ; créer une affiche bilingue.
\end{itemize}
\end{aretenir}

\begin{exercice}[Évaluation formative — Devoirs / Travail personnel]
\begin{enumerate}
\item Rédiger un paragraphe de 80--100 caractères chinois : « Mon téléphone idéal » (marque, prix, fonctions, pourquoi). Utiliser \zh{比}, \zh{没有}, \zh{觉得...因为...}.
\item Préparer 3 questions en chinois pour interviewer un camarade sur ses usages du téléphone (mobile money, jeux, études, réseaux sociaux).
\item Recherche : Trouver le chiffre d'affaires 2023 de Transsion en Afrique (source fiable) et le comparer à celui de Huawei Consumer BG. Présenter en 2 phrases chinoises la semaine prochaine.
\end{enumerate}
\end{exercice}

\newpage

\coursec{Méthodes et exercices résolus}

\courssub{Méthodes et exercices résolus}

\begin{methode}[Comparer les pratiques numériques Cameroun / Chine : structure \cle{比} et \cle{跟……一样/不一样}]
\textbf{Objectif :} Produire des phrases comparatives correctes pour opposer ou rapprocher les usages des TIC (applications, paiement mobile, commerce en ligne) dans les deux pays.
\begin{enumerate}
    \item \textbf{Identifier les deux termes de la comparaison} (A et B) : ex. \zh{微信} (WeChat) vs WhatsApp / Orange Money.
    \item \textbf{Choisir la structure adaptée} :
        \begin{itemize}
            \item Différence marquée (supériorité) : \textbf{A + 比 + B + adjectif}. \emph{Pas de \zh{很} avant l'adjectif dans cette structure.}
            \item Égalité : \textbf{A + 跟/和 + B + 一样 + adjectif/nom}.
            \item Inégalité : \textbf{A + 跟/和 + B + 不一样} ou \textbf{A + 没有 + B + 这么/那么 + adjectif}.
        \end{itemize}
    \item \textbf{Placer le complément de degré} (beaucoup / un peu) \textbf{après} l'adjectif : \zh{比……方便得多}、\zh{比……快一点儿}。
    \item \textbf{Vérifier l'ordre des mots} : en chinois, le terme de référence (B) se place \textbf{avant} l'adjectif, contrairement au français (« plus... que »).
\end{enumerate}
\cle{Structure à retenir} : \zh{中国的移动支付比喀麦隆的普及得多。} (Zhōngguó de yídòng zhīfù bǐ Kāmàilóng de pǔjí de duō.) — « Le paiement mobile est bien plus répandu en Chine qu'au Cameroun. »
\end{methode}

\begin{exoresolu}[Comparer WeChat et WhatsApp / Orange Money]
\textbf{Énoncé :} Traduis en chinois les phrases suivantes en utilisant la structure de comparaison adaptée.
\begin{enumerate}
    \item WeChat est plus complet que WhatsApp.
    \item Au Cameroun, Orange Money est aussi répandu que le Mobile Money au Ghana. (Utilise \zh{跟……一样}.)
    \item Les applications chinoises ne sont pas aussi chères que les applications occidentales. (Utilise \zh{没有……那么}.)
    \item Alipay est beaucoup plus pratique que l'argent liquide.
\end{enumerate}
\tcblower
\textbf{Solution détaillée :}
\begin{enumerate}
    \item \textbf{Analyse :} comparaison de supériorité. A = \zh{微信} (Wēixìn), B = WhatsApp, adjectif = \zh{全面} (quánmiàn, complet). Degré : « plus » $\rightarrow$ \zh{得多}.
    \textbf{Réponse :} \zh{微信比WhatsApp全面得多。} (Wēixìn bǐ WhatsApp quánmiàn de duō.)
    \item \textbf{Analyse :} égalité. A = \zh{喀麦隆的Orange Money}, B = \zh{加纳的移动货币} (Jiānà de yídòng huòbì). Adjectif = \zh{普及} (pǔjí, répandu). Structure : A 跟 B 一样 + adjectif.
    \textbf{Réponse :} \zh{喀麦隆的Orange Money跟加纳的移动货币一样普及。} (Kāmàilóng de Orange Money gēn Jiānà de yídòng huòbì yíyàng pǔjí.)
    \item \textbf{Analyse :} négation de l'égalité. A = \zh{中国的应用} (Zhōngguó de yìngyòng), B = \zh{西方的应用} (Xīfāng de yìngyòng). Adjectif = \zh{贵} (guì, cher). Structure : A 没有 B 那么 + adjectif.
    \textbf{Réponse :} \zh{中国的应用没有西方的应用那么贵。} (Zhōngguó de yìngyòng méiyǒu Xīfāng de yìngyòng nàme guì.)
    \item \textbf{Analyse :} supériorité forte. A = \zh{支付宝} (Zhīfùbǎo), B = \zh{现金} (xiànjīn, argent liquide). Adjectif = \zh{方便} (fāngbiàn). Degré : « beaucoup plus » $\rightarrow$ \zh{得多}.
    \textbf{Réponse :} \zh{支付宝比现金方便得多。} (Zhīfùbǎo bǐ xiànjīn fāngbiàn de duō.)
\end{enumerate}
\textbf{Conclusion :} respecter l'ordre \textbf{A + 比/跟/没有 + B + (一样/那么) + adjectif + (得多/一点儿)}. Ne jamais dire \zh{比B很……} : la présence de \zh{很} dans une comparaison avec \zh{比} est une faute classique.
\end{exoresolu}

\begin{methode}[Exprimer son avis et argumenter : \cle{觉得}, \cle{认为}, \cle{因为……所以……}, \cle{一方面……另一方面……}]
\textbf{Objectif :} rédiger un paragraphe d'opinion structuré sur l'impact des TIC chinoises (TikTok, Huawei, Transsion) ou camerounaises (applications locales, fintechs).
\begin{enumerate}
    \item \textbf{Introduire l'opinion} : \zh{我觉得} (wǒ juéde, ressenti personnel) ou \zh{我认为} (wǒ rènwéi, avis réfléchi et argumenté) + proposition complète.
    \item \textbf{Donner la cause} : \zh{因为} (yīnwèi) + raison, \zh{所以} (suǒyǐ) + conséquence. \emph{En chinois, on utilise souvent les deux ensemble dans la même phrase, contrairement au français où « parce que » et « donc » ne se combinent pas.}
    \item \textbf{Nuancer} : utiliser \zh{一方面……，另一方面……} (yīfāngmiàn..., lìng yīfāngmiàn...) pour présenter les avantages ET les inconvénients (ex. praticité vs vie privée).
    \item \textbf{Conclure} : \zh{因此} (yīncǐ, c'est pourquoi) / \zh{总之} (zǒngzhī, en résumé) + avis final ou suggestion (\zh{我建议} wǒ jiànyì, je suggère).
\end{enumerate}
\cle{Connecteurs logiques clés} : \zh{而且} (érqiě, de plus), \zh{但是} (dànshì, mais), \zh{不过} (búguò, cependant), \zh{例如} (lìrú, par exemple).
\end{methode}

\begin{exoresolu}[Rédaction d'un avis argumenté : le commerce en direct au Cameroun]
\textbf{Énoncé :} Rédige un court paragraphe (60-80 caractères) en chinois pour exprimer ton avis sur l'arrivée du commerce en direct (\zh{直播带货}) via les réseaux sociaux au Cameroun. Utilise \zh{我认为}, \zh{因为……所以……} et \zh{一方面……另一方面……}.
\tcblower
\textbf{Solution détaillée :}

\textbf{Plan de rédaction :}
\begin{enumerate}
    \item Opinion : c'est une opportunité, mais il faut des règles.
    \item Argument 1 (avantage) : aide les petits commerçants à vendre directement.
    \item Argument 2 (inconvénient) : risques de contrefaçon et de protection des données.
    \item Conclusion : suggestion d'encadrement.
\end{enumerate}

\textbf{Brouillon chinois (avec pinyin) :}

\zh{我认为直播带货来到喀麦隆是好事，但是也有缺点。} (Wǒ rènwéi zhíbò dàihuò láidào Kāmàilóng shì hǎoshì, dànshì yě yǒu quēdiǎn.)

\zh{一方面，它能帮助小商贩直接卖货，很方便；} (Yīfāngmiàn, tā néng bāngzhù xiǎo shāngfàn zhíjiē màihuò, hěn fāngbiàn ;)

\zh{另一方面，假货和隐私问题让人担心。} (Lìng yīfāngmiàn, jiǎhuò hé yǐnsī wèntí ràng rén dānxīn.)

\zh{所以我建议政府加强监管。} (Suǒyǐ wǒ jiànyì zhèngfǔ jiāqiáng jiānguǎn.)

\textbf{Vérification grammaticale :}
\begin{itemize}
    \item \zh{直播带货} (zhíbò dàihuò) : terme figé « vente en direct par streaming ».
    \item \zh{来到} (láidào) : verbe directionnel « arriver à ».
    \item \zh{让人担心} (ràng rén dānxīn) : construction causative « rendre inquiet ».
    \item La structure \zh{一方面……另一方面……} est bien respectée, avec le sujet \zh{它} exprimé dans le premier membre.
\end{itemize}

\textbf{Version finale soignée :}

\zh{我认为直播带货进入喀麦隆市场有利有弊。一方面，它让小商贩直接面对客户，降低了成本；另一方面，产品质量和数据安全难以保证。因此，我建议相关部门制定明确的法规。}

(Wǒ rènwéi zhíbò dàihuò jìnrù Kāmàilóng shìchǎng yǒulì yǒubì. Yīfāngmiàn, tā ràng xiǎo shāngfàn zhíjiē miànduì kèhù, jiàngdī le chéngběn ; lìng yīfāngmiàn, chǎnpǐn zhìliàng hé shùjù ānquán nányǐ bǎozhèng. Yīncǐ, wǒ jiànyì xiāngguān bùmén zhìdìng míngquè de fǎguī.)

« Je pense que l'entrée du commerce en direct sur le marché camerounais a des avantages et des inconvénients. D'un côté, il permet aux petits commerçants de faire face directement aux clients et réduit les coûts ; de l'autre, la qualité des produits et la sécurité des données sont difficiles à garantir. C'est pourquoi je suggère que les services concernés établissent des règles claires. »
\end{exoresolu}

\begin{methode}[Comprendre un texte socioculturel : repérer les marqueurs de faits, chiffres et contrastes]
\textbf{Objectif :} répondre à des questions de compréhension fine sur un article comparant l'écosystème TIC Chine / Cameroun (infrastructures, taux de pénétration, modèles économiques).
\begin{enumerate}
    \item \textbf{Lecture globale} : identifier le \cle{sujet} (qui/quoi), la \cle{source}, la \cle{date} et la \cle{thèse principale} (Chine leader / Cameroun en rattrapage / complémentarité).
    \item \textbf{Lecture détaillée — vocabulaire clé} : surligner les noms propres (marques : \zh{华为} Huáwéi, \zh{传音} Chuányīn/Transsion), les chiffres (\zh{占比} zhànbǐ, proportion ; \zh{普及率} pǔjílǜ, taux de pénétration), les verbes d'évolution (\zh{超过} chāoguò, dépasser ; \zh{缩小} suōxiǎo, réduire [un écart]).
    \item \textbf{Repérer les connecteurs de contraste} : \zh{虽然……但是……} (suīrán... dànshì..., bien que... mais...), \zh{然而} (rán'ér, cependant), \zh{与……不同} (yǔ... bùtóng, différent de...). Ils signalent les points de comparaison explicites.
    \item \textbf{Types de questions à l'examen} :
        \begin{itemize}
            \item \textit{Droit au texte} : relever l'information exacte (chiffre, nom).
            \item \textit{Inférence} : déduire une cause ou une conséquence non écrite explicitement.
            \item \textit{Vocabulaire en contexte} : deviner le sens d'un mot grâce au contexte.
        \end{itemize}
\end{enumerate}
\cle{Réflexe} : ne jamais répondre par « oui/non ». Citer le segment de texte (recopier les caractères) puis reformuler en français.
\end{methode}

\begin{exoresolu}[Compréhension de texte : « Transsion, le roi du mobile en Afrique »]
\textbf{Énoncé :} lis le texte suivant et réponds aux questions en français.

\begin{textebox}[Article simulé — Journal du numérique en Afrique]
\zh{传音控股是中国手机品牌 在非洲的传奇。虽然华为、小米在中国很有名，但在非洲，传音的市场占有率长期超过50\%。它的成功秘诀在于“本地化”：针对非洲用户喜欢自拍的习惯，研发了多摄像头、超强夜拍功能；针对电力不稳定，优化了续航和快充。此外，传音还投资当地组装厂，创造就业。然而，随着三星和新品牌进入，竞争加剧，传音正在布局智能家居和移动支付生态。}\\
(Chuányīn kònggǔ shì Zhōngguó shǒujī pǐnpái zài Fēizhōu de chuánqí. Suīrán Huáwéi, Xiǎomǐ zài Zhōngguó hěn yǒumíng, dàn zài Fēizhōu, Chuányīn de shìchǎng zhànyǒulǜ chángqī chāoguò bǎifēnzhī wǔshí. Tā de chénggōng mìjué zàiyú “běndìhuà” : zhēnduì Fēizhōu yònghù xǐhuan zìpāi de xíguàn, yánfā le duō shèxiàngtóu, chāoqiáng yèpāi gōngnéng ; zhēnduì diànlì bù wěndìng, yōuhuà le xùháng hé kuàichōng. Cǐwài, Chuányīn hái tóuzī dāngdì zǔzhuāngchǎng, chuàngzào jiùyè. Rán'ér, suízhe Sānxīng hé xīn pǐnpái jìnrù, jìngzhēng jiājù, Chuányīn zhèngzài bùjú zhìnéng jiājū hé yídòng zhīfù shēngtài.)\\
\textbf{Traduction :} Transsion Holdings est la légende des marques chinoises de téléphones en Afrique. Bien que Huawei et Xiaomi soient célèbres en Chine, en Afrique, la part de marché de Transsion dépasse durablement 50\%. Son secret de réussite réside dans la « localisation » : pour répondre à l'habitude des selfies des utilisateurs africains, elle a développé des appareils photo multiples et un mode nuit très puissant ; pour répondre à l'instabilité de l'électricité, elle a optimisé l'autonomie et la charge rapide. De plus, Transsion investit dans des usines d'assemblage locales et crée des emplois. Cependant, avec l'entrée de Samsung et de nouvelles marques, la concurrence s'intensifie, et Transsion se lance dans la maison connectée et l'écosystème du paiement mobile.
\end{textebox}

\textbf{Questions :}
\begin{enumerate}
    \item Quelle est la part de marché de Transsion en Afrique mentionnée dans le texte ?
    \item Cite deux adaptations techniques (« localisation ») faites pour le marché africain.
    \item Quel connecteur introduit le contraste entre la notoriété en Chine et la domination en Afrique ?
    \item Selon le texte, quel est l'impact socio-économique direct des usines d'assemblage ?
    \item Quel défi futur la marque doit-elle relever ?
\end{enumerate}
\tcblower
\textbf{Solution détaillée :}
\begin{enumerate}
    \item \textbf{Relevé texte :} \zh{市场占有率长期超过50\%}. \textbf{Réponse :} plus de 50\% (position majoritaire durable).
    \item \textbf{Relevé texte :} 1. \zh{研发了多摄像头、超强夜拍功能} (appareils photo multiples et mode nuit puissant, adaptés au goût des selfies) ; 2. \zh{优化了续航和快充} (grande autonomie et charge rapide, adaptées à l'électricité instable).
    \item \textbf{Analyse grammaticale :} la phrase commence par \zh{虽然……但……}. \textbf{Réponse :} \zh{虽然} (suīrán), dans la structure \zh{虽然……但是/但……}.
    \item \textbf{Relevé texte :} \zh{投资当地组装厂，创造就业}. \textbf{Réponse :} la création d'emplois locaux (et un début d'industrialisation).
    \item \textbf{Relevé texte :} \zh{竞争加剧……正在布局智能家居和移动支付生态}. \textbf{Réponse :} l'intensification de la concurrence (Samsung, nouvelles marques) l'oblige à diversifier son offre (maison connectée, paiement mobile) pour conserver ses clients.
\end{enumerate}
\textbf{Conclusion :} les réponses sont \textbf{ancrées dans le texte} (caractères recopiés). L'inférence (question 5) s'appuie sur le connecteur \zh{然而} (rán'ér) pour lier concurrence et stratégie de diversification.
\end{exoresolu}

\begin{methode}[Préparer et présenter un mini-exposé oral (3 min) : \cle{开头 – 主体 – 结尾}]
\textbf{Objectif :} présenter un outil TIC chinois ou camerounais (ex. \zh{抖音} Dǒuyīn / TikTok, \zh{微信} Wēixìn, Orange Money, une application camerounaise) de façon structurée, fluide et avec une prononciation soignée (tons).
\begin{enumerate}
    \item \textbf{Choisir l'angle} : comparatif (Chine vs Cameroun) ou monographique (une réussite). Titre accrocheur : \zh{从微信看中喀移动支付差异} (Cóng Wēixìn kàn Zhōng-Kā yídòng zhīfù chāyì, « Les différences de paiement mobile Chine-Cameroun vues à travers WeChat »).
    \item \textbf{Structurer le discours} (fiches de mots-clés, pas de texte lu) :
        \begin{itemize}
            \item \textbf{开头} (kāitóu, ouverture, 30 s) : accroche (chiffre ou fait), sujet, plan annoncé (\zh{今天我要介绍……，分三点}).
            \item \textbf{主体} (zhǔtǐ, développement, 2 min) : 2 ou 3 points logiques (histoire/fonctions, modèle économique, impact social). Utiliser \zh{第一，……第二，……最后，……}.
            \item \textbf{结尾} (jiéwěi, conclusion, 30 s) : synthèse, avis personnel (\zh{我觉得……}), ouverture (\zh{未来……}), remerciements (\zh{谢谢大家}).
        \end{itemize}
    \item \textbf{Préparer le lexique technique} : écrire les mots-clés en caractères + pinyin sur une fiche (ex. \zh{小程序} xiǎochéngxù, mini-programme ; \zh{红包} hóngbāo, enveloppe rouge ; \zh{带货} dàihuò, vente en direct). S'entraîner sur les tons.
    \item \textbf{Supports visuels} : une diapositive = un mot-clé + une image. Pas de phrases entières en chinois à l'écran.
    \item \textbf{Répétition chronométrée} : parler \textbf{à voix haute}, s'enregistrer, vérifier la fluidité, les tons et le contact visuel.
\end{enumerate}
\cle{Phrases de liaison orales} : \zh{接下来} (jiēxiàlái, ensuite), \zh{举个例子} (jǔ gè lìzi, par exemple), \zh{值得注意的是} (zhíde zhùyì de shì, il convient de noter que).
\end{methode}

\begin{exoresolu}[Simulation d'exposé : « Orange Money vs WeChat Pay : deux modèles »]
\textbf{Énoncé :} tu dois présenter oralement (3 min) la différence de modèle entre Orange Money (Cameroun) et WeChat Pay (Chine). Rédige le \textbf{plan détaillé (mots-clés en chinois + pinyin)} et le \textbf{script de l'introduction et de la conclusion} (phrases complètes).
\tcblower
\textbf{Solution détaillée :}

\textbf{1. Plan détaillé (fiche de mots-clés) :}
\begin{center}
\begin{tabularx}{\linewidth}{|X|X|X|}
\hline
\rowcolor{popblueL} \textbf{Étape} & \textbf{Mots-clés (caractères + pinyin)} & \textbf{Idée directrice (français)} \\
\hline
\textbf{开头} & \zh{移动支付｜普及率｜对比} (yídòng zhīfù | pǔjílǜ | duìbǐ) & Fait d'accroche : le paiement mobile est quasi universel en Chine, en progression au Cameroun. Annonce du plan : 3 différences majeures. \\
\hline
\textbf{主体 1 : usages} & \zh{全场景｜社交+支付｜闭环} (quán chǎngjǐng | shèjiāo + zhīfù | bìhuán) & WeChat : écosystème tout-en-un (messagerie, achats, transport, démarches). Orange Money : centré sur le transfert, le retrait et le paiement de factures. \\
\hline
\textbf{主体 2 : modèle éco.} & \zh{免费转账｜流量变现｜手续费收入} (miǎnfèi zhuǎnzhàng | liúliàng biànxiàn | shǒuxùfèi shōurù) & WeChat : transferts gratuits entre particuliers, revenus via les services et la publicité. Orange Money : frais sur retraits et transferts (modèle proche du bancaire). \\
\hline
\textbf{主体 3 : inclusion} & \zh{银行账户｜实名制｜无银行账户} (yínháng zhànghù | shímíngzhì | wú yínháng zhànghù) & WeChat Pay : compte lié à une carte bancaire (public bancarisé). Orange Money : compte adossé au numéro de téléphone (public non bancarisé). \\
\hline
\textbf{结尾} & \zh{互补｜数字鸿沟｜合作前景} (hǔbǔ | shùzì hónggōu | hézuò qiánjǐng) & Synthèse : deux modèles complémentaires. Avis : Orange Money réduit la fracture numérique. Avenir : interopérabilité ? Merci. \\
\hline
\end{tabularx}
\end{center}

\textbf{2. Script de l'introduction (oral, environ 30 s) :}

\zh{大家好。今天我想和大家分享一个有趣的现象：为什么中国人出门只带手机，而很多喀麦隆人还离不开现金和Orange Money？}

(Dàjiā hǎo. Jīntiān wǒ xiǎng hé dàjiā fēnxiǎng yí gè yǒuqù de xiànxiàng : wèishénme Zhōngguó rén chūmén zhǐ dài shǒujī, ér hěn duō Kāmàilóng rén hái líbukāi xiànjīn hé Orange Money ?)

\zh{在中国，移动支付非常普及，微信支付和支付宝占主导地位。在喀麦隆，Orange Money最受欢迎，但模式完全不同。}

(Zài Zhōngguó, yídòng zhīfù fēicháng pǔjí, Wēixìn zhīfù hé Zhīfùbǎo zhàn zhǔdǎo dìwèi. Zài Kāmàilóng, Orange Money zuì shòu huānyíng, dàn móshì wánquán bùtóng.)

\zh{今天我将从三个方面对比：使用场景、商业模式和金融包容性。}

(Jīntiān wǒ jiāng cóng sān gè fāngmiàn duìbǐ : shǐyòng chǎngjǐng, shāngyè móshì hé jīnróng bāoróngxìng.)

\textbf{3. Script de la conclusion (oral, environ 30 s) :}

\zh{总之，微信支付是“超级应用”，绑定银行卡，适合高度数字化的社会；Orange Money是基于电话卡的“轻钱包”，解决了没有银行账户人群的困难。}

(Zǒngzhī, Wēixìn zhīfù shì “chāojí yìngyòng”, bǎngdìng yínhángkǎ, shìhé gāodù shùzìhuà de shèhuì ; Orange Money shì jīyú diànhuàkǎ de “qīng qiánbāo”, jiějué le méiyǒu yínháng zhànghù rénqún de kùnnan.)

\zh{我觉得两者不分高下，而是互补。未来如果能实现跨境互通，对中喀贸易和留学生都是巨大的便利。}

(Wǒ juéde liǎng zhě bù fēn gāoxià, érshì hǔbǔ. Wèilái rúguǒ néng shíxiàn kuàjìng hùtōng, duì Zhōng-Kā màoyì hé liúxuéshēng dōu shì jùdà de biànlì.)

\zh{谢谢大家的聆听，欢迎提问。}

(Xièxie dàjiā de língtīng, huānyíng tíwèn.)

\textbf{Conseil prononciation :} attention aux tons de \zh{普及率} (pǔjílǜ, 3-2-4), \zh{包容性} (bāoróngxìng, 1-2-4), \zh{绑定} (bǎngdìng, 3-4). Rappel de sandhi : \zh{不} passe au 2\textsuperscript{e} ton devant un 4\textsuperscript{e} ton (\zh{不是} bú shì), et \zh{一} devient yí devant un 4\textsuperscript{e} ton (\zh{一个} yí gè).
\end{exoresolu}

\begin{methode}[Traduire un texte socioculturel (version/thème) : termes culturels, sigles et ordre des mots]
\textbf{Objectif :} traduire un paragraphe comparatif (français $\leftrightarrow$ chinois) en respectant la terminologie TIC, les noms propres et l'ordre des mots (compléments de lieu, de temps et de moyen avant le verbe).
\begin{enumerate}
    \item \textbf{Analyser le texte source} : repérer les \cle{réalités culturelles} (termes gardés ou expliqués : \zh{红包} hóngbāo, \zh{直播带货} zhíbò dàihuò), les \cle{sigles} (5G, IA $\rightarrow$ \zh{人工智能} réngōng zhìnéng), les \cle{tournures passives} françaises (« est utilisé par » $\rightarrow$ actif en chinois, ou \zh{被} bèi si la nuance de subi est forte).
    \item \textbf{Stratégie terminologique} :
        \begin{itemize}
            \item Marque connue : garder le nom officiel (\zh{华为} Huáwéi, \zh{传音} Chuányīn, Orange Money).
            \item Concept chinois spécifique : terme chinois + explication courte (\zh{小程序} xiǎochéngxù, mini-programme).
            \item Concept camerounais : traduction descriptive + nom propre (\zh{移动货币} yídòng huòbì pour « mobile money »).
        \end{itemize}
    \item \textbf{Réordonnancement syntaxique (français $\rightarrow$ chinois)} :
        \begin{itemize}
            \item Temps / lieu / moyen $\rightarrow$ \textbf{avant} le verbe. (Fr. « Il paie \textbf{avec son téléphone} \textbf{tous les jours} » $\rightarrow$ \zh{他每天用手机付款。} Tā měitiān yòng shǒujī fùkuǎn.)
            \item Subordonnée relative $\rightarrow$ \textbf{devant} le nom qu'elle qualifie, avec \zh{的}. (Fr. « l'application \textbf{que j'utilise} » $\rightarrow$ \zh{我用的应用} wǒ yòng de yìngyòng.)
        \end{itemize}
    \item \textbf{Relecture croisée} : vérifier les classificateurs (\zh{个} gè ; \zh{款} kuǎn pour une application ou un modèle de téléphone ; \zh{家} jiā pour une entreprise), les particules d'aspect (\zh{了} le, \zh{过} guo) et la cohérence globale.
\end{enumerate}
\cle{Piège classique} : « l'application est téléchargée par des millions d'utilisateurs » se rend plus naturellement par l'actif \zh{几百万用户下载了这个应用} (jǐ bǎiwàn yònghù xiàzài le zhè ge yìngyòng) que par le passif avec \zh{被}, réservé aux situations subies ou défavorables.
\end{methode}

\begin{exoresolu}[Thème : « L'écosystème Huawei au Cameroun »]
\textbf{Énoncé :} traduis le paragraphe suivant en chinois (caractères + pinyin).

\textbf{Texte source :}
« Au Cameroun, Huawei n'est pas seulement un vendeur de smartphones. Depuis 2018, l'entreprise a construit des centres de données et formé plus de 500 ingénieurs locaux grâce à son programme "Seeds for the Future" (未来种子). Bien que la 5G ne soit pas encore commercialisée, la 4G de Huawei couvre déjà une grande partie de la population. Les jeunes Camerounais apprécient particulièrement la solidité et le rapport qualité-prix des téléphones Tecno (marque de Transsion), mais Huawei reste la référence pour le haut de gamme. »

\tcblower
\textbf{Solution détaillée :}

\textbf{Analyse segment par segment :}
\begin{enumerate}
    \item \textit{Au Cameroun} $\rightarrow$ \zh{在喀麦隆} (zài Kāmàilóng) : complément de lieu en tête de phrase.
    \item \textit{Huawei n'est pas seulement un vendeur de smartphones} $\rightarrow$ \zh{华为不仅是智能手机销售商} (Huáwéi bùjǐn shì zhìnéng shǒujī xiāoshòushāng). Structure \zh{不仅……还/也……} (non seulement... mais aussi...).
    \item \textit{Depuis 2018} $\rightarrow$ \zh{自2018年以来} (zì 2018 nián yǐlái) : complément de temps avant le verbe.
    \item \textit{l'entreprise a construit des centres de données} $\rightarrow$ \zh{该企业建设了数据中心} (gāi qǐyè jiànshè le shùjù zhōngxīn). \zh{了} marque l'accompli.
    \item \textit{et formé plus de 500 ingénieurs locaux grâce à son programme...} $\rightarrow$ \zh{还通过“未来种子”计划培养了500多名当地工程师} (hái tōngguò “Wèilái Zhǒngzi” jìhuà péiyǎng le 500 duō míng dāngdì gōngchéngshī). \zh{通过} = « par le biais de » ; \zh{名} est le classificateur des personnes (registre soutenu) ; \zh{500多名} = « plus de 500 ».
    \item \textit{Bien que la 5G ne soit pas encore commercialisée} $\rightarrow$ \zh{虽然5G还没有商用} (suīrán 5G hái méiyǒu shāngyòng). \zh{商用} = « mise en service commerciale ».
    \item \textit{la 4G de Huawei couvre déjà une grande partie de la population} $\rightarrow$ \zh{华为的4G网络已经覆盖了大部分人口} (Huáwéi de 4G wǎngluò yǐjīng fùgài le dà bùfen rénkǒu).
    \item \textit{Les jeunes Camerounais apprécient particulièrement la solidité et le rapport qualité-prix des téléphones Tecno} $\rightarrow$ \zh{喀麦隆年轻人特别看重传音旗下Tecno手机的耐用性和性价比} (Kāmàilóng niánqīngrén tèbié kànzhòng Chuányīn qíxià Tecno shǒujī de nàiyòngxìng hé xìngjiàbǐ). \zh{旗下} qíxià = « filiale de » ; \zh{性价比} xìngjiàbǐ = « rapport qualité-prix » ; \zh{耐用性} nàiyòngxìng = « solidité ».
    \item \textit{mais Huawei reste la référence pour le haut de gamme} $\rightarrow$ \zh{但华为依然是高端机型的标杆} (dàn Huáwéi yīrán shì gāoduān jīxíng de biāogān). \zh{标杆} biāogān = « référence » ; \zh{高端机型} gāoduān jīxíng = « modèles haut de gamme ».
\end{enumerate}

\textbf{Version finale assemblée :}

\zh{在喀麦隆，华为不仅是智能手机销售商。自2018年以来，该企业建设了数据中心，还通过“未来种子”计划培养了500多名当地工程师。虽然5G还没有商用，但华为的4G网络已经覆盖了大部分人口。喀麦隆年轻人特别看重传音旗下Tecno手机的耐用性和性价比，但华为依然是高端机型的标杆。}

(Zài Kāmàilóng, Huáwéi bùjǐn shì zhìnéng shǒujī xiāoshòushāng. Zì 2018 nián yǐlái, gāi qǐyè jiànshè le shùjù zhōngxīn, hái tōngguò “Wèilái Zhǒngzi” jìhuà péiyǎng le 500 duō míng dāngdì gōngchéngshī. Suīrán 5G hái méiyǒu shāngyòng, dàn Huáwéi de 4G wǎngluò yǐjīng fùgài le dà bùfen rénkǒu. Kāmàilóng niánqīngrén tèbié kànzhòng Chuányīn qíxià Tecno shǒujī de nàiyòngxìng hé xìngjiàbǐ, dàn Huáwéi yīrán shì gāoduān jīxíng de biāogān.)

\textbf{Points de vigilance :}
\begin{itemize}
    \item \zh{销售商} (xiāoshòushāng, distributeur) convient mieux que \zh{卖家} (màijiā, vendeur particulier).
    \item \zh{在喀麦隆} (complément de lieu) se place en tout début de phrase.
    \item \zh{500多名} : \zh{多} placé après le nombre exprime « plus de ».
    \item Pas de \zh{被} inutile : le chinois préfère l'actif quand l'agent est connu.
\end{itemize}
\end{exoresolu}

\begin{attention}[Les 5 erreurs qui coûtent des points à l'examen — spécial TIC et comparatif]
\begin{enumerate}
    \item \textbf{Tons et homophones « tech » :} confondre \zh{数} (shù, chiffre/donnée) et \zh{树} (shù, arbre) ; \zh{码} (mǎ, code) et \zh{马} (mǎ, cheval). \emph{Ex. : \zh{扫码} (sǎomǎ, scanner un code) $\neq$ \zh{扫马} (sǎomǎ, « balayer le cheval »).} \textbf{Règle :} apprendre le vocabulaire technique par paires minimales tonales.
    \item \textbf{Caractères visuellement proches :} \zh{程} (chéng, programme — clé \cle{禾}, le céréale) vs \zh{呈} (chéng, présenter — clé \cle{口}, la bouche). \emph{Piège : \zh{程序} (chéngxù, programme) s'écrit avec la clé \cle{禾}, pas avec \cle{口}.}
    \item \textbf{Ordre des mots : compléments de lieu et de moyen AVANT le verbe.} \emph{Erreur :} \zh{我用微信付款了在超市。} \textbf{Correct :} \zh{我在超市用微信付款了。} (Wǒ zài chāoshì yòng Wēixìn fùkuǎn le.) \textbf{Règle d'or :} sujet + quand + où + avec quoi/comment + verbe.
    \item \textbf{Mots-outils : \cle{了} vs \cle{过}.} \emph{Erreur :} \zh{我下载了微信两年。} \textbf{Correct :} \zh{我下载微信两年了。} (Wǒ xiàzài Wēixìn liǎng nián le — je l'utilise depuis deux ans) ou \zh{我下载过微信。} (Wǒ xiàzài guo Wēixìn — je l'ai déjà téléchargé [expérience]). \zh{了} en fin de phrase = changement d'état ou situation qui dure ; \zh{过} = expérience vécue.
    \item \textbf{Classificateurs et calques du français :} on dit \zh{一款应用} (yì kuǎn yìngyòng, une application) et \zh{一条微信} (yì tiáo Wēixìn, un message WeChat). \emph{Calque à éviter :} « faire une vidéo » ne se traduit pas par \zh{做视频} mais par \zh{拍视频} (pāi shìpín, filmer) ou \zh{发视频} (fā shìpín, publier une vidéo). « Partager » (un contenu) = \zh{分享} (fēnxiǎng), et non \zh{共享} (gòngxiǎng, mise en commun de ressources).
\end{enumerate}
\end{attention}

\newpage

\coursec{Exercices}

\courssub{Je vérifie mes connaissances}

\begin{exercice}[QCM : les TIC en Chine et au Cameroun]
Entoure la bonne réponse.
\begin{enumerate}
\item 微信 (Wēixìn) est une application :
\begin{enumerate}
\item chinoise \item américaine \item camerounaise
\end{enumerate}
\item Le « mobile money » se dit en chinois :
\begin{enumerate}
\item 手机支付 (shǒujī zhīfù) \item 手机品牌 (shǒujī pǐnpái) \item 网络新闻 (wǎngluò xīnwén)
\end{enumerate}
\item 华为 (Huáwéi) est :
\begin{enumerate}
\item une marque de téléphone chinoise \item une ville du Cameroun \item une application de paiement
\end{enumerate}
\item 方便 (fāngbiàn) signifie :
\begin{enumerate}
\item cher \item pratique \item lent
\end{enumerate}
\end{enumerate}
\end{exercice}

\begin{exercice}[Vrai ou faux ? Justifie ta réponse]
Dis si chaque phrase est vraie ou fausse, puis justifie en une phrase en français.
\begin{enumerate}
\item 在中国，很多人用微信支付。 (Zài Zhōngguó, hěn duō rén yòng Wēixìn zhīfù.)
\item 喀麦隆没有中国的手机品牌。 (Kāmàilóng méiyǒu Zhōngguó de shǒujī pǐnpái.)
\item 手机支付很方便。 (Shǒujī zhīfù hěn fāngbiàn.)
\item 只有年轻人才用手机。 (Zhǐyǒu niánqīngrén cái yòng shǒujī.)
\end{enumerate}
\end{exercice}

\begin{exercice}[Vocabulaire : retrouve la traduction]
Associe chaque mot chinois à sa traduction française.
\begin{center}
\begin{tabularx}{\linewidth}{|X|X|}
\hline
\rowcolor{popblueL} Mot chinois & Traductions possibles \\
\hline
1. 品牌 (pǐnpái) & a. Internet / le réseau \\
\hline
2. 网络 (wǎngluò) & b. la marque \\
\hline
3. 支付 (zhīfù) & c. populaire, à la mode \\
\hline
4. 流行 (liúxíng) & d. payer / le paiement \\
\hline
5. 应用 (yìngyòng) & e. l'application \\
\hline
\end{tabularx}
\end{center}
\end{exercice}

\begin{exercice}[Pinyin : associe les sons aux caractères]
Associe chaque pinyin au bon caractère.
\begin{enumerate}
\item shǒujī \hspace{1cm} a. 电脑
\item diànnǎo \hspace{1cm} b. 手机
\item ruǎnjiàn \hspace{1cm} c. 短信
\item duǎnxìn \hspace{1cm} d. 软件
\end{enumerate}
Puis écris la traduction française de chaque mot.
\end{exercice}

\courssub{Je m'entraîne}

\begin{exercice}[Traduction vers le français]
Traduis les phrases suivantes en français.
\begin{enumerate}
\item 在喀麦隆，很多人用中国品牌的手机。 (Zài Kāmàilóng, hěn duō rén yòng Zhōngguó pǐnpái de shǒujī.)
\item 移动支付又便宜又方便。 (Yídòng zhīfù yòu piányi yòu fāngbiàn.)
\item 我觉得微信比短信好用。 (Wǒ juéde Wēixìn bǐ duǎnxìn hǎoyòng.)
\item 现在，网络在我们生活中很重要。 (Xiànzài, wǎngluò zài wǒmen shēnghuó zhōng hěn zhòngyào.)
\end{enumerate}
\end{exercice}

\begin{exercice}[Texte à trous : banque de mots]
Complète le texte avec les mots de la banque : \zh{品牌、网络、支付、方便、流行} (pǐnpái, wǎngluò, zhīfù, fāngbiàn, liúxíng).

\begin{textebox}[Texte à trous]
在雅温得，很多年轻人用中国\underline{\hspace{2cm}}的手机。\\
(Zài Yǎwēndé, hěn duō niánqīngrén yòng Zhōngguó \underline{\hspace{2cm}} de shǒujī.)\\
手机\underline{\hspace{2cm}}越来越流行，因为很\underline{\hspace{2cm}}。\\
(Shǒujī \underline{\hspace{2cm}} yuèláiyuè liúxíng, yīnwèi hěn \underline{\hspace{2cm}}.)\\
有了\underline{\hspace{2cm}}，我们可以上网、聊天、买东西。\\
(Yǒule \underline{\hspace{2cm}}, wǒmen kěyǐ shàngwǎng, liáotiān, mǎi dōngxi.)\\
现在手机支付在喀麦隆也很\underline{\hspace{2cm}}。\\
(Xiànzài shǒujī zhīfù zài Kāmàilóng yě hěn \underline{\hspace{2cm}}.)
\end{textebox}
\end{exercice}

\begin{exercice}[Remets les mots dans l'ordre]
Remets les mots dans le bon ordre pour faire des phrases correctes, puis traduis-les en français.
\begin{enumerate}
\item 用 / 我 / 手机 / 中国品牌的
\item 比 / 微信支付 / 现金 / 方便
\item 觉得 / 我 / 移动支付 / 很 / 有用
\item 在 / 流行 / 手机支付 / 杜阿拉 / 很
\end{enumerate}
\end{exercice}

\begin{exercice}[Transformation : comparer avec 比]
Transforme chaque paire de phrases en une seule phrase de comparaison avec \zh{比} (bǐ), comme dans l'exemple.

Exemple : 微信很好用。短信不太好用。 → 微信比短信好用。 (Wēixìn bǐ duǎnxìn hǎoyòng.)
\begin{enumerate}
\item 手机支付很方便。现金不太方便。
\item 华为手机很便宜。苹果手机很贵。
\item 网络很快。写信很慢。
\item 这个应用很流行。那个应用不太流行。
\end{enumerate}
\end{exercice}

\courssub{Je me prépare au Bac}

\begin{exercice}[Compréhension écrite — type Bac]
Lis le texte suivant, puis réponds aux questions.

\begin{textebox}[Le téléphone portable au Cameroun et en Chine]
在喀麦隆，手机越来越重要。很多人用中国品牌的手机，因为又便宜又好用。\\
(Zài Kāmàilóng, shǒujī yuèláiyuè zhòngyào. Hěn duō rén yòng Zhōngguó pǐnpái de shǒujī, yīnwèi yòu piányi yòu hǎoyòng.)\\
在杜阿拉和雅温得，手机支付很流行：人们用手机付钱、转钱。\\
(Zài Dùālālā hé Yǎwēndé, shǒujī zhīfù hěn liúxíng : rénmen yòng shǒujī fùqián, zhuǎnqián.)\\
这很方便，也很安全。\\
(Zhè hěn fāngbiàn, yě hěn ānquán.)\\
在中国，人们常用微信。\\
(Zài Zhōngguó, rénmen cháng yòng Wēixìn.)\\
微信可以聊天，也可以支付。\\
(Wēixìn kěyǐ liáotiān, yě kěyǐ zhīfù.)\\
很多中国人出门不带现金，只带手机。\\
(Hěn duō Zhōngguórén chūmén bù dài xiànjīn, zhǐ dài shǒujī.)\\
我觉得，喀麦隆和中国一样，手机改变了人们的生活。\\
(Wǒ juéde, Kāmàilóng hé Zhōngguó yíyàng, shǒujī gǎibiànle rénmen de shēnghuó.)
\end{textebox}

Questions :
\begin{enumerate}
\item Vrai ou faux ? Justifie avec une phrase du texte : « Les téléphones de marques chinoises sont chers au Cameroun. »
\item Vrai ou faux ? Justifie : « En Chine, tout le monde sort avec beaucoup d'argent liquide. »
\item 在喀麦隆，人们用手机支付做什么？ (Réponds en chinois.)
\item 微信可以做什么？ (Réponds en chinois : deux choses.)
\item Question de langue : relève dans le texte une comparaison avec \zh{比} ou une structure \zh{又……又……}, recopie-la et traduis-la en français.
\end{enumerate}
\end{exercice}

\begin{exercice}[Traduction — thème et version]
\textbf{Version.} Traduis en français :
\begin{enumerate}
\item 在喀麦隆，很多人用中国品牌的手机，因为又便宜又好用。\\
(Zài Kāmàilóng, hěn duō rén yòng Zhōngguó pǐnpái de shǒujī, yīnwèi yòu piányi yòu hǎoyòng.)
\item 在中国，微信非常流行，人们用它聊天和支付。\\
(Zài Zhōngguó, Wēixìn fēicháng liúxíng, rénmen yòng tā liáotiān hé zhīfù.)
\end{enumerate}
\textbf{Thème.} Traduis en chinois (caractères + pinyin) :
\begin{enumerate}
\item Je pense que le paiement mobile est très pratique.
\item À Douala, beaucoup de jeunes utilisent Internet tous les jours.
\end{enumerate}
\end{exercice}

\begin{exercice}[Expression écrite et orale — type Bac]
\textbf{A. Expression écrite.} Rédige un texte de \textbf{6 à 8 phrases} en chinois (environ 60 à 80 caractères) sur le sujet : \zh{我最喜欢的应用} (Wǒ zuì xǐhuan de yìngyòng — « Mon application préférée »).

Consignes obligatoires :
\begin{itemize}
\item utilise au moins une comparaison avec \zh{比} (bǐ) ;
\item donne ton avis avec \zh{我觉得} (wǒ juéde) ;
\item emploie au moins trois mots du vocabulaire de la leçon : \zh{品牌、网络、支付、方便、流行、应用}.
\end{itemize}

\textbf{B. Expression orale.} Prépare un mini-exposé de \textbf{2 minutes} comparant un usage numérique au Cameroun et en Chine (exemple : le mobile money au Cameroun et WeChat Pay en Chine). Prévois : une introduction, deux points de comparaison, ta conclusion personnelle. Le jury te posera ensuite deux questions en chinois, par exemple : \zh{你常用什么应用？} (Nǐ cháng yòng shénme yìngyòng ?) et \zh{你觉得手机支付安全吗？} (Nǐ juéde shǒujī zhīfù ānquán ma ?).
\end{exercice}

\newpage

\coursec{Corrigés des exercices}

\begin{corrige}[Exercice 1 — QCM : les TIC en Chine et au Cameroun]
\begin{enumerate}
\item \textbf{Réponse a.} 微信 (Wēixìn, « WeChat ») est une application chinoise, créée par l'entreprise 腾讯 (Téngxùn, Tencent). Elle permet de discuter, de payer et de partager des photos.
\item \textbf{Réponse a.} Le « mobile money » se dit \zh{手机支付} (shǒujī zhīfù, littéralement « paiement par téléphone ») ou \zh{移动支付} (yídòng zhīfù, « paiement mobile »). 手机品牌 (shǒujī pǐnpái) signifie « marque de téléphone » et 网络新闻 (wǎngluò xīnwén) « les informations sur Internet ».
\item \textbf{Réponse a.} 华为 (Huáwéi) est une marque de téléphone chinoise, très présente au Cameroun. Ses téléphones sont appréciés parce qu'ils sont 又便宜又好用 (yòu piányi yòu hǎoyòng, « à la fois bon marché et pratiques »).
\item \textbf{Réponse b.} 方便 (fāngbiàn) signifie « pratique, commode ». « Cher » se dit 贵 (guì) et « lent » 慢 (màn).
\end{enumerate}
\textbf{Point de vigilance :} ne pas confondre 支付 (zhīfù, « payer ») et 品牌 (pǐnpái, « marque ») : les deux mots apparaissent souvent ensemble dans le thème des TIC.
\end{corrige}

\begin{corrige}[Exercice 2 — Vrai ou faux ? Justifie ta réponse]
\begin{enumerate}
\item \textbf{Vrai.} En Chine, le paiement par WeChat est très répandu : beaucoup de Chinois paient leurs achats avec leur téléphone, même au marché.
\item \textbf{Faux.} Au Cameroun, on trouve beaucoup de marques de téléphone chinoises, comme Huawei, Tecno ou Xiaomi, surtout dans les villes comme Douala et Yaoundé.
\item \textbf{Vrai.} Le paiement par téléphone est pratique : on peut payer ou transférer de l'argent rapidement, sans se déplacer à la banque.
\item \textbf{Faux.} Aujourd'hui, les personnes de tous les âges utilisent le téléphone portable : les jeunes, les adultes et même beaucoup de personnes âgées.
\end{enumerate}
\textbf{Point de vigilance :} dans la phrase 4, la structure 只有……才…… (zhǐyǒu... cái..., « seulement... que... ») exprime une restriction : « seuls les jeunes utilisent le téléphone ». C'est cette généralisation qui rend la phrase fausse.
\end{corrige}

\begin{corrige}[Exercice 3 — Vocabulaire : retrouve la traduction]
\begin{enumerate}
\item 品牌 (pǐnpái) → \textbf{b.} la marque
\item 网络 (wǎngluò) → \textbf{a.} Internet / le réseau
\item 支付 (zhīfù) → \textbf{d.} payer / le paiement
\item 流行 (liúxíng) → \textbf{c.} populaire, à la mode
\item 应用 (yìngyòng) → \textbf{e.} l'application
\end{enumerate}
\textbf{Point de vigilance :} 流行 (liúxíng) est un adjectif (« être à la mode ») : on dit 手机支付很流行 (shǒujī zhīfù hěn liúxíng, « le paiement mobile est très populaire »), jamais avec 很 + verbe d'action.
\end{corrige}

\begin{corrige}[Exercice 4 — Pinyin : associe les sons aux caractères]
\begin{enumerate}
\item shǒujī → \textbf{b.} 手机 : le téléphone (portable)
\item diànnǎo → \textbf{a.} 电脑 : l'ordinateur
\item ruǎnjiàn → \textbf{d.} 软件 : le logiciel
\item duǎnxìn → \textbf{c.} 短信 : le SMS, le texto
\end{enumerate}
\textbf{Point de vigilance :} attention aux tons : 手机 (shǒujī, 3\textsuperscript{e} + 1\textsuperscript{er} ton) ne se prononce pas comme 收集 (shōují, « collectionner »). De même, 电脑 (diànnǎo) porte deux 3\textsuperscript{e} tons : dans la prononciation, le premier devient un 2\textsuperscript{e} ton (diánnǎo), mais l'écriture pinyin garde les tons d'origine.
\end{corrige}

\begin{corrige}[Exercice 5 — Traduction vers le français]
\begin{enumerate}
\item 在喀麦隆，很多人用中国品牌的手机。 → « Au Cameroun, beaucoup de gens utilisent des téléphones de marques chinoises. »
\item 移动支付又便宜又方便。 → « Le paiement mobile est à la fois bon marché et pratique. »
\item 我觉得微信比短信好用。 → « Je trouve que WeChat est plus pratique que les SMS. »
\item 现在，网络在我们生活中很重要。 → « Aujourd'hui, Internet est très important dans notre vie. »
\end{enumerate}
\textbf{Points de vigilance :}
\begin{itemize}
\item 又……又…… (yòu... yòu...) se traduit par « à la fois... et... » ; il ne s'emploie qu'avec des adjectifs ou des verbes psychologiques.
\item 比 (bǐ) se traduit par « plus... que » : A + 比 + B + adjectif.
\item 在……中 (zài... zhōng) signifie « dans, au sein de » : 在我们生活中, « dans notre vie ».
\end{itemize}
\end{corrige}

\begin{corrige}[Exercice 6 — Texte à trous : banque de mots]
Texte complété :
\begin{enumerate}
\item 在雅温得，很多年轻人用中国\textbf{品牌}的手机。 (Zài Yǎwēndé, hěn duō niánqīngrén yòng Zhōngguó pǐnpái de shǒujī.) → « À Yaoundé, beaucoup de jeunes utilisent des téléphones de marques chinoises. »
\item 手机\textbf{支付}越来越流行，因为很\textbf{方便}。 (Shǒujī zhīfù yuèláiyuè liúxíng, yīnwèi hěn fāngbiàn.) → « Le paiement par téléphone est de plus en plus populaire, parce que c'est très pratique. »
\item 有了\textbf{网络}，我们可以上网、聊天、买东西。 (Yǒule wǎngluò, wǒmen kěyǐ shàngwǎng, liáotiān, mǎi dōngxi.) → « Grâce à Internet, nous pouvons aller en ligne, discuter et faire des achats. »
\item 现在手机支付在喀麦隆也很\textbf{流行}。 (Xiànzài shǒujī zhīfù zài Kāmàilóng yě hěn liúxíng.) → « Aujourd'hui, le paiement mobile est aussi très répandu au Cameroun. »
\end{enumerate}
\textbf{Justification :} 中国品牌 (marque chinoise) exige le nom 品牌 ; 手机支付 (paiement par téléphone) exige 支付 ; après 很, il faut un adjectif, donc 方便 ; 有了 + nom exige 网络 ; la dernière phrase décrit un phénomène à la mode, donc 流行.
\end{corrige}

\begin{corrige}[Exercice 7 — Remets les mots dans l'ordre]
\begin{enumerate}
\item 我用中国品牌的手机。 (Wǒ yòng Zhōngguó pǐnpái de shǒujī.) → « J'utilise un téléphone de marque chinoise. » Ordre : Sujet + Verbe + Complément (le déterminant 中国品牌的 se place avant le nom 手机).
\item 微信支付比现金方便。 (Wēixìn zhīfù bǐ xiànjīn fāngbiàn.) → « Le paiement WeChat est plus pratique que l'argent liquide. » Structure : A + 比 + B + adjectif.
\item 我觉得移动支付很有用。 (Wǒ juéde yídòng zhīfù hěn yǒuyòng.) → « Je trouve que le paiement mobile est très utile. » 觉得 introduit une proposition subordonnée, sans mot de liaison.
\item 在杜阿拉，手机支付很流行。 (Zài Dùālālā, shǒujī zhīfù hěn liúxíng.) → « À Douala, le paiement mobile est très populaire. » Le complément de lieu avec 在 se place en tête de phrase (ou avant le verbe).
\end{enumerate}
\textbf{Point de vigilance :} erreur fréquente des francophones : placer le lieu après le verbe comme en français. En chinois, on dit 在杜阿拉手机支付很流行 ou 手机支付在杜阿拉很流行, jamais le lieu en fin de phrase.
\end{corrige}

\begin{corrige}[Exercice 8 — Transformation : comparer avec 比]
\begin{enumerate}
\item 手机支付比现金方便。 (Shǒujī zhīfù bǐ xiànjīn fāngbiàn.) → « Le paiement par téléphone est plus pratique que l'argent liquide. »
\item 华为手机比苹果手机便宜。 (Huáwéi shǒujī bǐ Píngguǒ shǒujī piányi.) → « Les téléphones Huawei sont moins chers que les iPhone. »
\item 网络比写信快。 (Wǎngluò bǐ xiěxìn kuài.) → « Internet est plus rapide que la lettre. »
\item 这个应用比那个应用流行。 (Zhège yìngyòng bǐ nàge yìngyòng liúxíng.) → « Cette application est plus populaire que celle-là. »
\end{enumerate}
\textbf{Justification grammaticale :} la structure de comparaison est \cle{A + 比 + B + adjectif}. L'adjectif seul suffit : on ne met pas 很 après 比 (on ne dit pas 比现金很方便). Pour nuancer, on peut ajouter 一点儿 (un peu) ou 得多 (beaucoup) après l'adjectif : 手机支付比现金方便得多 (« beaucoup plus pratique »).
\end{corrige}

\begin{corrige}[Exercice 9 — Compréhension écrite — type Bac]
\textbf{Barème indicatif (sur 10 points) :} questions 1 et 2 : 2 points chacune (1 point pour vrai/faux, 1 point pour la justification) ; questions 3 et 4 : 2 points chacune ; question 5 : 2 points.

\begin{enumerate}
\item \textbf{Faux.} Le texte dit : 很多人用中国品牌的手机，因为又便宜又好用。 (Hěn duō rén yòng Zhōngguó pǐnpái de shǒujī, yīnwèi yòu piányi yòu hǎoyòng.) → « Beaucoup de gens utilisent des téléphones de marques chinoises parce qu'ils sont à la fois bon marché et pratiques. » Ils sont donc bon marché, pas chers.
\item \textbf{Faux.} Le texte dit : 很多中国人出门不带现金，只带手机。 (Hěn duō Zhōngguórén chūmén bù dài xiànjīn, zhǐ dài shǒujī.) → « Beaucoup de Chinois ne prennent pas d'argent liquide quand ils sortent, seulement leur téléphone. »
\item 人们用手机付钱、转钱。 (Rénmen yòng shǒujī fùqián, zhuǎnqián.) → « Les gens utilisent le téléphone pour payer et transférer de l'argent. »
\item 微信可以聊天，也可以支付。 (Wēixìn kěyǐ liáotiān, yě kěyǐ zhīfù.) → « WeChat permet de discuter et aussi de payer. »
\item Exemple de structure relevée : 又便宜又好用 (yòu piányi yòu hǎoyòng) → « à la fois bon marché et pratique ». Autre réponse possible : 手机改变了人们的生活 (shǒujī gǎibiànle rénmen de shēnghuó) → « le téléphone a changé la vie des gens » (avec 了 de changement d'état).
\end{enumerate}
\textbf{Points de vigilance :} pour les questions « réponds en chinois », recopier la phrase du texte est accepté si elle répond exactement à la question ; ne pas traduire mot à mot le français. Pour la question 5, recopier la structure complète avec sa traduction : une structure sans traduction ne vaut que la moitié des points.
\end{corrige}

\begin{corrige}[Exercice 10 — Traduction — thème et version]
\textbf{Version.}
\begin{enumerate}
\item 在喀麦隆，很多人用中国品牌的手机，因为又便宜又好用。 → « Au Cameroun, beaucoup de gens utilisent des téléphones de marques chinoises, parce qu'ils sont à la fois bon marché et pratiques. »
\item 在中国，微信非常流行，人们用它聊天和支付。 → « En Chine, WeChat est très populaire : les gens l'utilisent pour discuter et payer. »
\end{enumerate}
\textbf{Thème.}
\begin{enumerate}
\item « Je pense que le paiement mobile est très pratique. » → 我觉得移动支付很方便。 (Wǒ juéde yídòng zhīfù hěn fāngbiàn.)
\item « À Douala, beaucoup de jeunes utilisent Internet tous les jours. » → 在杜阿拉，很多年轻人每天用网络。 (Zài Dùālālā, hěn duō niánqīngrén měitiān yòng wǎngluò.)
\end{enumerate}
\textbf{Points de vigilance :}
\begin{itemize}
\item Après 我觉得, pas de mot de liaison : on dit 我觉得移动支付很方便, jamais 我觉得 que...
\item Le complément de temps 每天 (měitiān, « chaque jour ») se place avant le verbe 用, après le sujet.
\item « Utiliser Internet » se dit 用网络 ou 上网 (shàngwǎng, « aller sur Internet ») : 很多年轻人每天上网 est également correct.
\end{itemize}
\textbf{Barème indicatif (sur 10 points) :} 2,5 points par phrase ; enlever 0,5 point par faute de caractère ou d'ordre des mots, 0,25 point par faute de pinyin.
\end{corrige}

\begin{corrige}[Exercice 11 — Expression écrite et orale — type Bac]
\textbf{A. Expression écrite — proposition de rédaction modèle.}

我最喜欢的应用是微信。\\
(Wǒ zuì xǐhuan de yìngyòng shì Wēixìn.)\\
« Mon application préférée est WeChat. »

它比短信好用，也比打电话便宜。\\
(Tā bǐ duǎnxìn hǎoyòng, yě bǐ dǎ diànhuà piányi.)\\
« Elle est plus pratique que les SMS et moins chère que les appels. »

有了网络，我可以用微信和朋友聊天、发照片。\\
(Yǒule wǎngluò, wǒ kěyǐ yòng Wēixìn hé péngyou liáotiān, fā zhàopiàn.)\\
« Grâce à Internet, je peux discuter avec mes amis et envoyer des photos. »

在中国，很多人还用微信支付，这很方便。\\
(Zài Zhōngguó, hěn duō rén hái yòng Wēixìn zhīfù, zhè hěn fāngbiàn.)\\
« En Chine, beaucoup de gens paient aussi avec WeChat, c'est très pratique. »

我觉得，这个应用很有用，也很流行。\\
(Wǒ juéde, zhège yìngyòng hěn yǒuyòng, yě hěn liúxíng.)\\
« Je trouve que cette application est très utile et très populaire. »

Vérification des consignes : une comparaison avec 比 (它比短信好用), un avis avec 我觉得, et quatre mots du vocabulaire de la leçon (应用、网络、支付、方便、流行). Environ 75 caractères.

\textbf{B. Expression orale — plan et éléments de réponse.}
\begin{itemize}
\item Introduction : 今天我说一说手机支付。 (Jīntiān wǒ shuō yi shuō shǒujī zhīfù.) → « Aujourd'hui, je vais parler du paiement mobile. »
\item Point 1 (Cameroun) : 在喀麦隆，很多人用手机转钱、付钱，这很方便。 → « Au Cameroun, beaucoup de gens transfèrent et paient avec le téléphone. »
\item Point 2 (Chine) : 在中国，人们常用微信支付，出门不带现金。 → « En Chine, on paie souvent avec WeChat, on sort sans argent liquide. »
\item Conclusion : 我觉得，手机支付改变了我们的生活。 → « Je pense que le paiement mobile a changé notre vie. »
\end{itemize}
Réponses possibles aux questions du jury :
\begin{itemize}
\item 你常用什么应用？ → 我常用微信和WhatsApp。 (Wǒ cháng yòng Wēixìn hé WhatsApp.) → « J'utilise souvent WeChat et WhatsApp. »
\item 你觉得手机支付安全吗？ → 我觉得手机支付很安全，也很方便。 (Wǒ juéde shǒujī zhīfù hěn ānquán, yě hěn fāngbiàn.) → « Je trouve que le paiement mobile est sûr et pratique. »
\end{itemize}
\textbf{Barème indicatif (sur 20 points) :} respect des consignes (comparaison + avis + vocabulaire) : 6 points ; correction grammaticale et caractères : 6 points ; richesse du vocabulaire : 4 points ; prononciation et aisance à l'oral : 4 points.
\textbf{Points de vigilance :} ne pas dire 比短信很好用 (pas de 很 après 比) ; prononcer 应用 (yìngyòng) avec le 4\textsuperscript{e} ton sur yìng ; à l'oral, parler lentement et soigner les tons des mots clés : 支付 (zhīfù), 方便 (fāngbiàn), 流行 (liúxíng).
\end{corrige}

\newpage

\coursec{Fiche bilan}

\begin{popfigure}
\begin{tikzpicture}[mindmap, grow cyclic, every node/.style={concept, text=white, font=\small}, root concept/.append style={concept color=popblue, font=\bfseries}, level 1/.append style={level distance=3.4cm, sibling angle=51.4}, level 1 concept/.append style={font=\footnotesize}]
\node[root concept]{TIC : marques chinoises et camerounaises}
  child[concept color=poporange]{ node[align=center]{Lexique TIC\\ \zh{手机、网络、软件}} }
  child[concept color=popgreen]{ node[align=center]{Marques chinoises\\ \zh{华为、小米、传音}} }
  child[concept color=poppurple]{ node[align=center]{Marques camerounaises\\ \zh{喀麦隆的品牌}} }
  child[concept color=poppink]{ node[align=center]{Grammaire\\ comparaison \zh{比}} }
  child[concept color=popgold]{ node[align=center]{Donner son avis\\ \zh{我认为……}} }
  child[concept color=popmuted]{ node[align=center]{Comparer\\ Cameroun / Chine} }
  child[concept color=popblue!60]{ node[align=center]{Mini-exposé\\ et réflexion} };
\end{tikzpicture}
\legende{Carte mentale de la leçon : les outils TIC de marques chinoises et camerounaises}
\end{popfigure}

\begin{bilan}
\textbf{1. Lexique clé à connaître par cœur}

\begin{center}
\begin{tabularx}{\linewidth}{|X|X|X|X|}
\hline
\rowcolor{popblueL} Caractères & Pinyin & Nature & Traduction \\
\hline
手机 & shǒujī & n. & téléphone portable \\
\hline
网络 & wǎngluò & n. & Internet, réseau \\
\hline
软件 & ruǎnjiàn & n. & logiciel, application \\
\hline
品牌 & pǐnpái & n. & marque \\
\hline
华为 & Huáwéi & n. propre & Huawei \\
\hline
小米 & Xiǎomǐ & n. propre & Xiaomi \\
\hline
传音 & Chuányīn & n. propre & Transsion (Tecno, Itel) \\
\hline
社交媒体 & shèjiāo méitǐ & n. & réseaux sociaux \\
\hline
移动支付 & yídòng zhīfù & n. & paiement mobile \\
\hline
便宜 & piányi & adj. & bon marché \\
\hline
方便 & fāngbiàn & adj. & pratique, commode \\
\hline
发展 & fāzhǎn & v./n. & (se) développer ; développement \\
\hline
\end{tabularx}
\end{center}

\textbf{2. Grammaire à connaître par cœur}

\begin{center}
\begin{tabularx}{\linewidth}{|X|X|X|}
\hline
\rowcolor{popblueL} Structure & Exemple & Traduction \\
\hline
A + 比 + B + adj. & 小米手机比苹果手机便宜。 \emph{Xiǎomǐ shǒujī bǐ Píngguǒ shǒujī piányi.} & Les téléphones Xiaomi sont moins chers que les iPhone. \\
\hline
A + 没有 + B + adj. & 这个品牌没有华为有名。 \emph{Zhège pǐnpái méiyǒu Huáwéi yǒumíng.} & Cette marque n'est pas aussi connue que Huawei. \\
\hline
越来越 + adj. & 用移动支付的人越来越多。 \emph{Yòng yídòng zhīfù de rén yuè lái yuè duō.} & De plus en plus de gens utilisent le paiement mobile. \\
\hline
我认为 / 我觉得 + phrase & 我认为网络很有用。 \emph{Wǒ rènwéi wǎngluò hěn yǒuyòng.} & Je pense qu'Internet est très utile. \\
\hline
因为……所以…… & 因为方便，所以大家都喜欢。 \emph{Yīnwèi fāngbiàn, suǒyǐ dàjiā dōu xǐhuan.} & Parce que c'est pratique, tout le monde aime. \\
\hline
\end{tabularx}
\end{center}

\textbf{3. Faits socioculturels à retenir}
\begin{itemize}
  \item En Chine : géants des TIC (Huawei, Xiaomi, Tencent avec WeChat \zh{微信}, Alibaba) ; le paiement mobile y est très répandu.
  \item Au Cameroun : forte présence des marques chinoises (Tecno, Itel, Huawei) ; le Mobile Money (MTN, Orange) transforme la vie quotidienne ; des start-up camerounaises émergent à Yaoundé et Douala.
  \item Ton respectueux et factuel : comparer sans hiérarchiser, valoriser les deux pays.
\end{itemize}

\textbf{4. Méthodes express}
\begin{itemize}
  \item \textbf{Comparer} : repérer les deux termes (A et B), choisir \zh{比} (supériorité) ou \zh{没有} (infériorité), placer l'adjectif à la fin.
  \item \textbf{Donner son avis} : \zh{我认为 / 我觉得} + phrase complète + justification avec \zh{因为……所以……}.
  \item \textbf{Mini-exposé} : introduction (\zh{大家好，今天我说一说……}) → 3 idées (marques, usages, avis) → conclusion (\zh{谢谢大家！}).
\end{itemize}
\end{bilan}

\begin{aretenir}[Checklist avant le Bac]
\begin{itemize}
  \item[$\square$] Je connais le lexique TIC : \zh{手机、网络、软件、品牌、移动支付、社交媒体}.
  \item[$\square$] Je sais écrire et prononcer les marques : \zh{华为、小米、传音}.
  \item[$\square$] Je maîtrise la comparaison avec \zh{比} : A + 比 + B + adjectif.
  \item[$\square$] Je sais nier la comparaison avec \zh{没有} (et jamais \zh{不} devant \zh{有}).
  \item[$\square$] J'emploie \zh{越来越} + adjectif pour exprimer une évolution.
  \item[$\square$] Je donne mon avis avec \zh{我认为} / \zh{我觉得} sans oublier le sujet.
  \item[$\square$] Je justifie avec \zh{因为……所以……} (les deux mots peuvent se combiner en chinois).
  \item[$\square$] Je place correctement les tons : shǒujī (3-1), wǎngluò (3-4), piányi (2-2 léger).
  \item[$\square$] Je connais les clés : 扌(main) dans \zh{手机}, 纟(soie) dans \zh{网络}, 车(voiture) dans \zh{软件}.
  \item[$\square$] Je peux comparer Cameroun et Chine de façon factuelle et respectueuse.
  \item[$\square$] Je sais présenter un mini-exposé de 5 phrases avec introduction et conclusion.
  \item[$\square$] Je relis ma production : ordre des mots, place de \zh{很} devant l'adjectif, ponctuation chinoise 。，
\end{itemize}
\end{aretenir}

\findecours

\end{document}
